% !TEX program = pdflatex
% Final working revision: 12 September 2026.
% Build with pdfLaTeX; bibliography is inline (no BibTeX step).
% Preserve the established statement/equation labels and numbering.
% Style: amsart 11pt, Latin Modern, 1.1in margins.
% Revision history and reproducibility checks are documented separately.
\documentclass[11pt,reqno]{amsart}
\usepackage[T1]{fontenc}
\usepackage[utf8]{inputenc}
\usepackage{lmodern}
\usepackage{amsmath,amssymb,amsthm,mathtools}
\usepackage{booktabs,array,longtable}
\usepackage[margin=1.1in]{geometry}
\usepackage{microtype}
\usepackage{xcolor}
\definecolor{linkcol}{rgb}{0.10,0.20,0.45}
\usepackage[colorlinks=true,linkcolor=linkcol,citecolor=linkcol,
            urlcolor=linkcol]{hyperref}
\hypersetup{
  pdftitle={A proof of Nivat's conjecture in its convex form},
  pdfauthor={Guancheng Pan},
  pdfsubject={Star configurations and low convex pattern complexity},
  pdfkeywords={Nivat conjecture, periodicity, convex window, zonotope},
  pdfdisplaydoctitle=true,bookmarksnumbered=true}


% Reproducible output, as in the other preprints: omit the creation and
% modification dates, the pdfTeX banner and the non-deterministic trailer id.
\pdfinfoomitdate=1
\pdfsuppressptexinfo=-1
\pdftrailerid{}

% ---------------------------------------------------------------------------
% Numbering.  One counter for all statements, run within sections.  Sections are
% numbered from 0, so that the notation section is Section 0 and every statement
% carries the number it carries in the manuscript.
% ---------------------------------------------------------------------------
\theoremstyle{plain}
\newtheorem{thm}{Theorem}[section]
\newtheorem{lem}[thm]{Lemma}
\newtheorem{prop}[thm]{Proposition}
\newtheorem{cor}[thm]{Corollary}
\theoremstyle{definition}
\newtheorem{dfn}[thm]{Definition}
\theoremstyle{remark}
\newtheorem{rem}[thm]{Remark}
\newtheorem{obs}[thm]{Observation}
\theoremstyle{plain}
\newtheorem*{thmA}{Theorem A}
\newtheorem*{thmB}{Theorem B}
% The manuscript states its main result in Section 0 as an unnumbered
% "Theorem T" and restates it, numbered, as Theorem 7.3.  Keeping that
% arrangement is what gives Remark 0.1 its number.
\newtheorem*{thmT}{Theorem T}
% Primed remarks.  They are companions to the results they sit beside and are
% numbered by hand precisely so that no later cross-reference number moves.
\theoremstyle{remark}
\newtheorem*{remgalois}{Remark 2.1$'$}
\newtheorem*{remnogen}{Remark 3.1$'$}
\theoremstyle{definition}
\newtheorem*{dfnhalf}{Definition}

% Equations are numbered within sections.  Because sections start at 0 and the
% equation counter is reset at each \section, the automatic numbers coincide
% with the manuscript's -- (0.1), (0.2), (1.1), (2.0), ..., (8.1) -- with one
% adjustment: Section 2's first display is (2.0), so its counter starts at -1.
% This is done with the counter rather than with \tag so that every equation
% gets its own hyperref destination.
\numberwithin{equation}{section}

% ---------------------------------------------------------------------------
% Notation
% ---------------------------------------------------------------------------
\DeclareMathOperator{\Per}{Per}
\DeclareMathOperator{\Conv}{Conv}
\DeclareMathOperator{\supp}{supp}
\DeclareMathOperator{\Newt}{Newt}
\DeclareMathOperator{\dist}{dist}
\DeclareMathOperator{\sgn}{sign}
\DeclareMathOperator{\spn}{span}
\DeclareMathOperator{\Gal}{Gal}
\DeclareMathOperator{\Orb}{Orb}
\DeclareMathOperator{\diam}{diam}
\DeclareMathOperator{\lcm}{lcm}
\DeclareMathOperator{\ord}{ord}
\DeclareMathOperator{\Tr}{Tr}
\newcommand{\R}{\mathbb{R}}
\newcommand{\Z}{\mathbb{Z}}
\newcommand{\N}{\mathbb{N}}
\newcommand{\C}{\mathbb{C}}
\newcommand{\Q}{\mathbb{Q}}
\newcommand{\Fp}{\mathbb{F}_p}
\newcommand{\one}{\mathbb{1}}
\newcommand{\cA}{\mathcal{A}}
\newcommand{\cE}{\mathcal{E}}
\newcommand{\cH}{\mathcal{H}}
\newcommand{\cK}{\mathcal{K}}
\newcommand{\cL}{\mathcal{L}}
\newcommand{\cR}{\mathcal{R}}
\newcommand{\cS}{\mathcal{S}}
\newcommand{\cT}{\mathcal{T}}
\newcommand{\cU}{\mathcal{U}}
\newcommand{\cW}{\mathcal{W}}
\newcommand{\dbr}[1]{[\![#1]\!]}
\newcommand{\su}{\mathsf{u}}
\newcommand{\sw}{\mathsf{w}}

\allowdisplaybreaks
\setlength{\parskip}{0.15\baselineskip}
\raggedbottom
\emergencystretch=2.5em
\setlength{\tabcolsep}{4pt}
\predisplaypenalty=10000
\medmuskip=4mu plus 2mu minus 1mu
\clubpenalty=10000
\widowpenalty=10000
\displaywidowpenalty=10000

\title[A proof of Nivat's conjecture in its convex form]
      {A proof of Nivat's conjecture\\
       in its convex form}
% Author information updated on 2026-09-30.
\author[G. Pan]{Guancheng Pan}
\makeatletter
\renewcommand{\@setauthors}{%
  \begin{center}
    {\large Guancheng Pan}\\[1.2ex]
    {\small Chu Kochen Honors College, Zhejiang University}\\[0.8ex]
    {\small\texttt{3250101086@zju.edu.cn}}
  \end{center}}
\makeatother

\date{1 October 2026}
% The subject classification and the key words are printed as a centred block
% under the abstract, not in amsart's first-page footer, so that page 1 reads
% the same way in all of the authors' preprints.

\begin{document}

\begin{abstract}
We prove Nivat's conjecture in its convex form: if a configuration
$\xi\colon\Z^2\to\cA$ over a finite alphabet satisfies $P_\xi(S)\le|S|$ for
some non-empty finite lattice-convex window $S$, that is, one with
$S=\Conv(S)\cap\Z^2$, then $\xi$ is periodic.  Rectangles are lattice-convex,
so this contains Nivat's conjecture itself: $P_\xi(m,n)\le mn$ for some
$m,n\ge1$ forces $\xi$ to be periodic.  Convexity is exactly the right hypothesis:
by Khetan's counterexample the statement fails for a non-convex window of
full affine span.
The whole proof, including every result we quote from the literature, has been
formalised in Lean~4 with Mathlib; the kernel reports that the final theorem
depends on no axiom beyond the three standard ones of Lean's logic.

The proof has two parts.  The first occupies \S\S\ref{sec:notation}--\ref{sec:quadratic}: we call
$\theta=\sum_{i=1}^mF_i\colon\Z^2\to\Fp$ a \emph{star
configuration} when $m\ge2$, each $F_i$ has a non-zero period $k_iv_i$, is not doubly
periodic, and agrees with doubly periodic fields on the two half-planes
$\{\pi_i<\ell_i\}$ and $\{\pi_i>r_i\}$, the primitive directions $v_i$ being
pairwise non-parallel; and we prove that $P_\theta(S)\ge|S|+1$ for every
non-empty finite lattice-convex $S$.  The proof attaches to $\theta$ an
\emph{exceptional spectrum} $\Lambda_i$ of roots of unity in each direction, a
Laurent polynomial $A=\prod_iA_i(T^{v_i})$, and the zonotope
$Z=\Newt(A)=\sum_i[0,d_iv_i]$.  For a spectrum-preserving scalar encoding
$\eta=w(\theta)$, every Laurent polynomial $f$ with $f(T)\eta$ constant is
divisible by $A$.  This bounds the affine deficit by $|R_Z(S)|$:
$\dim U_S\ge|S|+1-|R_Z(S)|$.  In the remaining case of a finite-difference
dichotomy, the same zonotope supplies $|R_Z(S)|$ quadratic observables
independent modulo $U_S$, giving the required complexity bound.

The second part (\S\ref{sec:reduction}) reduces the convex conjecture to the
first, and so proves it.  On paper it quotes the annihilator and periodic
decomposition theorem of Kari and Szabados \cite{KS} and Colle's results
\cite{Colle23a} on one-sided nonexpansive directions and on full periodicity over
a region; in the formalisation these are proved as well.
Appendix~\ref{app:szabados} proves in full the $\Fp$ and convex form of
Szabados' theorem, and a further appendix reports computational checks of the
intermediate assertions.  Appendix~\ref{app:lean} describes the formalisation
and Appendix~\ref{app:ai} how the work was done: the mathematics and the Lean
development were produced with the author's AI research system.
\end{abstract}

\maketitle

\begin{center}\small
\emph{2020 Mathematics Subject Classification.}\ Primary 37B51;
Secondary 37B10, 68R15, 52C05, 13F20.\\
\emph{Key words and phrases.}\ Nivat's conjecture; low pattern complexity;
convex window; periodic decomposition; star configuration; annihilator;
Newton polygon; zonotope.
\end{center}

\section*{Introduction}

\subsection*{The problem}

Let $\cA$ be a finite alphabet and $\xi\colon\Z^2\to\cA$ a configuration.  For a
finite $Q\subset\Z^2$ write $P_\xi(Q)$ for the number of distinct patterns
$\xi|_{u+Q}$, $u\in\Z^2$, and $P_\xi(m,n)=P_\xi(\dbr{m}\times\dbr{n})$ for the
rectangular complexity, $\dbr{m}=\{0,\dots,m-1\}$.  The Morse--Hedlund theorem
\cite{MH38} says that a bi-infinite word is periodic as soon as its complexity
function satisfies $P(n)\le n$ for a single $n$.  Nivat asked, in his 1997 ICALP
address \cite{Nivat}, whether the two-dimensional analogue holds:

\begin{quote}
\emph{if $P_\xi(m,n)\le mn$ for some $m,n\ge1$, is $\xi$ periodic?}
\end{quote}

This paper answers the question in the affirmative, and proves more: the
same conclusion holds for every lattice-convex window (Theorem~B below).
Before this, partial results fell into two families.  On the
combinatorial and dynamical side, Sander--Tijdeman \cite{ST},
Epifanio--Koskas--Mignosi \cite{EKM} and Quas--Zamboni \cite{QZ} proved
periodicity under complexity bounds of the form $P_\xi(m,n)\le mn/C$; Cyr and
Kra \cite{CyrKra15} brought the constant down to $2$ by way of nonexpansive
subdynamics and settled in \cite{CyrKra16} the short-rectangle case
$P_\xi(k,n)\le kn$ for $k\le3$; Colle and Garibaldi \cite{CGalpha} gave a
version corrected by the size of the alphabet.  On the algebraic side,
Kari and Szabados \cite{KS} observed that low complexity forces the existence of
a non-trivial annihilating Laurent polynomial and deduced that the configuration
decomposes as a sum of periodic configurations over $\Z$; Szabados \cite{Sz}
settled the case of two components; Colle \cite{Colle23a,Colle23b} developed
the interaction between minimal periodic decompositions and one-sided
nonexpansive directions; and Colle and Garibaldi \cite{CG} recently proved, in
the framework of $\Z_p$-minimal periodic decompositions, a structure theorem —
the orbit closure of an aperiodic configuration of low convex complexity
contains an aperiodic configuration each of whose periodic components is fully
periodic on two disjoint half-planes — and deduced from it the case $k\le4$.
Kari and Moutot \cite{KM20,KM23} proved that the orbit closure of a
configuration of low convex complexity contains a periodic configuration, and
the related decidability statement for low-complexity tilings.

The \emph{convex} form of the conjecture replaces the rectangle by an arbitrary
window that is convex in the only sense available on a lattice: a finite
$S\subset\Z^2$ is \emph{lattice-convex} if $S=\Conv(S)\cap\Z^2$.  A
configuration has \emph{low convex pattern complexity} if $P_\xi(S)\le|S|$ for
some non-empty finite lattice-convex $S$.  The convexity hypothesis is not
decoration.  Khetan \cite{Khe} has constructed an exact cluster
$F\subset\Z^2$ of cardinality $8$ with full affine span and an $F$-tiling whose
orbit closure contains no $1$-periodic $F$-tiling; since every $F$-tiling has
low complexity with respect to the window $-F$, this is a counterexample to the
statement for non-convex windows of full affine span.  The convex form is
therefore the natural general statement, it contains the rectangular one, and
it is the one we prove.

\subsection*{Results}

The main result is Theorem~B, the convex Nivat conjecture.  It rests on
Theorem~A, a complexity lower bound for star configurations proved from
scratch; the reduction of Theorem~B to Theorem~A quotes only results of
\cite{KS} and \cite{Colle23a}, and these too are proved in the formalisation.

Call $\theta\colon\Z^2\to\Fp$ a \emph{star configuration} if
$\theta=\sum_{i=1}^mF_i$ with $m\ge2$, where each $F_i$ is periodic in a
direction $v_i$, is not doubly periodic, and coincides with doubly periodic
fields on two disjoint half-planes bounded by lines of the form
$\pi_i=\det(v_i,\cdot)=\text{const}$, the primitive
directions $v_1,\dots,v_m$ being pairwise non-parallel.  The precise axioms
\textup{(S1)}--\textup{(S3)} are in \S\ref{ssec:star}.  Informally: $m$
one-dimensionally periodic fields, each of which is doubly periodic outside a
strip, superimposed in $m$ pairwise distinct directions.

\begin{thmA}[Theorem~\ref{thm:main}]
Let $\theta$ be a star configuration and let $S\subset\Z^2$ be non-empty,
finite and lattice-convex.  Then $P_\theta(S)\ge|S|+1$.
\end{thmA}

Theorem~A is proved in \S\S\ref{sec:notation}--\ref{sec:quadratic} using no
input from the literature beyond the facts that $\C[T_1^{\pm},T_2^{\pm}]$ is a
unique factorisation domain, that Newton polygons are additive under
multiplication (Ostrowski), the Chinese remainder theorem, Hilbert's
Nullstellensatz, and the existence of unimodular triangulations of lattice
polygons.

\begin{thmB}[Convex Nivat conjecture; Theorem~\ref{thm:convex-nivat}]
Let $\cA$ be finite and $\xi\colon\Z^2\to\cA$.  If $P_\xi(S)\le|S|$
for some non-empty finite lattice-convex $S$, then $\xi$ is periodic.
In particular \textup{(}Nivat's conjecture\textup{)}, if $P_\xi(m,n)\le mn$ for
some $m,n\ge1$, then $\xi$ is periodic.
\end{thmB}

Both statements are checked in Lean~4: Theorem~B is \texttt{Nivat.convex\_nivat}
and its rectangular case is \texttt{Nivat.nivat\_conjecture} (see
\emph{Machine verification} below).

Theorem~B is proved by reducing it to Theorem~A (\S\ref{sec:reduction}).  The
chain runs as follows: the annihilator and the
periodic decomposition of Kari and Szabados \cite{KS}; the two results of Colle
\cite{Colle23a} — a line that is one-sidedly nonexpansive in both directions,
and an aperiodic configuration $\vartheta$ in the orbit closure which is fully
periodic on a region; and then the two steps of this paper.  The first is that
each periodic component is fully periodic on one half-plane $U_i$
(Proposition~\ref{prop:firsthalf}: extend the full periodicity over the region
to a doubly periodic field, isolate the component with a difference operator
$Q_i$, and run a finite-state argument).  The second is that each component is
fully periodic on a half-plane $V_i$ on the far side
(Theorem~\ref{thm:secondhalf}: push to the limit along the direction of some
component; the limit is a sum of two periodic configurations in distinct
directions, hence periodic by Theorem~\ref{thm:szabados} of
Appendix~\ref{app:szabados}, and therefore every limit on the far side of the
component under examination is doubly periodic; that there are only finitely
many such limits then forces the far side itself to agree with a doubly periodic
field).  Finally Lemma~\ref{lem:normalise} normalises such a $\vartheta$ into a
star configuration, contradicting Theorem~A.  Theorem~\ref{thm:structure} —
each component is fully periodic on two disjoint half-planes — coincides
formally with the structure theorem of \cite{CG}; the proof given here is
different from the one there.  After the region has been extracted, the
half-plane arguments do not require the $\Fp$-decomposition to be minimal;
the region extraction uses a counterexample of minimal $\Z$-order, as specified
in Remark~\ref{rem:minimal}.

\subsection*{Machine verification}

The complete proof of Theorem~B is formalised in Lean~4 \cite{Lean4} on top of
Mathlib \cite{Mathlib}.  The formal theorem \texttt{Nivat.convex\_nivat} says:
for every type $\alpha$ with \texttt{[Finite $\alpha$]}, every
$\xi\colon\Z\times\Z\to\alpha$ and every non-empty \texttt{Finset} $S$ of
$\Z\times\Z$ that is lattice-convex, $P_\xi(S)\le|S|$ implies that $\xi$ has a
non-zero period; \texttt{Nivat.nivat\_conjecture} is the case
$S=[1,n]\times[1,k]$.  Here $P_\xi(S)$ is literally the number of distinct
restrictions of $\xi$ to the translates $u+S$, $u\in\Z^2$, and lattice
convexity is $S=\Conv(S)\cap\Z^2$ with the real convex hull of Mathlib.

The Lean kernel reports that both theorems depend only on \texttt{propext},
\texttt{Classical.choice} and \texttt{Quot.sound}, the axioms of Lean's own
logic.  There is no \texttt{sorry}, and the development declares no axiom of its
own.  In particular nothing is taken from the literature on trust: the
annihilator and decomposition theorems of Kari and Szabados and the results of
Colle used in \S\ref{sec:reduction}, with the theorems of Boyle--Lind and
Cyr--Kra behind them, are proved in the formalisation, which by volume is
mostly about them (some $192{,}000$ of its $216{,}000$ lines).  The development
has $449$ modules and builds from a clean checkout; its source is attached to
this paper as an ancillary file.  Appendix~\ref{app:lean} explains how the
formal definitions match ours and how to check the build.

Formalising the quoted results also tested them.  One step of the source did not
survive: in the proof of Lemma~3.5(i) of the arXiv version of \cite{Colle23a},
the swept sets used as shells are in general neither enveloped nor squeezed
between the required strips.  The formalisation replaces that construction by a
different one; the lemma itself, and with it everything we quote, stands
(Appendix~\ref{app:lean}).

\subsection*{The shape of the proof of Theorem~A}

Let $H_i$ be a period of $F_i$ along $v_i$ that also preserves every tail field,
and let $D=\prod_{i=1}^m(T^{H_i}-1)$.

\emph{A dichotomy} (\S\ref{sec:firstcase}).  Applying $D$ to a function that
depends on $\theta$ through a finite window kills everything except a bounded
region, because outside the union of the pairwise intersections of $m$ strips in
pairwise distinct directions at most one component is still varying.  If
$D\one[\theta=a]\ne0$ for some colour $a$, the theorem follows for
\emph{arbitrary} finite $S$, convex or not, by a one-line integral-domain
argument.  The remaining case is $D\one[\theta=a]\equiv0$ for every colour $a$.

\emph{The exceptional spectrum and its zonotope} (\S\ref{sec:spectrum}).
Pushing $\theta$ to infinity along $\pm H_i$ isolates the $i$-th component
against a doubly periodic background, producing configurations
$\Xi_i^\pm=F_i+B_i^\pm$ in the orbit closure together with their own pure tails
$G_i^{\sigma,L},G_i^{\sigma,R}$.  Their colour-indicator differences are invariant
under $H_i$, hence decompose under the $\kappa_i$-th roots of unity; the set of
roots that actually occur is the exceptional spectrum $\Lambda_i$, non-empty
exactly because $F_i$ is not doubly periodic.  Set
$A=\prod_iA_i(T^{v_i})$ with $A_i(t)=\prod_{\lambda\in\Lambda_i}(t-\lambda)$,
and $Z=\Newt(A)=\sum_i[0,d_iv_i]$, $d_i=|\Lambda_i|$.  The key divisibility is
Lemma~\ref{lem:divisibility}: if $f(T)\eta$ is constant for a spectrum-preserving
scalar encoding $\eta=w(\theta)$ as in Lemma~\ref{lem:encoding}, then $A\mid f$.
Counting, this says that a window $S$
carries at most $|R_Z(S)|$ independent affine relations, where
$R_Z(S)=\{r:r+Z\subseteq\Conv(S)\}$: the \emph{cost} of the zonotope.

\emph{The quadratic witness} (\S\S\ref{sec:background}--\ref{sec:zonotope}).
The affine bound may fall short by $|R_Z(S)|$.  We supply that many additional
quadratic observables.  This requires a single vector $d$ with
$D_z[\eta(z)\eta(z+d)]\ne0$ \emph{and} $d\in(Z-Z)\cap\Z^2$.  Existence of some
such $d$ is \S\ref{sec:twopoint}; confinement to $Z-Z$ comes from a pair of
recursions in the two variables $(d,z)$ (\S\ref{sec:quotient}) which say that
the Laurent transform of the witness kills the ideal $(A(Y),A(X/Y))$ in
$K[Y^{\pm}]$, $K=\C(X_1,X_2)$, whose quotient is spanned by monomials with
exponents in $Z-Z$; and the passage from $d\in Z-Z$ to a \emph{pair} of lattice
points of $Z$ differing by $d$ is a unimodular triangulation argument
(\S\ref{sec:zonotope}).

\emph{The cancellation} (\S\ref{sec:quadratic}).  Each $r\in R_Z(S)$ now yields
a quadratic observable $\Phi_r$ on the pattern set, and these are independent
modulo the linear ones.  Their number compensates for the full possible affine
deficit; that is Remark~\ref{rem:balance}.

\subsection*{Organisation}

\S\ref{sec:notation} fixes notation, the axioms of a star configuration, and the
main theorem.  \S\ref{sec:firstcase} proves the dichotomy.
\S\ref{sec:spectrum} constructs $\Lambda_i$, $A$, $Z$ and the dimension budget.
\S\ref{sec:background} shows that $A$ applied to the colour field of $\theta$ is
a doubly periodic background.  \S\ref{sec:twopoint} produces a non-vanishing
two-point difference, \S\ref{sec:quotient} confines it to $Z-Z$, and
\S\ref{sec:zonotope} converts that into a pair of lattice points of $Z$.
\S\ref{sec:quadratic} assembles Theorem~A, and \S\ref{sec:reduction} deduces
Theorem~B: \S\ref{ssec:halfplanes} sets up half-planes, regions and full
periodicity; \S\ref{ssec:external} collects the results that are quoted;
\S\ref{ssec:firsthalf} proves the first half-plane
(Proposition~\ref{prop:firsthalf}); \S\ref{ssec:secondhalf} proves the second
(Theorem~\ref{thm:secondhalf});
\S\ref{ssec:normalisation} normalises; \S\ref{ssec:nivat} proves Theorem~B.
Appendix~\ref{app:skeleton} is a one-page skeleton of the argument.
Appendix~\ref{app:edge} briefly describes, without proof and for the record
only, a different route that was explored and is not used here.
Appendix~\ref{app:numerics} reports the numerical verification.
Appendix~\ref{app:szabados} proves the $\Fp$ and convex form of Szabados'
theorem.  Appendix~\ref{app:lean} describes the Lean formalisation and
Appendix~\ref{app:ai} records how the work was done with an AI system.

\smallskip\noindent\textbf{Dependencies.}
\S\S\ref{sec:notation}--\ref{sec:quadratic} are self-contained.
\S\ref{sec:reduction} quotes \cite{KS} (annihilators and periodic
decompositions), \cite[Theorem~1.9]{Colle23a}, and the proof of
\cite[Lemma~4.6]{Colle23a}; the formalisation proves all of these.
Appendix~\ref{app:szabados} uses nothing beyond the one-dimensional
Morse--Hedlund theorem \cite{MH38}.

\setcounter{section}{-1}
\section{Notation, hypotheses and the main theorem}\label{sec:notation}

\subsection{Star configurations}\label{ssec:star}

Let $p$ be a prime and $m\ge2$.  Let $v_1,\dots,v_m\in\Z^2$ be \emph{primitive}
(coprime coordinates) and \emph{pairwise non-parallel}.  Write
\[\pi_i(z)=\det(v_i,z)\qquad(z\in\Z^2),\]
a surjective linear form $\Z^2\to\Z$, surjective because $v_i$ is primitive.

Call $G\colon\Z^2\to\Fp$ \emph{doubly periodic} if its period group
$\Per(G)=\{u\in\Z^2:T^uG=G\}$ contains two linearly independent vectors;
equivalently, $\Per(G)$ is a finite-index sublattice of $\Z^2$.

Let $F_i\colon\Z^2\to\Fp$ satisfy
\begin{itemize}
\item[\textup{(S1)}] $F_i$ has period $k_iv_i$ for some integer $k_i\ge1$;
\item[\textup{(S2)}] $F_i$ is \emph{not} doubly periodic;
\item[\textup{(S3)}] there are doubly periodic fields $L_i,R_i$ and integers
      $\ell_i\le r_i+1$ with
      \[F_i=L_i\ \text{ on }\ \{\pi_i<\ell_i\},\qquad
        F_i=R_i\ \text{ on }\ \{\pi_i>r_i\}.\]
\end{itemize}
(If one starts with $\ell_i>r_i+1$, shrinking $\ell_i$ to $r_i+1$ keeps
\textup{(S3)}.)  We call $\{\pi_i<\ell_i\}$ and $\{\pi_i>r_i\}$ the \emph{left}
and \emph{right tail regions} of component $i$, call $L_i,R_i$ its \emph{left}
and \emph{right tails}, and call $\{\ell_i\le\pi_i\le r_i\}$ the
\emph{transition strip} (possibly empty).  We do not require $L_i=R_i$, do not
require either to be constant, and do not require $k_i=1$.  Put
\[\theta=\sum_{i=1}^mF_i\colon\ \Z^2\to\Fp,\]
and call $\theta$ a \emph{star configuration}.

\subsection{Complexity and the main theorem}\label{ssec:main}

For finite $Q\subset\Z^2$ let $\theta|_{u+Q}$ denote the function
$q\mapsto\theta(u+q)$ and
\[P_\theta(Q)=\#\{\theta|_{u+Q}:u\in\Z^2\}.\]
Call a finite $S\subset\Z^2$ \emph{lattice-convex} if $S=\Conv(S)\cap\Z^2$.

\begin{thmT}
Let $\theta$ be a star configuration and let $S\subset\Z^2$ be non-empty,
finite and lattice-convex.  Then
\[P_\theta(S)\ \ge\ |S|+1 .\]
\end{thmT}

\begin{rem}[Periodic backgrounds]\label{rem:background}
If a doubly periodic field $B$ is added to $\theta$, absorb $B$ into $F_1$ and
replace the tangential period of $F_1$ by a multiple that also preserves $B$;
such a multiple exists because $\Per(B)$ has finite index, so
$\Per(B)\cap\Z k_1v_1\ne0$.  Then $F_1+B$ is still not doubly periodic
(otherwise $F_1=(F_1+B)-B$ would be), and $L_1+B$, $R_1+B$ are still doubly
periodic.  Hence Theorem~T covers an arbitrary doubly periodic
background.
\end{rem}

\subsection{Translations, differences and basic conventions}\label{ssec:conv}

Write $T^uf(z)=f(z+u)$.  A \emph{Laurent polynomial} is a finite sum
$f=\sum_uc_uT^u$ with complex coefficients unless stated otherwise, with support
$\supp(f)=\{u:c_u\ne0\}$ and Newton polygon $\Newt(f)=\Conv(\supp(f))$.  The
ring $\C[T_1^{\pm},T_2^{\pm}]$ is an integral domain and a unique factorisation
domain, and its units are exactly the monomials $cT^u$ with $c\in\C^\times$.

\emph{Colour indicators, the encoding, the differences and all linear algebra
take place in characteristic zero}; only the summation of the original
components happens in $\Fp$.  In particular $\one[\theta=a]$ takes values in
$\{0,1\}\subset\Z$.

\smallskip\noindent\textbf{The common period lattice and the vectors $H_i$.}
Put
\[\Gamma:=\bigcap_{j=1}^m\bigl(\Per(L_j)\cap\Per(R_j)\bigr),\]
a finite intersection of finite-index sublattices, hence of finite index; choose
$N\in\Z_{>0}$ with $N\Z^2\subseteq\Gamma$.  For each $i$ we have
$Nk_iv_i\in\Gamma$, so we may put
\[H_i:=\kappa_iv_i,\qquad
  \kappa_i:=\min\{\kappa\in k_i\Z_{>0}:\ \kappa v_i\in\Gamma\}.\]
Thus \emph{$H_i$ preserves $F_i$ (by \textup{(S1)}) and preserves every
$L_j,R_j$, $j=1,\dots,m$}.

\smallskip\noindent\textbf{The difference operator and its formal exponents.}
Put
\[D=\prod_{i=1}^m\bigl(T^{H_i}-1\bigr),\qquad
  I_a=\one[\theta=a],\qquad e_\theta=(I_a)_{a\in\Fp}.\]
Then $D\ne0$: each factor is non-zero and $\C[T^{\pm}]$ is an integral domain.
For $C\subseteq[m]:=\{1,\dots,m\}$ write $H_C=\sum_{i\in C}H_i$, and call
\[\cE=\{H_C:C\subseteq[m]\},\qquad
  \cE_i=\{H_C:C\subseteq[m]\setminus\{i\}\}\]
the \emph{formal exponent sets}.  For an arbitrary function $g$ one has the
formal expansion
\begin{equation}
(Dg)(z)=\sum_{C\subseteq[m]}(-1)^{m-|C|}\,g\bigl(z+H_C\bigr),
\label{eq:Dexp}
\end{equation}
valid as an operator identity irrespective of whether the vectors $H_C$ for
different $C$ coincide.  Thus $\supp(D)\subseteq\cE$, and for $m\ge3$ the
inclusion may be strict, since distinct subsets may have equal sums (for
instance if $H_3=H_1+H_2$).  \emph{Every pairing argument made on a formal
expansion below is indexed by subsets $C$, never by the points $H_C$.}
Likewise, for $Q_i:=\prod_{j\ne i}(T^{H_j}-1)$ one has $D=(T^{H_i}-1)Q_i$ and
\begin{equation}
(Q_ig)(z)=\sum_{C\subseteq[m]\setminus\{i\}}(-1)^{m-1-|C|}\,g\bigl(z+H_C\bigr).
\label{eq:Qexp}
\end{equation}

\section{Finite differences and the first case}\label{sec:firstcase}

\subsection{Differencing turns a local function into one of finite support}

\begin{lem}\label{lem:finite-support}
Let $h(z)=\varphi\bigl(\theta|_{z+W}\bigr)$ with $W\subset\Z^2$ finite.  Then
$Dh$ has finite support.
\end{lem}

\begin{proof}
If $W=\varnothing$, then $h$ is constant and $Dh=0$.  Assume henceforth that
$W\ne\varnothing$.
Put $M=\cE+W$, a finite set, and $\diam_i=\max_{s\in M}|\pi_i(s)|$.  Define the
\emph{widened strips}
\[\Sigma_i=\bigl\{z:\ \ell_i-\diam_i\le\pi_i(z)\le r_i+\diam_i\bigr\},\]
and call direction $i$ \emph{active at the anchor $z$} if $z\in\Sigma_i$.  If
$i$ is not active at $z$, then for every $s\in M$ the values $\pi_i(z+s)$ and
$\pi_i(z)$ lie both in $\{\pi_i<\ell_i\}$ or both in $\{\pi_i>r_i\}$, so on
$z+M$ the field $F_i$ is identically one pure tail
$\varepsilon_i\in\{L_i,R_i\}$.

\smallskip\noindent\emph{Claim: if at most one direction is active at $z$, then
$Dh(z)=0$.}  Let $J$ be the set of active directions, $|J|\le1$.  Take
$j_0\in J$ if $J\ne\varnothing$, and $j_0=1$ otherwise.  On $z+M$,
\[\theta=G+B,\qquad
  G=\begin{cases}F_{j_0},&J=\{j_0\}\\ 0,&J=\varnothing\end{cases},\qquad
  B=\sum_{j\notin J}\varepsilon_j ,\]
where $B$ is doubly periodic and preserved by every $H_l$, and $G$ is preserved
by $H_{j_0}$ by \textup{(S1)}.  Hence $G+B$ is preserved by $H_{j_0}$.  Expand
$D=(T^{H_{j_0}}-1)Q_{j_0}$ using \eqref{eq:Qexp}:
\[Dh(z)=\sum_{C\subseteq[m]\setminus\{j_0\}}(-1)^{m-1-|C|}
  \Bigl[h\bigl(z+H_C+H_{j_0}\bigr)-h\bigl(z+H_C\bigr)\Bigr].\]
Fix $C$.  Since $H_C$ and $H_C+H_{j_0}$ both lie in $\cE$, the two windows
$z+H_C+W$ and $z+H_C+H_{j_0}+W$ are contained in $z+M$, where $\theta=G+B$.  So
for $y\in W$,
\[\theta\bigl(z+H_C+H_{j_0}+y\bigr)=(G+B)\bigl(z+H_C+y+H_{j_0}\bigr)
 =(G+B)\bigl(z+H_C+y\bigr)=\theta\bigl(z+H_C+y\bigr),\]
so the two $W$-patterns agree and the bracket vanishes.  This proves the claim.

\smallskip\noindent\emph{Conclusion.}  Therefore
$\supp(Dh)\subseteq\bigcup_{i\ne j}(\Sigma_i\cap\Sigma_j)$.  As
$v_i\not\parallel v_j$, the intersection of two strips of bounded width in
non-parallel directions is bounded, so $Dh$ has finite support.
\end{proof}

\subsection{The integral-domain argument}

\begin{lem}\label{lem:domain}
Let $J\colon\Z^2\to\C$ be non-zero with finite support and let $f$ be a non-zero
Laurent polynomial.  Then $f(T)J\ne0$.
\end{lem}

\begin{proof}
Write $\widehat J(X)=\sum_zJ(z)X^{-z}\in\C[X_1^{\pm},X_2^{\pm}]$, which is
non-zero.  For $f=\sum_uc_uT^u$,
\[\sum_z\bigl(f(T)J\bigr)(z)X^{-z}
  =\sum_uc_u\sum_zJ(z+u)X^{-z}
  =\sum_uc_u\sum_{z'}J(z')X^{-z'+u}
  =\Bigl(\sum_uc_uX^{u}\Bigr)\widehat J(X).\]
Both factors are non-zero Laurent polynomials and the ring is an integral
domain, so the product is non-zero.
\end{proof}

\subsection{The first case}

\begin{prop}\label{prop:case-A}
If $DI_a\ne0$ for some $a\in\Fp$, then $P_\theta(S)\ge|S|+1$ for
\emph{every} non-empty finite $S\subset\Z^2$; no convexity is needed.
\end{prop}

\begin{proof}
Suppose $P_\theta(S)\le|S|$.  Write $\cL=\{\theta|_{u+S}:u\in\Z^2\}$, so
$|\cL|\le|S|$.  On $\cL$ consider the $|S|+1$ functions
\[\mathbf 1;\qquad \gamma\mapsto\one[\gamma(s)=a]\quad(s\in S).\]
They lie in a space of dimension $\le|\cL|\le|S|$, hence are linearly
dependent: there are rationals $c_0,(c_s)_{s\in S}$, not all zero, with
\[\sum_{s\in S}c_s\,\one[\gamma(s)=a]=c_0\qquad\text{for all }\gamma\in\cL .\]
Substituting $\gamma=\theta|_{u+S}$ gives $q(T)I_a=c_0$ with
$q=\sum_sc_sT^s$.  If $q=0$ then all $c_s=0$ and hence $c_0=0$, a
contradiction; so $q\ne0$.

Apply $D$ to both sides.  The right-hand side gives $Dc_0=0$, since a constant
is killed by $T^{H_1}-1$; the left-hand side gives $q(T)\,DI_a$.  By
Lemma~\ref{lem:finite-support} and the hypothesis, $DI_a$ is non-zero with
finite support, so $q(T)DI_a\ne0$ by Lemma~\ref{lem:domain}.  Contradiction.
\end{proof}

\noindent\emph{From now on we assume}
\begin{equation}
DI_a=0\quad\text{for every }a\in\Fp. \label{eq:caseB}
\end{equation}
Consequently $D\,w(\theta)=\sum_aw(a)DI_a=0$ for every function
$w\colon\Fp\to\Z$.

\begin{rem}[The dichotomy does not depend on the choice of the $H_i$]
\label{rem:dichotomy-stable}
Replacing $H_i$ by a positive multiple $n_iH_i$ produces
$D'=\prod_i(T^{n_iH_i}-1)$, which is a multiple of $D$ because
$T^{nH}-1=(T^H-1)(1+T^H+\cdots+T^{(n-1)H})$.  Hence $DI_a=0$ implies
$D'I_a=0$; conversely, if $DI_a\ne0$ then it is non-zero of finite support by
Lemma~\ref{lem:finite-support}, and Lemma~\ref{lem:domain} gives
$D'I_a=(D'/D)(T)\,DI_a\ne0$.  Any two choices satisfying the conventions of
\S\ref{ssec:conv}, namely $\kappa_i\in k_i\Z_{>0}$ and
$H_i=\kappa_iv_i\in\Gamma$, can be compared through a common positive
multiple in each direction.  Thus \eqref{eq:caseB} is independent of these
choices.  Each $H_i$ must preserve its component $F_i$ as well as every tail
field; preservation of the tails alone is insufficient.
\end{rem}

\section{Exceptional spectra, the encoding, and the dimension budget for
affine relations}\label{sec:spectrum}

\setcounter{thm}{-1}
\setcounter{equation}{-1}% the first display of this section is (2.0)
\subsection{Isolating configurations}\label{ssec:isolating}

Fix $i$.  For $j\ne i$ we have
$\pi_j(H_i)=\det(v_j,H_i)=\kappa_i\det(v_j,v_i)\ne0$, so
$\pi_j(nH_i)\to\pm\infty$; hence on any fixed finite set $F_j$ is eventually
identically equal, as $n\to+\infty$, to
\[\varepsilon_j^+=\begin{cases}R_j,&\pi_j(H_i)>0,\\
                               L_j,&\pi_j(H_i)<0,\end{cases}\]
while $F_i$ is preserved by $H_i$.  So $T^{nH_i}\theta$ stabilises on every
finite set, with limit
\[\Xi_i^+=F_i+B_i^+,\qquad B_i^+=\sum_{j\ne i}\varepsilon_j^+ .\]
Letting $n\to-\infty$ gives $\Xi_i^-=F_i+B_i^-$, with the opposite tail
choices: $\varepsilon_j^-$ is $L_j$ when $\pi_j(H_i)>0$ and $R_j$ otherwise.
Both belong to the orbit closure $X_\theta=\overline{\{T^u\theta:u\in\Z^2\}}$ in
the product topology.  Since $\kappa_i>0$ we have
$\sgn\pi_j(H_i)=\sgn\pi_j(v_i)$, so that
\begin{equation}
\varepsilon_j^\sigma=\begin{cases}R_j,&\sigma\pi_j(v_i)>0,\\
                                  L_j,&\sigma\pi_j(v_i)<0.\end{cases}
\label{eq:tailchoice}
\end{equation}

By \textup{(S3)},
\[\Xi_i^\sigma=G_i^{\sigma,L}:=L_i+B_i^\sigma\ \text{ on }\ \{\pi_i<\ell_i\},
  \qquad
  \Xi_i^\sigma=G_i^{\sigma,R}:=R_i+B_i^\sigma\ \text{ on }\ \{\pi_i>r_i\},\]
and $G_i^{\sigma,L},G_i^{\sigma,R}$ are doubly periodic and preserved by
$\Gamma$.  They too belong to $X_\theta$: take $(v_i,u_i)$ a basis of $\Z^2$
with $\det(v_i,u_i)=1$ and put $u=Nu_i\in\Gamma$, so that $\pi_i(u)=N>0$.  For
fixed $z$ and $n$ large, $\pi_i(z+nu)>r_i$, whence
\[\bigl(T^{nu}\Xi_i^\sigma\bigr)(z)=R_i(z+nu)+B_i^\sigma(z+nu)
  =R_i(z)+B_i^\sigma(z)=G_i^{\sigma,R}(z),\]
that is, $T^{nu}\Xi_i^\sigma\to G_i^{\sigma,R}$; using $-u$ instead gives
$G_i^{\sigma,L}$.

We record that \emph{$H_i$ preserves $\Xi_i^\sigma$, and every $H_j$ preserves
each of $G_i^{\sigma,L},G_i^{\sigma,R}$}, the latter being sums of tails.

\subsection{Spectral projections, exceptional difference fields, and
\texorpdfstring{$\Lambda_i$}{Lambda\_i}}\label{ssec:spectra}

Take an integral basis $(v_i,u_i)$ with $\det(v_i,u_i)=1$ and write
$z=sv_i+tu_i$; then $\pi_i(z)=t$ and $s=\det(z,u_i)$.  Call
$\{t=\text{const}\}$ a \emph{transverse row}.  Recall $H_i=\kappa_iv_i$.

\begin{lem}[Spectral projection]\label{lem:spectral}
Let $d\colon\Z^2\to\C$ satisfy $T^{H_i}d=d$.  For $\lambda^{\kappa_i}=1$ put
\[P_\lambda:=\frac1{\kappa_i}\sum_{k=0}^{\kappa_i-1}\lambda^{-k}\,T^{kv_i},
  \qquad
  \widehat d_t(\lambda):=\frac1{\kappa_i}\sum_{k=0}^{\kappa_i-1}\lambda^{-k}\,
  d(kv_i+tu_i).\]
\begin{enumerate}
\item $(P_\lambda d)(sv_i+tu_i)=\lambda^s\,\widehat d_t(\lambda)$, and
      $\sum_{\lambda}P_\lambda d=d$;
\item $T^{v_i}P_\lambda d=\lambda\,P_\lambda d$;
\item $P_\lambda$ is a Laurent polynomial in $T$, hence commutes with every
      $T^u$ and with every Laurent polynomial $f(T)$;
\item for $f(T)=\sum_{\alpha,\beta}f_{\alpha\beta}T^{\alpha v_i+\beta u_i}$,
\[\bigl(f(T)P_\lambda d\bigr)(sv_i+tu_i)
  =\lambda^{s}\sum_{\alpha,\beta}f_{\alpha\beta}\lambda^{\alpha}\,
    \widehat d_{t+\beta}(\lambda)
  =\lambda^s\bigl[f(\lambda,Y)\,\widehat d_{\bullet}(\lambda)\bigr](t),\]
      where $f(\lambda,Y):=\sum_{\alpha,\beta}f_{\alpha\beta}\lambda^\alpha
      Y^\beta$ acts, as a Laurent polynomial in the single variable $Y$, on
      functions of $t$, with $Y$ the shift $t\mapsto t+1$.
\end{enumerate}
\end{lem}

\begin{proof}
Parts 1 and 2 are the discrete Fourier decomposition on $\Z/\kappa_i$, carried
out row by row:
$(P_\lambda d)(sv_i+tu_i)=\frac1{\kappa_i}\sum_k\lambda^{-k}
 d((s+k)v_i+tu_i)=\lambda^s\cdot\frac1{\kappa_i}\sum_{k'}\lambda^{-k'}
 d(k'v_i+tu_i)$, re-indexing by $\kappa_i$-periodicity, while
$\sum_\lambda\lambda^{s-k}=\kappa_i\one[k\equiv s]$ gives completeness.  Part 3
is clear, and part 4 is a direct computation from part 1 and the definition of
$T^{\alpha v_i+\beta u_i}$.
\end{proof}

Say that $\lambda$ \emph{occurs in} $d$ if $P_\lambda d\ne0$, that is, if
$\widehat d_t(\lambda)\ne0$ for some row $t$.

Define the \emph{exceptional difference fields}
\[d_{i,\sigma,\epsilon,a}(z)=\one\bigl[\Xi_i^\sigma(z)=a\bigr]
  -\one\bigl[G_i^{\sigma,\epsilon}(z)=a\bigr]\in\{-1,0,1\}
  \qquad(\sigma\in\{+,-\},\ \epsilon\in\{L,R\},\ a\in\Fp).\]

\begin{itemize}
\item[\textup{(P1)}] $d_{i,\sigma,R,a}\equiv0$ on $\{\pi_i>r_i\}$, and
      $d_{i,\sigma,L,a}\equiv0$ on $\{\pi_i<\ell_i\}$.
\item[\textup{(P2)}] Both $\Xi_i^\sigma$ and $G_i^{\sigma,\epsilon}$ are
      preserved by $H_i$, so $T^{H_i}d_{i,\sigma,\epsilon,a}
      =d_{i,\sigma,\epsilon,a}$ and Lemma~\ref{lem:spectral} applies.
\end{itemize}

Put
\[\Lambda_i=\bigl\{\lambda:\ \lambda^{\kappa_i}=1,\ \exists\,\sigma,\epsilon,a
  \text{ such that }\lambda\text{ occurs in }d_{i,\sigma,\epsilon,a}\bigr\}.\]

\begin{lem}\label{lem:Lambda-nonempty}
$\Lambda_i\ne\varnothing$.
\end{lem}

\begin{proof}
Otherwise every $d_{i,\sigma,\epsilon,a}$ vanishes identically, by part 1 of
Lemma~\ref{lem:spectral}.  Taking $\sigma=+$ and $\epsilon=R$ gives
$\Xi_i^+=G_i^{+,R}$; subtracting $B_i^+$ gives $F_i=R_i$, so $F_i$ is doubly
periodic, contradicting \textup{(S2)}.
\end{proof}

Put
\begin{align}
A_i(t)&=\prod_{\lambda\in\Lambda_i}(t-\lambda),\qquad
  d_i=\deg A_i=|\Lambda_i|\ge1, \label{eq:Ai}\\
A(T)&=\prod_{i=1}^mA_i\bigl(T^{v_i}\bigr),\qquad
  Z=\sum_{i=1}^m[0,d_iv_i]\subset\R^2 . \label{eq:AZ}
\end{align}

\begin{itemize}
\item[\textup{(P3)}] For all $\sigma,\epsilon,a$ one has
      $A_i(T^{v_i})\,d_{i,\sigma,\epsilon,a}=0$.  Indeed, by
      Lemma~\ref{lem:spectral} only the terms with $\lambda\in\Lambda_i$ in
      $d=\sum_\lambda P_\lambda d$ are non-zero, and
      $A_i(T^{v_i})P_\lambda d=A_i(\lambda)P_\lambda d=0$.
\end{itemize}

The elements of $\Lambda_i$ are distinct $\kappa_i$-th roots of unity, so $A_i$
is monic with $A_i(0)=\prod_{\lambda}(-\lambda)\ne0$.  Hence
\begin{equation}
\Newt\bigl(A_i(T^{v_i})\bigr)=[0,d_iv_i],\qquad
\Newt(A)=\sum_i[0,d_iv_i]=Z, \label{eq:NewtA}
\end{equation}
the second equality by Minkowski additivity of Newton polygons (Ostrowski: over
an integral domain, $\Newt(fg)=\Newt(f)+\Newt(g)$).

\begin{remgalois}[Galois stability of $\Lambda_i$]
Each $d_{i,\sigma,\epsilon,a}$ is integer-valued, so for every
$\tau\in\Gal(\Q(\zeta_{\kappa_i})/\Q)$ one has
$\tau\bigl(\widehat d_t(\lambda)\bigr)=\widehat d_t(\tau\lambda)$.  Hence
$\lambda$ occurs if and only if $\tau\lambda$ occurs, $\Lambda_i$ is a union of
Galois orbits, $A_i$ is a product of distinct cyclotomic polynomials,
$A_i\in\Z[t]$ and $A\in\Z[T^{\pm}]$.  Nothing below depends on this
(Lemma~\ref{lem:budget} handles the coefficient field by invariance of rank
under field extension), but it does mean that every object of
\S\ref{sec:spectrum} can be discussed over $\Q$.
\end{remgalois}

\subsection{A spectrum-preserving encoding}\label{ssec:encoding}

\begin{lem}\label{lem:encoding}
There is an injection $w\colon\Fp\to\Z$ such that for every $i$ and every
$\lambda\in\Lambda_i$, the scalar difference field
\[\delta_{i,\sigma,\epsilon}:=w(\Xi_i^\sigma)-w\bigl(G_i^{\sigma,\epsilon}\bigr)
  =\sum_aw(a)\,d_{i,\sigma,\epsilon,a}\]
still has $\lambda$ occurring in it, for some pair $\sigma,\epsilon$.
\end{lem}

\begin{proof}
Regard $(w(a))_{a\in\Fp}$ as a vector in $\Q^p$.  For each $\lambda\in\Lambda_i$
there are, by definition, $\sigma,\epsilon,t$ for which the vector
$\bigl(\widehat{d_{i,\sigma,\epsilon,a}}\bigr)_t(\lambda)\in\C^p$, indexed by
$a$, is non-zero.  By linearity of the Fourier transform,
\[\bigl(\widehat{\delta_{i,\sigma,\epsilon}}\bigr)_t(\lambda)
  =\sum_aw(a)\,\bigl(\widehat{d_{i,\sigma,\epsilon,a}}\bigr)_t(\lambda)\]
is a non-zero complex linear form in $w$; its zero set inside $\Q^p$ is a proper
$\Q$-subspace of $\Q^p$, since otherwise the form would vanish on a $\Q$-basis
of $\Q^p$ and hence identically.  There are $\sum_i|\Lambda_i|$ such conditions.
Injectivity excludes a further $\binom p2$ proper hyperplanes
$\{w(a)=w(a')\}$.  Finitely many proper subspaces cannot cover $\Q^p$, so a
rational $w$ satisfying all the conditions exists; multiplying by a common
denominator makes it integer-valued without affecting non-vanishing or
injectivity.
\end{proof}

Fix such a $w$ and put $\eta=w(\theta)\colon\Z^2\to\Z$.  By \eqref{eq:caseB},
$D\eta=0$.

\subsection{One-sided vanishing}\label{ssec:onesided}

\begin{lem}\label{lem:onesided}
Let $d\colon\Z^2\to\C$ satisfy $T^{H_i}d=d$, vanish identically on a half-plane
$\{\pi_i>c\}$ or on a half-plane $\{\pi_i<c\}$, and let $\lambda$ occur in $d$.
If a Laurent polynomial $f$ satisfies $f(T)d=0$, then
$\bigl(T^{v_i}-\lambda\bigr)\mid f$.
\end{lem}

\begin{proof}
By part 3 of Lemma~\ref{lem:spectral}, $f(T)P_\lambda d=P_\lambda f(T)d=0$.
Write $c_\lambda(t)=\widehat d_t(\lambda)$, so $c_\lambda(t)=0$ whenever $d$
vanishes on row $t$, and $g(Y):=f(\lambda,Y)=\sum_\beta g_\beta Y^\beta$.  By
part 4 of Lemma~\ref{lem:spectral},
\[0=\bigl(f(T)P_\lambda d\bigr)(sv_i+tu_i)
  =\lambda^s\sum_\beta g_\beta\,c_\lambda(t+\beta)\qquad\text{for all }s,t,\]
that is, $g(Y)c_\lambda=0$; and $c_\lambda\not\equiv0$ because $\lambda$
occurs.

\emph{Case: $d$ vanishes on $\{\pi_i>c\}$.}  Then $c_\lambda(t)=0$ for $t>c$, so
$t_0:=\max\{t:c_\lambda(t)\ne0\}$ exists.  If $g\ne0$, let
$\beta_-=\min\{\beta:g_\beta\ne0\}$ and evaluate at $t=t_0-\beta_-$: for
$\beta>\beta_-$ one has $t_0-\beta_-+\beta>t_0$ and $c_\lambda$ vanishes there,
so only the term $\beta=\beta_-$ survives, giving
$g_{\beta_-}c_\lambda(t_0)=0$, a contradiction.

\emph{Case: $d$ vanishes on $\{\pi_i<c\}$.}  Symmetrically, take
$t_0=\min\{t:c_\lambda(t)\ne0\}$ and $\beta_+=\max\{\beta:g_\beta\ne0\}$ and
evaluate at $t=t_0-\beta_+$; only the term $\beta=\beta_+$ survives.

So $g\equiv0$.  Write $f$ as an element of $\C[X^{\pm},Y^{\pm}]$ with
$X=T^{v_i}$ and $Y=T^{u_i}$; since $(v_i,u_i)$ is a basis of $\Z^2$ this is an
automorphism of the Laurent ring.  With $f=\sum_\beta Y^\beta f_\beta(X)$, the
identity $g\equiv0$ says $f_\beta(\lambda)=0$ for every $\beta$, whence
$(X-\lambda)\mid f_\beta(X)$ — a one-variable Laurent polynomial vanishing at a
non-zero point is divisible by $X-\lambda$ — and therefore $(X-\lambda)\mid f$.
\end{proof}

\subsection{Divisibility}\label{ssec:divisibility}

\begin{lem}\label{lem:divisibility}
Let $f$ be a non-zero Laurent polynomial with complex coefficients such that
$f(T)\eta=c$ is constant.  Then $A\mid f$.
\end{lem}

\begin{proof}
\emph{(a) The relation passes to the orbit closure.}  The map
$z\mapsto\bigl(f(T)\eta\bigr)(z)-c$ depends on $\theta$ only through its values
on $z+\supp(f)$, so it is a local function of $\theta$.  It vanishes on every
translate of $\theta$; local functions are continuous for the product topology,
so it vanishes on every member $\chi$ of $X_\theta$: $f(T)w(\chi)=c$.

\emph{(b) Subtraction.}  Take $\chi=\Xi_i^\sigma$ and $\chi=G_i^{\sigma,\epsilon}$,
both in $X_\theta$ by \S\ref{ssec:isolating}, and subtract the two identities:
$f(T)\delta_{i,\sigma,\epsilon}=0$.

\emph{(c) One prime factor at a time.}  Fix $\lambda\in\Lambda_i$.  By
Lemma~\ref{lem:encoding} there are $\sigma,\epsilon$ with $\lambda$ occurring in
$\delta_{i,\sigma,\epsilon}$; by \textup{(P1)} this difference field vanishes on
a half-plane; by \textup{(P2)} it is preserved by $H_i$.  Lemma~\ref{lem:onesided}
gives $T^{v_i}-\lambda\mid f$.

\emph{(d) Assembling.}  Under the automorphism $X=T^{v_i}$, $Y=T^{u_i}$ the
element $T^{v_i}-\lambda$ becomes $X-\lambda$, and
$\C[X^{\pm},Y^{\pm}]/(X-\lambda)\cong\C[Y^{\pm}]$ is an integral domain, so
$T^{v_i}-\lambda$ is prime.  Primes coming from different directions $i\ne k$
are non-associate: being associate would force
$T^{v_i}-\lambda=cT^u(T^{v_k}-\mu)$, and comparing supports gives
$\{0,v_i\}=\{u,u+v_k\}$, hence $v_i=\pm v_k$, contradicting non-parallelism.
Primes from the same direction with $\lambda\ne\lambda'$ are also
non-associate: comparing supports forces $u=0$ and then comparing coefficients
forces $\lambda=\lambda'$.  In a unique factorisation domain, pairwise
non-associate primes each dividing $f$ have their product dividing $f$:
$A=\prod_i\prod_{\lambda\in\Lambda_i}(T^{v_i}-\lambda)\mid f$.
\end{proof}

\subsection{The dimension budget for affine relations}\label{ssec:budget}

For a lattice-convex $S$ put
\[R_Z(S)=\bigl\{r\in\Z^2:\ r+Z\subseteq\Conv(S)\bigr\},\]
\[U_S=\spn_{\Q}\Bigl\{\mathbf 1,\ \ \gamma\mapsto w\bigl(\gamma(s)\bigr)\ (s\in S)
      \Bigr\}\ \subseteq\ \{\cL\to\Q\}.\]

\begin{lem}\label{lem:budget}
$\dim_{\Q}U_S\ \ge\ |S|-|R_Z(S)|+1$.
\end{lem}

\begin{proof}
$U_S$ is generated by $|S|+1$ elements, so
$\dim_{\Q}U_S=(|S|+1)-\dim_{\Q}\cR$, where
\[\cR=\Bigl\{(c_0,(c_s)_{s\in S})\in\Q^{1+|S|}:\
   \sum_{s\in S}c_s\,\eta(u+s)=c_0\ \text{ for all }u\in\Z^2\Bigr\}.\]
$\cR$ is the $\Q$-solution space of a (possibly infinite) homogeneous linear
system with integer coefficients; writing $\cR_{\C}$ for its $\C$-solution
space, $\cR_{\C}=\cR\otimes_{\Q}\C$ and
$\dim_{\C}\cR_{\C}=\dim_{\Q}\cR$, the dimension of the solution space of a
rational linear system being unchanged by extension of the base field.

Let $(c_0,(c_s))\in\cR_{\C}$ be non-zero and put $f=\sum_sc_sT^s$.  If $f=0$
then $c_0=0$ and the element is zero; so $f\ne0$, and $f(T)\eta=c_0$ with
$\supp(f)\subseteq S$.  By Lemma~\ref{lem:divisibility}, valid for complex $f$,
we may write $f=Ag$ with $g\ne0$.  By \eqref{eq:NewtA} and Newton-polygon
additivity,
\[\Newt(f)=\Newt(A)+\Newt(g)=Z+\Newt(g).\]
From $\supp(f)\subseteq S$ we get $\Newt(f)\subseteq\Conv(S)$, so for every
$r\in\supp(g)$,
\[r+Z\subseteq\Newt(g)+Z=\Newt(f)\subseteq\Conv(S),\]
that is, $r\in R_Z(S)$.  Hence $\supp(g)\subseteq R_Z(S)$.

The map $\cR_{\C}\to\{g:\supp(g)\subseteq R_Z(S)\}$, $(c_0,f)\mapsto g=f/A$, is
$\C$-linear and injective ($c_0$ is determined by $f$, $A\ne0$, and division in
an integral domain is unique), so
\[\dim_{\Q}\cR=\dim_{\C}\cR_{\C}\le|R_Z(S)| .\qedhere\]
\end{proof}

\begin{rem}\label{rem:empty-RZ}
By lattice-convexity, $r\in R_Z(S)$ implies
$(r+Z)\cap\Z^2\subseteq\Conv(S)\cap\Z^2=S$.  If $R_Z(S)=\varnothing$, then
Lemma~\ref{lem:budget} gives $\dim U_S\ge|S|+1$; but $U_S$ is a space of
functions on the finite set $\cL$, so
$P_\theta(S)=|\cL|\ge\dim U_S\ge|S|+1$ and the main theorem holds.  This case
covers point and segment windows, since $Z$ has non-empty interior ($m\ge2$ and
$v_1\not\parallel v_2$).  \emph{From now on we assume $R_Z(S)\ne\varnothing$.}
\end{rem}

\section{\texorpdfstring{$A$}{A} applied to the colour field yields a doubly
periodic background}\label{sec:background}

\setcounter{thm}{-1}
\subsection{Realised sectors and adjacent sectors}\label{ssec:sectors}

The complement $\R^2\setminus\bigcup_i\R v_i$ has exactly $2m$ connected
components, open cones: order the $2m$ rays $\{\sigma v_i:\sigma=\pm,\,i\in[m]\}$
by angle; the components are the \emph{sectors} between consecutive rays.  On
each sector every $\pi_i$ has constant sign; write $\epsilon(K)\in\{+,-\}^m$ for
the sign vector of the sector $K$, and call these $2m$ sign vectors
\emph{realised}.  (For $m\ge3$ the remaining $2^m-2m$ sign vectors are not
realised.)  For a realised $\epsilon$ put
\[\Theta_\epsilon:=\sum_i\varepsilon_i^{\epsilon_i},\qquad
  \varepsilon_i^{+}=R_i,\quad \varepsilon_i^{-}=L_i,\]
which is doubly periodic and preserved by every $H_j$, as is its colour field
$e_{\Theta_\epsilon}=(\one[\Theta_\epsilon=a])_a$.

\begin{lem}[Sectors adjacent along a ray]\label{lem:adjacent}
Let the ray $\sigma v_i$ be the common boundary ray of two sectors $K,K'$.
Then $\epsilon(K)$ and $\epsilon(K')$ differ exactly in the coordinate $i$, and
for $j\ne i$ one has
$\epsilon(K)_j=\epsilon(K')_j=\sgn\bigl(\sigma\pi_j(v_i)\bigr)$.  Consequently
\[\{\Theta_{\epsilon(K)},\Theta_{\epsilon(K')}\}=\{G_i^{\sigma,L},\ G_i^{\sigma,R}\}.\]
In the angular order, the $2m$ sectors form a $2m$-cycle under the relation of
sharing a boundary ray.
\end{lem}

\begin{proof}
The ray $\sigma v_i$ lies on $\R v_i$ and on no $\R v_j$ with $j\ne i$, the
directions being pairwise non-parallel.  A small circular arc centred at
$\sigma v_i$ crosses only $\R v_i$, so the sign vectors of the two sectors on
either side differ only in coordinate $i$; for $j\ne i$, the sign of $\pi_j$ is
constant on the arc and equals
$\sgn\pi_j(\sigma v_i)=\sgn(\sigma\pi_j(v_i))$.  By \eqref{eq:tailchoice},
$\varepsilon_j^{\epsilon(K)_j}=\varepsilon_j^\sigma$ is precisely the tail
choice made in $B_i^\sigma$, so $\Theta_{\epsilon(K)}$ and
$\Theta_{\epsilon(K')}$ are $R_i+B_i^\sigma$ and $L_i+B_i^\sigma$ in the order
dictated by $\epsilon_i$.  The cycle structure comes from the angular ordering.
\end{proof}

\begin{lem}\label{lem:Ae-sectors}
For every ray $\sigma v_i$,
\[A(T)\,e_{G_i^{\sigma,R}}=A(T)\,e_{G_i^{\sigma,L}} .\]
Consequently $A(T)e_{\Theta_\epsilon}$ takes one and the same value for all
\emph{realised} $\epsilon$; write
\[b=(b_a)_{a\in\Fp}:=A(T)e_{\Theta_\epsilon}.\]
This $b$ is doubly periodic and preserved by every $H_j$.
\end{lem}

\begin{proof}
Colour by colour,
\[\one[G_i^{\sigma,R}=a]-\one[G_i^{\sigma,L}=a]=d_{i,\sigma,L,a}-d_{i,\sigma,R,a},\]
which is killed by $A_i(T^{v_i})$ by \textup{(P3)}, hence by $A$.  By
Lemma~\ref{lem:adjacent} the backgrounds of two sectors adjacent along a ray are
exactly $G_i^{\sigma,R}$ and $G_i^{\sigma,L}$, so their values of $Ae$ agree;
and the $2m$ sectors form a cycle under this adjacency, hence a connected graph,
so $Ae_{\Theta_\epsilon}$ is the same for all realised sectors.  Double
periodicity and $H_j$-invariance follow from the corresponding properties of
$\Theta_\epsilon$ together with the fact that $A$ commutes with translations.
\end{proof}

\begin{remnogen}[The lemma does not extend to unrealised tail
choices]
The generalisation ``$Ae_{\Theta_\varepsilon}=Ae_{\Theta_{\varepsilon'}}$ for
all $\varepsilon,\varepsilon'\in\{L,R\}^m$, realised or not'' is \emph{false}.
The reason is that $\Lambda_i$ is defined using only the backgrounds $B_i^\pm$,
whereas a mixed tail choice $B$ may have periods along $v_i$ that $B_i^\pm$ do
not have, so the spectrum of $\one[R_i+B=a]-\one[L_i+B=a]$ can fall outside
$\Lambda_i$.  An explicit counterexample is in Appendix~\ref{app:counterexample}.
Lemma~\ref{lem:Ae-finite} uses only realised sectors, so the main line is
unaffected.
\end{remnogen}

\subsection{\texorpdfstring{$Ae_\theta-b$}{Ae-b} has finite
support}\label{ssec:Ae-finite}

\begin{lem}\label{lem:Ae-finite}
$Ae_\theta-b$ has finite support.
\end{lem}

\begin{proof}
Write $M_A=\supp(A)$.  Choose $c_i\ge\max(|\ell_i|,|r_i|)$ so that
$\{\pi_i<-c_i\}\subseteq\{\pi_i<\ell_i\}$ and
$\{\pi_i>c_i\}\subseteq\{\pi_i>r_i\}$, and put
\[\rho_i=c_i+\max_{s\in M_A}|\pi_i(s)|,\qquad \Sigma_i=\{z:|\pi_i(z)|\le\rho_i\}.\]
Call $i$ \emph{active at $z$} if $z\in\Sigma_i$; otherwise all of $z+M_A$ lies
in one and the same tail region of component $i$, the one on the side
$\sgn\pi_i(z)$.

\emph{(a) No active direction.}  Then $|\pi_i(z)|>\rho_i>0$ for every $i$; put
$\epsilon_i=\sgn\pi_i(z)$.  Regarded as a point of $\R^2$, $z$ lies on no
$\R v_i$ and has sign vector $\epsilon$, so $\epsilon$ is realised; and
$\theta=\Theta_\epsilon$ on $z+M_A$.  Hence
$(Ae_\theta)(z)=(Ae_{\Theta_\epsilon})(z)=b(z)$.

\emph{(b) At least two active directions.}  Then $z\in\Sigma_i\cap\Sigma_j$ with
$i\ne j$, an intersection of two strips of bounded width in non-parallel
directions, which is bounded and so contains only finitely many lattice points.

\emph{(c) Exactly one active direction $i$.}  Write $z=tv_i+qu_i$ in the basis
$(v_i,u_i)$, so $\pi_i(z)=q$ and $|q|\le\rho_i$.  For $j\ne i$,
\[\pi_j(z)=t\,\pi_j(v_i)+q\,\pi_j(u_i),\qquad \pi_j(v_i)=\det(v_j,v_i)\ne0 .\]
If
\begin{equation}
|t|\ >\ \frac{\rho_i\max_{j\ne i}|\pi_j(u_i)|+\max_{j\ne i}\rho_j}
             {\min_{j\ne i}|\pi_j(v_i)|}, \label{eq:tbound}
\end{equation}
then for every $j\ne i$ and every $s\in M_A$,
\[|\pi_j(z+s)|\ \ge\ |t|\,|\pi_j(v_i)|-|q|\,|\pi_j(u_i)|-|\pi_j(s)|
  \ >\ \rho_j-|\pi_j(s)|\ \ge\ \rho_j-\max_{s'\in M_A}|\pi_j(s')|=c_j,\]
and, because $|t|\,|\pi_j(v_i)|>|q|\,|\pi_j(u_i)|+|\pi_j(s)|$, also
$\sgn\pi_j(z+s)=\sgn\bigl(t\,\pi_j(v_i)\bigr)$.  So all of $z+M_A$ lies in the
tail region of component $j$ determined by $\sgn(t\pi_j(v_i))$.  By
\eqref{eq:tailchoice} these are exactly the tail choices made in
$B_i^{\sgn(t)}$, so on $z+M_A$
\[\theta=F_i+B_i^{\sgn(t)}=\Xi_i^{\sigma},\qquad\sigma=\sgn(t),\]
whence $(Ae_\theta)(z)=\bigl(Ae_{\Xi_i^{\sigma}}\bigr)(z)$.  By \textup{(P3)},
$A_i(T^{v_i})$ kills
$e_{\Xi_i^\sigma}-e_{G_i^{\sigma,R}}=\bigl(d_{i,\sigma,R,a}\bigr)_a$, so
\[\bigl(Ae_{\Xi_i^\sigma}\bigr)(z)=\bigl(Ae_{G_i^{\sigma,R}}\bigr)(z)=b(z),\]
the last step because $G_i^{\sigma,R}$ is the background of one of the realised
sectors adjacent along the ray $\sigma v_i$ (Lemma~\ref{lem:adjacent}).

Only finitely many $t$ violate \eqref{eq:tbound}, and these combine with
$|q|\le\rho_i$ to give finitely many lattice points.  Combining (a), (b), (c),
$\supp(Ae_\theta-b)$ is finite.
\end{proof}

\subsection{Conclusion}

\begin{prop}\label{prop:Ae-b}
$A(T)e_\theta=b$.  In particular
\[A\eta=b_\eta:=\sum_aw(a)b_a\]
is a doubly periodic field preserved by every $H_i$.
\end{prop}

\begin{proof}
By \eqref{eq:caseB}, $De_\theta=0$; and $b$ is preserved by every $H_i$, so
$Db=0$.  Hence $D(Ae_\theta-b)=A\,De_\theta-Db=0$.  By
Lemma~\ref{lem:Ae-finite} this function has finite support, so
Lemma~\ref{lem:domain} (with $D\ne0$) forces it to vanish.
\end{proof}

\section{At least one two-point difference is non-zero}\label{sec:twopoint}

Put $Q_i=\prod_{j\ne i}\bigl(T^{H_j}-1\bigr)$ as in \eqref{eq:Qexp}, and
$f_i=Q_i\,\eta$.  From $D=(T^{H_i}-1)Q_i$ and $D\eta=0$ it follows that $f_i$
has period $H_i$.

\begin{lem}\label{lem:fi}
The support of $f_i$ is contained in a strip of bounded width in the direction
$v_i$, and $f_i\ne0$.
\end{lem}

\begin{proof}
\emph{(a) A global identity.}  We have $f_i=T^{nH_i}f_i$ for every $n$.  Let
$n\to+\infty$: on every fixed finite set $T^{nH_i}\theta\to\Xi_i^+$, the
convergence being eventual pointwise stabilisation, so
$T^{nH_i}\eta\to w(\Xi_i^+)$ and therefore
\begin{equation}
f_i=Q_i\,w\bigl(\Xi_i^+\bigr). \label{eq:fi-Xi}
\end{equation}
Each factor $T^{H_j}-1$ with $j\ne i$ annihilates $w(G_i^{+,R})$, since $H_j$
preserves that pure-tail field, so
\begin{equation}
f_i=Q_i\,\delta_i,\qquad
\delta_i:=w(\Xi_i^+)-w\bigl(G_i^{+,R}\bigr). \label{eq:fi-delta}
\end{equation}
Using $G_i^{+,L}$ instead gives $f_i=Q_i\,\delta_i'$ with
$\delta_i'=w(\Xi_i^+)-w(G_i^{+,L})$.

\emph{(b) A strip of bounded width.}  By \textup{(P1)}, $\delta_i\equiv0$ on
$\{\pi_i>r_i\}$ and $\delta_i'\equiv0$ on $\{\pi_i<\ell_i\}$.  Put
$c=\max_{s\in\cE_i}|\pi_i(s)|$.  By \eqref{eq:fi-delta} and \eqref{eq:Qexp},
$f_i(z)\ne0$ forces $\delta_i(z+s)\ne0$ for some $s\in\cE_i$, whence
$\pi_i(z)\le r_i+c$; the other identity gives $\pi_i(z)\ge\ell_i-c$.  So
$\supp(f_i)\subseteq\{\ell_i-c\le\pi_i\le r_i+c\}$.

\emph{(c) Non-vanishing.}  First, $\delta_i\ne0$: otherwise
$w(\Xi_i^+)=w(G_i^{+,R})$, and injectivity of $w$ gives $\Xi_i^+=G_i^{+,R}$;
subtracting $B_i^+$ gives $F_i=R_i$, which is doubly periodic, contradicting
\textup{(S2)}.

Since $\delta_i$ vanishes on $\{\pi_i>r_i\}$ but is not identically zero,
\[\pi^*=\max\{\pi_i(z):\delta_i(z)\ne0\}\ (\le r_i)\]
exists.  Choose $z^*$ with $\delta_i(z^*)\ne0$ and $\pi_i(z^*)=\pi^*$.

Expand $Q_i$ formally as in \eqref{eq:Qexp}.  We have
$\pi_i(H_C)=\sum_{j\in C}\pi_i(H_j)$ with every $\pi_i(H_j)\ne0$ for $j\ne i$.
Put
\[C_{\min}=\{j\ne i:\pi_i(H_j)<0\}.\]
For any $C\subseteq[m]\setminus\{i\}$ with $C\ne C_{\min}$,
\[\pi_i\bigl(H_C\bigr)-\pi_i\bigl(H_{C_{\min}}\bigr)
  =\sum_{j\in C\setminus C_{\min}}\pi_i(H_j)
   -\sum_{j\in C_{\min}\setminus C}\pi_i(H_j)>0,\]
because the terms of the first sum are positive, those of the second are
negative, and at least one of the two sums is non-empty.  So
$\pi_i(H_{C_{\min}})$ is the \emph{unique} minimum.  This is a statement about
the $\pi_i$-projections and is unaffected by any coincidences among the points
$H_C$ themselves.

Evaluate at $z=z^*-H_{C_{\min}}$:
\[f_i(z)=\sum_{C}(-1)^{m-1-|C|}\delta_i\bigl(z+H_C\bigr).\]
The term $C=C_{\min}$ contributes $\pm\delta_i(z^*)\ne0$; for $C\ne C_{\min}$,
\[\pi_i\bigl(z+H_C\bigr)=\pi^*+\Bigl(\pi_i(H_C)-\pi_i(H_{C_{\min}})\Bigr)>\pi^*,\]
so $\delta_i$ vanishes there.  Hence $f_i(z)=\pm\delta_i(z^*)\ne0$.
\end{proof}

\begin{lem}\label{lem:J-nonzero}
For every $d$ the function $J(d,\cdot):=D_z\bigl[\eta(z)\eta(z+d)\bigr]$ has
finite support; and there exists $d$ with $J(d,\cdot)\ne0$.
\end{lem}

\begin{proof}
\emph{Finite support:} $z\mapsto\eta(z)\eta(z+d)$ is a local function of
$\theta$, with window $W=\{0,d\}$; apply Lemma~\ref{lem:finite-support}.

\emph{Non-vanishing:} choose $i\ne j$, possible since $m\ge2$.  By
Lemma~\ref{lem:fi} choose $x_i,x_j$ with $f_i(x_i)\ne0$ and $f_j(x_j)\ne0$.  Put
$d_0=x_j-x_i$ and
\[C(z)=f_i(z)\,f_j(z+d_0).\]
Then $C(x_i)=f_i(x_i)f_j(x_j)\ne0$.  Now $\supp(f_i)$ lies in a strip of bounded
width in the direction $v_i$ and $\supp\bigl(f_j(\cdot+d_0)\bigr)$ in one of
direction $v_j$, with $v_i\not\parallel v_j$, so $\supp(C)$ is bounded.  By
Lemma~\ref{lem:domain}, $DC\ne0$.

On the other hand, expanding formally by \eqref{eq:Qexp},
\[C(z)=\Bigl(\sum_{C_1}\pm\eta\bigl(z+H_{C_1}\bigr)\Bigr)
       \Bigl(\sum_{C_2}\pm\eta\bigl(z+d_0+H_{C_2}\bigr)\Bigr)
      =\sum_{C_1,C_2}\pm\,\eta(z+u_1)\eta(z+u_2),\]
with $u_1=H_{C_1}$ and $u_2=d_0+H_{C_2}$.  By translation invariance of $D$,
\[DC=\sum_{C_1,C_2}\pm\,T^{u_1}\Bigl(D_z\bigl[\eta(z)\eta(z+u_2-u_1)\bigr]\Bigr)
    =\sum_{C_1,C_2}\pm\,T^{u_1}J\bigl(u_2-u_1,\cdot\bigr).\]
If every $J(d,\cdot)$ vanished, then $DC=0$, a contradiction.
\end{proof}

\section{The bi-recursion and the quotient algebra}\label{sec:quotient}

\subsection{The bi-recursion}

Regard $J$ as a function of the pair $(d,z)$:
$J(d,z)=D_z\bigl[\eta(z)\eta(z+d)\bigr]$.  Write $T_d^sJ(d,z)=J(d+s,z)$ and
$T_z^sJ(d,z)=J(d,z+s)$.

\begin{lem}\label{lem:birecursion}
\begin{equation}
A(T_d)\,J=0,\qquad A\bigl(T_zT_d^{-1}\bigr)\,J=0 . \label{eq:birec}
\end{equation}
\end{lem}

\begin{proof}
\emph{First identity.}  For any Laurent polynomial $g=\sum_sg_sT^s$,
\[g(T_d)J(d,z)=\sum_sg_s\,D_z\bigl[\eta(z)\eta(z+d+s)\bigr]
             =D_z\Bigl[\eta(z)\bigl(g(T)\eta\bigr)(z+d)\Bigr].\]
Take $g=A$; by Proposition~\ref{prop:Ae-b}, $(A\eta)(z+d)=b_\eta(z+d)$.
Expanding $D_z$ by \eqref{eq:Dexp},
\[D_z\bigl[\eta(z)b_\eta(z+d)\bigr]
  =\sum_C(-1)^{m-|C|}\,\eta\bigl(z+H_C\bigr)\,b_\eta\bigl(z+H_C+d\bigr).\]
Each $H_C$ is a sum of vectors $H_i$ and so preserves $b_\eta$; therefore
$b_\eta(z+H_C+d)=b_\eta(z+d)$ and
\[=b_\eta(z+d)\sum_C(-1)^{m-|C|}\,\eta\bigl(z+H_C\bigr)
  =b_\eta(z+d)\,\bigl(D\eta\bigr)(z)=0 .\]

\emph{Second identity.}  The operator $\bigl(T_zT_d^{-1}\bigr)^s$ sends $(d,z)$
to $(d-s,z+s)$: the first observation point moves, $z\mapsto z+s$, while the
second, $z+d$, does not.  Hence
\begin{align*}
g\bigl(T_zT_d^{-1}\bigr)J(d,z)
 &=\sum_sg_s\,D_{z'}\bigl[\eta(z')\eta(z'+d-s)\bigr]\Big|_{z'=z+s}\\
 &=\sum_C(-1)^{m-|C|}\sum_sg_s\,\eta\bigl(z+s+H_C\bigr)\,
    \eta\bigl(z+H_C+d\bigr)\\
 &=\sum_C(-1)^{m-|C|}\bigl(g(T)\eta\bigr)\bigl(z+H_C\bigr)\,
    \eta\bigl(z+H_C+d\bigr).
\end{align*}
Take $g=A$: then $(A\eta)(z+H_C)=b_\eta(z+H_C)=b_\eta(z)$, so the sum equals
\[b_\eta(z)\sum_C(-1)^{m-|C|}\,\eta\bigl(z+H_C+d\bigr)
 =b_\eta(z)\cdot\bigl(D\eta\bigr)(z+d)=0 .\qedhere\]
\end{proof}

\subsection{The quotient algebra}\label{ssec:quotient}

Let $X=(X_1,X_2)$ be \emph{independent indeterminates}, $K=\C(X_1,X_2)$, and
$R=K[Y_1^{\pm1},Y_2^{\pm1}]$.  For $u\in\Z^2$ write $Y^u=Y_1^{u_1}Y_2^{u_2}$ and
$X^u=X_1^{u_1}X_2^{u_2}$; the monomials $\{Y^u\}_{u\in\Z^2}$ are a $K$-basis of
$R$.  Put
\[a_i=A_i\bigl(Y^{v_i}\bigr),\quad c_i=A_i\bigl(X^{v_i}Y^{-v_i}\bigr),\quad
  a=\prod_ia_i=A(Y),\quad c=\prod_ic_i=A(X/Y).\]

\begin{lem}\label{lem:quotient}
As a $K$-vector space, $R/(a,c)$ is spanned by the images of
$\{Y^d:d\in(Z-Z)\cap\Z^2\}$.
\end{lem}

\begin{proof}
\textbf{Step 1: no common zeros, and the Chinese remainder theorem.}  Extend to
$\overline K$ and test in $(\overline K^\times)^2$.  Note that $\C$ is
algebraically closed, so an element of $\overline K$ that is algebraic over $\C$
lies in $\C$, whereas $X^{u}$ with $u\ne0$ is transcendental over $\C$.

\emph{(i) $a_i$ and $c_i$ have no common zero.}  From $a_i=0$ we get
$Y^{v_i}=\lambda\in\Lambda_i$; from $c_i=0$ we get
$X^{v_i}Y^{-v_i}=\mu\in\Lambda_i$.  Multiplying, $X^{v_i}=\lambda\mu\in\C$, a
contradiction.

\emph{(ii) If $a_i=a_k=0$ with $i\ne k$, then every $c_j\ne0$.}  From
$Y^{v_i}=\lambda$, $Y^{v_k}=\lambda'$ and $\Delta=\det(v_i,v_k)\ne0$, the
elements $Y_1^{\Delta}$ and $Y_2^{\Delta}$ are products of integer powers of
$\lambda,\lambda'$, so $Y_1,Y_2$ are algebraic over $\C$ and hence lie in
$\C^\times$.  If some $c_j=0$ then $X^{v_j}=\mu Y^{v_j}\in\C$, a contradiction.

\emph{(iii) If $c_j=c_l=0$ with $j\ne l$, then every $a_i\ne0$.}  Put $W=X/Y$,
that is $W_1=X_1/Y_1$ and $W_2=X_2/Y_2$.  From $W^{v_j}=\mu$ and
$W^{v_l}=\mu'$ we get, as in (ii), $W_1,W_2\in\C^\times$, that is
$Y_1=\alpha X_1$ and $Y_2=\beta X_2$ with $\alpha,\beta\in\C^\times$.  If some
$a_i=0$ then $Y^{v_i}=\lambda$, that is
$\alpha^{v_{i1}}\beta^{v_{i2}}X^{v_i}=\lambda$ and $X^{v_i}\in\C$, a
contradiction.

By (ii), for $i\ne k$ the ideal $(a_i,a_k,c)$ has no zero in
$(\overline K^\times)^2$, so it equals $(1)$ in $\overline K[Y^{\pm}]$ by the
Nullstellensatz, and descends to $R$ by faithful flatness of
$R\to\overline K[Y^\pm]$: $(a_i,a_k,c)=R$.  That is, the $a_i$ are pairwise
coprime in $R/(c)$.  By the Chinese remainder theorem,
\[R/(a,c)=R\big/\bigl(\textstyle\prod_ia_i,\ c\bigr)\ \simeq\ \prod_iR/(a_i,c).\]
In $R/(a_i)$: by (i), $(a_i,c_i)=R$, so $c_i$ is a unit; by (iii), for
$j\ne l$ with $j,l\ne i$ one has $(a_i,c_j,c_l)=R$, so the $c_j$ are pairwise
coprime in $R/(a_i)$.  A second application of the Chinese remainder theorem
gives
\[R/(a_i,c)=R\big/\bigl(a_i,\ \textstyle\prod_jc_j\bigr)
  \simeq\prod_{j\ne i}R/(a_i,c_j),\]
and combining, with the isomorphism given by the natural projections,
\begin{equation}
R/(a,c)\ \simeq\ \prod_{i\ne j}R/(a_i,c_j). \label{eq:CRT}
\end{equation}

\textbf{Step 2: a monomial spanning set for each factor.}  Fix $i\ne j$.  Put
$\su=Y^{v_i}$ and $\sw=Y^{-v_j}$, let
$\mathfrak L=\Z v_i+\Z(-v_j)$ and $\nu=[\Z^2:\mathfrak L]=|\det(v_i,v_j)|$.
Take the $\nu$ lattice points $r_1,\dots,r_\nu$ of the half-open fundamental
domain
\[\Pi=\{\alpha v_i-\beta v_j:\ 0\le\alpha<1,\ 0\le\beta<1\},\]
a complete set of representatives for $\Z^2/\mathfrak L$.  Then $R$ is free as a
$K[\su^{\pm},\sw^{\pm}]$-module with basis $\{Y^{r_n}\}$.

The ideal $(a_i,c_j)$ is generated by the elements $a_i=A_i(\su)$ and
$c_j=A_j\bigl(X^{v_j}\sw\bigr)$ of the subring $K[\su^\pm,\sw^\pm]$, so
\[R/(a_i,c_j)\simeq\bigoplus_{n=1}^{\nu}Y^{r_n}\cdot
  K[\su^\pm,\sw^\pm]\big/\bigl(A_i(\su),A_j(X^{v_j}\sw)\bigr).\]
Now $A_i(\su)$ has degree $d_i$ in $\su$ with both extreme coefficients non-zero
($A_i$ is monic and $A_i(0)\ne0$), so any integer power of $\su$ reduces into
the $K$-span of $\{\su^0,\dots,\su^{d_i-1}\}$; and $A_j(X^{v_j}\sw)$ has degree
$d_j$ in $\sw$ with leading coefficient $X^{d_jv_j}\ne0$ and constant term
$A_j(0)\ne0$, so any integer power of $\sw$ reduces into the span of
$\{\sw^0,\dots,\sw^{d_j-1}\}$.  Hence
\begin{equation}
\Bigl\{Y^{\,r_n+\alpha v_i-\beta v_j}:\ 1\le n\le \nu,\ 0\le\alpha<d_i,\
0\le\beta<d_j\Bigr\} \label{eq:spanning}
\end{equation}
spans $R/(a_i,c_j)$.  Since $r_n\in\Pi$ — its $v_i$-coefficient and its
$(-v_j)$-coefficient both lie in $[0,1)$ — all the exponents in
\eqref{eq:spanning} lie in
\begin{equation}
P_{ij}:=[0,d_iv_i]+[0,-d_jv_j]. \label{eq:Pij}
\end{equation}

\textbf{Step 3: projection onto a factor.}  Put
\[E_{ij}=\prod_{k\ne i}a_k\ \prod_{l\ne j}c_l .\]
In the decomposition \eqref{eq:CRT}:
\begin{itemize}
\item on a factor $(i',j')$ with $i'\ne i$: the factor $a_{i'}$ occurs in
      $E_{ij}$ and is $\equiv0$ there, so $E_{ij}\equiv0$;
\item on a factor $(i,j')$ with $j'\ne j$: the factor $c_{j'}$ occurs and is
      $\equiv0$, so $E_{ij}\equiv0$;
\item on the factor $(i,j)$: by (i), (ii), (iii), the ideals $(a_i,c_j,a_k)$ for
      $k\ne i$ and $(a_i,c_j,c_l)$ for $l\ne j$ have no zero in
      $(\overline K^\times)^2$, hence equal $R$ by the Nullstellensatz and
      faithfully flat descent.  Now $(a_i,c_j,x)=R$ means precisely that there
      is $w\in R$ with $wx\equiv1\pmod{(a_i,c_j)}$, so each $a_k$ and each $c_l$
      is a unit in $R/(a_i,c_j)$, and so is their product $E_{ij}$.
\end{itemize}

Therefore the image of $E_{ij}\cdot\eqref{eq:spanning}$ spans the factor
$(i,j)$ and vanishes on all the others.  Every monomial exponent occurring in
$E_{ij}\cdot Y^{e}$ lies in $\Newt(E_{ij})+e$, and
\begin{align*}
\Newt(E_{ij})+P_{ij}
 &\subseteq\sum_{k\ne i}[0,d_kv_k]+\sum_{l\ne j}[0,-d_lv_l]+P_{ij}\\
 &=\sum_{k}[0,d_kv_k]+\sum_{l}[0,-d_lv_l]=Z+(-Z)=Z-Z,
\end{align*}
using $\Newt(a_k)=[0,d_kv_k]$, $\Newt(c_l)=[0,-d_lv_l]$ and the fact that the
Newton polygon of a product is the Minkowski sum.  Taking the union over all
$i\ne j$ gives the spanning statement.
\end{proof}

\subsection{The witness lies in
\texorpdfstring{$Z-Z$}{Z-Z}}\label{ssec:witness}

\begin{prop}\label{prop:witness}
There is $d\in(Z-Z)\cap\Z^2$ with $d\ne0$ and $J(d,\cdot)\ne0$.
\end{prop}

\begin{proof}
By Lemma~\ref{lem:J-nonzero} each $J(d,\cdot)$ has finite support, so
\[\widehat J(d;X)=\sum_zJ(d,z)X^{-z}\ \in\ \C[X_1^{\pm},X_2^{\pm}]\subset K\]
is well defined, the supports for different $d$ being allowed to differ.  Define
a \emph{$K$-linear} functional $L\colon R\to K$ by setting
$L(Y^d)=\widehat J(d;X)$ on the $K$-basis and extending $K$-linearly.  That $L$
is $K$-linear, and not merely $\C$-linear, is what makes the next step work:
coefficients $X^s\in K$ may be pulled out of $L$.

\smallskip\noindent\emph{$L$ kills the ideal $(a,c)$.}  Write $A=\sum_sA_sT^s$,
so that $a=A(Y)=\sum_sA_sY^s$ and $c=A(X/Y)=\sum_sA_sX^sY^{-s}$.

By the first identity of \eqref{eq:birec}, for every $(d,z)$ we have
$\sum_sA_s\,J(d+s,z)=0$.  Taking the Laurent transform in $z$,
\[\sum_sA_s\,\widehat J(d+s;X)=0
  \ \Longrightarrow\ L\bigl(Y^{d}\cdot a\bigr)=\sum_sA_sL\bigl(Y^{d+s}\bigr)=0 .\]

By the second identity of \eqref{eq:birec}, for every $(d,z)$ we have
$\sum_sA_s\,J(d-s,z+s)=0$.  Taking the Laurent transform in $z$ and noting
\[\sum_zJ(d-s,z+s)X^{-z}=\sum_{z'}J(d-s,z')X^{-z'+s}=X^{s}\widehat J(d-s;X),\]
we get $\sum_sA_sX^s\widehat J(d-s;X)=0$, that is
$L\bigl(Y^d\cdot c\bigr)=\sum_sA_sX^sL(Y^{d-s})=0$.

For an arbitrary $g=\sum_dg_dY^d\in R$ with $g_d\in K$, $K$-linearity gives
$L(ga)=\sum_dg_dL(Y^da)=0$, and likewise $L(gc)=0$.  So $L$ vanishes on
$(a,c)$ and descends to $\bar L\colon R/(a,c)\to K$.

\smallskip\noindent\emph{Conclusion.}  If $J(d,\cdot)=0$ for every
$d\in(Z-Z)\cap\Z^2$, then $\bar L$ vanishes on a spanning set supplied by
Lemma~\ref{lem:quotient}, hence $\bar L=0$, hence $L=0$, hence
$\widehat J(d;X)=0$ for \emph{every} $d$, that is $J\equiv0$ — contradicting
Lemma~\ref{lem:J-nonzero}.

\smallskip\noindent\emph{$d\ne0$.}  We have $J(0,z)=D_z[\eta(z)^2]$, and
$\eta^2=w'(\theta)$ with $w'(a)=w(a)^2$ is a single-site colour function, so
$D\eta^2=0$ by \eqref{eq:caseB}.
\end{proof}

\section{Putting two points into one integral zonotope}\label{sec:zonotope}

\begin{lem}\label{lem:zonotope}
Let $Z=\sum_{i=1}^m[0,d_iv_i]$ with integers $d_i\ge1$, $v_i\in\Z^2$, $m\ge2$
and $v_1\not\parallel v_2$.  Then
\begin{equation}
(Z-Z)\cap\Z^2\ =\ \bigl(Z\cap\Z^2\bigr)-\bigl(Z\cap\Z^2\bigr).
\label{eq:ZZ}
\end{equation}
\end{lem}

\begin{proof}
The inclusion ``$\supseteq$'' is clear.

``$\subseteq$'': put $c_0=\sum_id_iv_i\in\Z^2$.  For $x=\sum_ix_i\in Z$ with
$x_i\in[0,d_iv_i]$ we have $c_0-x=\sum_i(d_iv_i-x_i)$ with
$d_iv_i-x_i\in[0,d_iv_i]$, so $c_0-x\in Z$; that is, $c_0-Z\subseteq Z$.
Applying this once more, $Z=c_0-(c_0-Z)\subseteq c_0-Z$, so $c_0-Z=Z$ and
$-Z=Z-c_0$.  Hence
\[Z-Z=Z+(-Z)=Z+Z-c_0=2Z-c_0 .\]

Let $d\in(Z-Z)\cap\Z^2$ and put $x=d+c_0\in 2Z\cap\Z^2$.

The vertices of $Z$ are subset sums of the vectors $d_iv_i$ and so are lattice
points; thus $Z$ is a two-dimensional lattice polygon.  Take a \emph{unimodular
triangulation} of $Z$: every two-dimensional lattice polygon can be subdivided
into lattice triangles of area $\tfrac12$ whose vertices are lattice points of
$Z$.  Since $x/2\in Z$, it lies in some unimodular triangle $\triangle ABC$
(boundary included), and $(B-A,C-A)$ is a basis of $\Z^2$.  Write
\[\tfrac12x=A+\alpha(B-A)+\beta(C-A),\qquad \alpha,\beta\ge0,\ \alpha+\beta\le1 .\]
Then $x-2A=(2\alpha)(B-A)+(2\beta)(C-A)$.  The left-hand side is a lattice point
and $(B-A,C-A)$ is a lattice basis, so $2\alpha,2\beta\in\Z_{\ge0}$ with
$2\alpha+2\beta\le2$.  The cases are:
\begin{center}
\begin{tabular}{@{}lll@{}}
\toprule
$(2\alpha,2\beta)$ & $x$ & decomposition\\
\midrule
$(0,0)$ & $2A$                        & $A+A$\\
$(1,0)$ & $2A+(B-A)=A+B$              & $A+B$\\
$(0,1)$ & $A+C$                       & $A+C$\\
$(2,0)$ & $2A+2(B-A)=2B$              & $B+B$\\
$(0,2)$ & $2C$                        & $C+C$\\
$(1,1)$ & $2A+(B-A)+(C-A)=B+C$        & $B+C$\\
\bottomrule
\end{tabular}
\end{center}
In every case $x=p+q$ with $p,q\in\{A,B,C\}\subseteq Z\cap\Z^2$.

Finally $d=x-c_0=p-(c_0-q)$, and $c_0-q\in c_0-Z=Z$ is a lattice point.
\end{proof}

\begin{cor}\label{cor:pair}
The vector $d$ of Proposition~\ref{prop:witness} can be written $d=q_1-q_0$ with
$q_0,q_1\in Z\cap\Z^2$.  Putting $q=q_0$, both $q$ and $q+d$ lie in
$Z\cap\Z^2$.
\end{cor}

\section{Quadratic observables make up the whole linear
deficit}\label{sec:quadratic}

Fix $d$ and $q$ as in Proposition~\ref{prop:witness} and
Corollary~\ref{cor:pair}, and put
\[h(z)=\eta(z)\,\eta(z+d),\qquad J:=Dh\ne0\ \text{of finite support}.\]

\begin{lem}\label{lem:capacity}
For every $r\in R_Z(S)$ one has $r+q\in S$ and $r+q+d\in S$.
\end{lem}

\begin{proof}
We have $r+Z\subseteq\Conv(S)$ and $q,q+d\in Z\cap\Z^2$, so
\[r+q,\ r+q+d\ \in\ (r+Z)\cap\Z^2\ \subseteq\ \Conv(S)\cap\Z^2=S .\qedhere\]
\end{proof}

Consequently, for each $r\in R_Z(S)$ the \emph{quadratic scalar observable}
\[\Phi_r\colon\ \cL\to\Q,\qquad
  \Phi_r(\gamma)=w\bigl(\gamma(r+q)\bigr)\cdot w\bigl(\gamma(r+q+d)\bigr)\]
is well defined, and
$\Phi_r\bigl(\theta|_{u+S}\bigr)=\eta(u+r+q)\,\eta(u+r+q+d)=h(u+r+q)$.

\begin{lem}\label{lem:independence}
The family $\{\Phi_r\}_{r\in R_Z(S)}$ is linearly independent in
$\{\cL\to\Q\}/U_S$.
\end{lem}

\begin{proof}
Suppose not: there are $(c_r)$, not all zero, and $\psi\in U_S$ with
$\sum_rc_r\Phi_r=\psi$ on $\cL$.  Substituting $\gamma=\theta|_{u+S}$,
\[\sum_{r\in R_Z(S)}c_r\,h(u+r+q)=c_0+\sum_{s\in S}\mu_s\,\eta(u+s)
  \qquad\text{for all }u .\]
Apply $D$ in $u$ to both sides.  The right-hand side gives
$Dc_0+\sum_s\mu_s\,T^s(D\eta)=0$ by \eqref{eq:caseB}; the left-hand side gives
$\bigl(\sum_rc_rT^{r+q}\bigr)(Dh)=\bigl(\sum_rc_rT^{r+q}\bigr)J$.

So $\bigl(\sum_rc_rT^{r+q}\bigr)J=0$.  The exponents $r+q$ for $r\in R_Z(S)$ are
distinct and the $(c_r)$ are not all zero, so that Laurent polynomial is
non-zero; and $J$ is non-zero of finite support.  This contradicts
Lemma~\ref{lem:domain}.
\end{proof}

\begin{thm}[Theorem~T]\label{thm:main}
$P_\theta(S)\ \ge\ |S|+1$.
\end{thm}

\begin{proof}
The case $R_Z(S)=\varnothing$ is Remark~\ref{rem:empty-RZ}.  Assume
$R_Z(S)\ne\varnothing$ and put
\[V=U_S+\spn_{\Q}\{\Phi_r:r\in R_Z(S)\}\ \subseteq\ \{\cL\to\Q\}.\]
By Lemma~\ref{lem:independence}, $\dim V=\dim U_S+|R_Z(S)|$; by
Lemma~\ref{lem:budget},
\[\dim V\ \ge\ \bigl(|S|-|R_Z(S)|+1\bigr)+|R_Z(S)|=|S|+1 .\]
On the other hand $V$ is a space of functions on the finite set $\cL$, so
$\dim V\le|\cL|=P_\theta(S)$.
\end{proof}

\begin{rem}[The balance]\label{rem:balance}
The zonotope $Z$ enters the proof twice and with opposite signs: it is the
\emph{cost} of the affine relations, Lemma~\ref{lem:budget} subtracting
$|R_Z(S)|$, and it is the \emph{capacity} of the quadratic witness,
Lemma~\ref{lem:capacity} supplying exactly $|R_Z(S)|$ positions.  The two cancel
exactly, independently of the shape of $S$ and of the size of $Z$.  The
sampled modular ranks in Appendix~\ref{app:numerics} attain the lower bound of
Lemma~\ref{lem:budget} on every reported Case-B instance.
\end{rem}

\section{Reduction to the convex Nivat conjecture}\label{sec:reduction}

This section starts from a non-periodic configuration $\xi$ of low convex
complexity, extracts a star configuration from $X_\xi$, and contradicts
Theorem~\ref{thm:main}.  The chain is
\[
\begin{array}{c}
\text{Theorem~\ref{thm:KS}~\cite{KS}}\ \longrightarrow\
\text{Lemma~\ref{lem:modp} (mod $p$)}\ \longrightarrow\
\text{Theorems~\ref{thm:coneed} and \ref{thm:colregion}~\cite{Colle23a}}\\[4pt]
\longrightarrow\ \text{Proposition~\ref{prop:firsthalf} (first half-plane,
this paper)}\\[4pt]
\longrightarrow\ \text{Theorem~\ref{thm:secondhalf} (second half-plane,
this paper)}\ \longrightarrow\ \text{Lemma~\ref{lem:normalise}}\
\longrightarrow\ \text{Theorem~\ref{thm:main}} .
\end{array}
\]
Every link is either a result of \cite{KS} or \cite{Colle23a}, or proved here in
full.

\begin{thm}[Structure input]\label{thm:structure}
Let $\xi$ be a counterexample of minimal order in the sense of
Remark~\ref{rem:minimal}.  Then there are a prime $p$ with $\cA\subset\dbr{p}$,
a non-periodic $\vartheta\in X_\xi$, an $m\ge2$ and $\Fp$-valued
$G_1,\dots,G_m$ with $\vartheta=\sum_iG_i$ over $\Fp$, each $G_i$ having a
non-zero period $h_i$, the $h_i$ pairwise non-parallel and no component doubly
periodic, such that for every $i$:
\begin{enumerate}
\item[\textup{(i)}] $G_i$ is fully periodic on a half-plane $U_i$ whose
boundary is parallel to $h_i$;
\item[\textup{(ii)}] $G_i$ is fully periodic on a half-plane $V_i$ disjoint
from $U_i$.
\end{enumerate}
\end{thm}

Part~(i) is proved in \S\S\ref{ssec:external}--\ref{ssec:firsthalf}
(Proposition~\ref{prop:firsthalf} and Corollary~\ref{cor:firsthalf}); part~(ii)
is proved in \S\ref{ssec:secondhalf} (Theorem~\ref{thm:secondhalf}).

\subsection{Half-planes, regions and full periodicity}\label{ssec:halfplanes}

\begin{dfnhalf}
A \emph{half-plane} is a set of the form
$\cH=\{z\in\Z^2:\langle z,n\rangle>c\}$ or
$\cH=\{z\in\Z^2:\langle z,n\rangle\ge c\}$ with $n\in\R^2\setminus\{0\}$ and
$c\in\R$; we call $n$ its normal.  For an integer vector $g$: if
$\langle g,n\rangle\ge0$ then $\cH+g\subseteq\cH$; if $\langle g,n\rangle>0$
then $\bigcup_{N\ge0}(\cH-Ng)=\Z^2$.
\end{dfnhalf}

Call $R\subset\Z^2$ a \emph{lattice-convex region} if $R=C\cap\Z^2$ for some
closed convex $C\subseteq\R^2$.  Writing $C_R:=\overline{\Conv}(R)$ for the
closed convex hull, one still has $R=C_R\cap\Z^2$, since $R\subseteq
C_R\subseteq C$.  Let
\[K_R:=\{w\in\R^2:\ C_R+w\subseteq C_R\}\]
be the recession cone of $C_R$, a closed convex cone.  If $h\in\Z^2$ satisfies
$R+h\subseteq R$, then $C_R+h=\overline{\Conv}(R+h)\subseteq C_R$, that is
$h\in K_R$.  A half-plane with rational boundary direction is a
lattice-convex region: after rescaling its normal to an integer vector, a strict
inequality can be replaced by a weak inequality with an integer threshold.
For such a half-plane, $K_R=\{w:\langle w,n\rangle\ge0\}$.

For a non-empty $\cU\subseteq\Z^2$ say that $G\colon\Z^2\to\Fp$ is \emph{fully
periodic on $\cU$} (\cite[Definition~2.1]{CG}) if
there are $\R$-linearly independent $h,h'\in\Z^2$ with $\cU+h\subseteq\cU$,
$\cU+h'\subseteq\cU$, $G(z+h)=G(z)$ and $G(z+h')=G(z)$ for all $z\in\cU$.  The
restriction of a doubly periodic field to any half-plane is fully periodic: take
two independent vectors of its period lattice pointing into the half-plane.
Conversely:

\begin{lem}[Global extension]\label{lem:extend}
Let $R$ be a lattice-convex region or a half-plane, and let $G$ be fully periodic
on $R$, with vectors $h,h'$.  Put $g:=h+h'$.  Then for every $z\in\Z^2$ one
has $z+Ng\in R$ for all large $N$, and there is a doubly
periodic field $\widetilde G$ with $\Per(\widetilde G)\supseteq\Z h+\Z h'$ and
$\widetilde G=G$ on $R$.  If $R$ is a lattice-convex region, then
$\dim K_R=2$ and $g\in\operatorname{int}K_R$; if $R$ is a half-plane with
normal $n$, then $\langle g,n\rangle>0$.
\end{lem}

\begin{proof}
First suppose $R$ is a lattice-convex region.  From $R+h\subseteq R$ and
$R+h'\subseteq R$ we get $h,h'\in K_R$; they are
linearly independent, so the closed sector they span lies in $K_R$, its interior
is non-empty and contains $g$, whence $g\in\operatorname{int}K_R$ and
$\dim K_R=2$.  Pick $x_0\in C_R$; by the definition of the recession cone
$x_0+K_R\subseteq C_R$.  For any $z$ we have
$z+Ng-x_0=N\bigl(g+\tfrac{z-x_0}{N}\bigr)$, and for $N$ large the vector in
brackets lies in a neighbourhood of $g$ contained in $K_R$; hence
$z+Ng\in C_R\cap\Z^2=R$.

For a half-plane $R$ with normal $n$, there is a direct argument, whether its
boundary direction is rational or irrational and whether its defining
inequality is strict or weak.  The inclusion $R+h\subseteq R$ implies
$\langle h,n\rangle\ge0$: otherwise, starting at any point of $R$, repeated
addition of $h$ would eventually leave $R$.  Likewise
$\langle h',n\rangle\ge0$.  Since $h,h'$ are independent, these two inner
products cannot both vanish, so $\langle g,n\rangle>0$.  Thus $z+Ng\in R$ for
all sufficiently large $N$.  In either case the extension is constructed as
follows.

\emph{Step 1.}  For $z\in R$ and $k\ge0$: $z+kh\in R$ and $G(z+kh)=G(z)$, by
induction on $k$; likewise for $h'$ and for $g$.

\emph{Definition.}  For $z\in\Z^2$ pick $N\ge0$ with $z+Ng\in R$ and set
$\widetilde G(z)=G(z+Ng)$.  This is well defined: if $N\le N'$ both work, then
$G(z+N'g)=G\bigl((z+Ng)+(N'-N)g\bigr)=G(z+Ng)$ by Step 1.  Taking $N=0$ shows
$\widetilde G=G$ on $R$.

\emph{Periodicity.}  Given $z$, pick $N$ with $z+Ng\in R$; then also
$z+h+Ng\in R$, so
$\widetilde G(z+h)=G\bigl((z+Ng)+h\bigr)=G(z+Ng)=\widetilde G(z)$ by Step 1.
The same for $h'$.
\end{proof}

An open half-plane with irrational boundary direction need not be a
lattice-convex region under the closed-convex-set definition above.  The direct
half-plane argument in Lemma~\ref{lem:extend} covers this case as well.

\subsection{The results quoted, and counterexamples of minimal
order}\label{ssec:external}

\begin{rem}[Counterexamples of minimal order]\label{rem:minimal}
For a non-periodic configuration $\xi$ of low convex complexity with
$\cA\subset\Z_+$, write $\ord(\xi)$ for the length of a $\Z$-minimal periodic
decomposition of $\xi$; by \cite{KS} a configuration of low complexity has a
non-trivial annihilator and hence a finite periodic decomposition, so $\ord$ is
finite.  If the convex Nivat conjecture has a counterexample, then one may take
a counterexample $\xi$ with $\ord(\xi)$ minimal, and for such a $\xi$ every
non-periodic $\xi'\in X_\xi$ is again a counterexample, with
$\ord(\xi')=\ord(\xi)$.

\emph{Proof.}  Every finite-window pattern of $\xi'$ is a pattern of $\xi$, so
$\xi'$ has low convex complexity; being non-periodic it is a counterexample, and
minimality gives $\ord(\xi')\ge\ord(\xi)=:m$.  Conversely write
$\xi=\sum_{i\le m}\xi_i$ minimal, with $\xi_i$ of period $h_i$; the directions
are pairwise distinct, since two components with parallel periods could be
merged, contradicting minimality.  Then $\varphi=\prod_i(X^{h_i}-1)$ annihilates
$\xi$; an annihilation relation is a linear identity on a finite window and
therefore passes to the orbit closure, so $\varphi$ annihilates $\xi'$, and by
\cite{KS} the configuration $\xi'$ has a periodic decomposition with at most $m$
terms, that is $\ord(\xi')\le m$.\hfill$\square$
\end{rem}

Consequently every result of \cite{Colle23a} proved under the standing
assumption that the non-periodic configurations of the orbit closure all have
the same order (\cite[\S4.1]{Colle23a}) — including
\cite[Theorem~1.10]{Colle23a} and the Claims of its
\S4 — is available for a counterexample $\xi$ of minimal order; it is under that
assumption that Theorem~\ref{thm:colregion} below quotes
\cite[Lemma~4.6]{Colle23a}.

\begin{thm}[Kari--Szabados \cite{KS}]\label{thm:KS}
Let $\xi\colon\Z^2\to\cA$ with $\cA\subset\Z$ finite.
\begin{enumerate}
\item[\textup{(a)}] If $P_\xi(S)\le|S|$ for some non-empty finite
$S\subset\Z^2$, then $\operatorname{Ann}_\Z(\xi)\ne0$.
\item[\textup{(b)}] If $\operatorname{Ann}_\Z(\xi)\ne0$, then there are pairwise
non-parallel $h_1,\dots,h_m\in\Z^2\setminus\{0\}$ such that
$\varphi=\prod_{i=1}^m(X^{h_i}-1)$ annihilates $\xi$; and every
$\eta\in\Z^{\Z^2}$ annihilated by such a $\varphi$ can be written
$\eta=\sum_i\eta_i$ with $\eta_i\in\Z^{\Z^2}$ of period $h_i$ (the values
possibly unbounded).
\end{enumerate}
\end{thm}

Both (a) and (b) are in \cite{KS}; the argument for (a) is the same as in
the proof of Proposition~\ref{prop:case-A}.

\begin{lem}[Passing to the orbit closure and reducing mod
$p$]\label{lem:modp}
Let $\xi$ be as in Theorem~\ref{thm:KS} with $\cA\subset\Z_+$, and let
$\varphi=\prod_{i=1}^m(X^{h_i}-1)$ annihilate $\xi$, with the $h_i$ non-zero
and pairwise non-parallel as in Theorem~\ref{thm:KS}(b).  Then for every
$\vartheta\in X_\xi$ and every prime $p>\max\cA$:
\begin{enumerate}
\item[\textup{(a)}] $\varphi$ annihilates $\vartheta$, so
$\vartheta=\sum_i\vartheta_i$ with $\vartheta_i\in\Z^{\Z^2}$ of period $h_i$;
\item[\textup{(b)}] putting $\bar\vartheta_i:=\vartheta_i\bmod p\colon\Z^2\to
\Fp$ one has $\sum_i\bar\vartheta_i=\vartheta$, the values of $\vartheta$ being
read in $\Fp$ (no colours are merged, since $\cA\subset\dbr{p}$), and each
$\bar\vartheta_i$ has period $h_i$;
\item[\textup{(c)}] after merging the components whose periods are parallel one
may assume the $h_i$ pairwise non-parallel; and if $\vartheta$ is non-periodic,
at least two components remain.
\end{enumerate}
\end{lem}

\begin{proof}
(a) At each $z$ the identity $\varphi\vartheta=0$ involves only
$\vartheta|_{z+\supp\varphi}$, so it is a local function of $\vartheta$; it
vanishes on every translate of $\xi$, and local functions are continuous for the
product topology, hence it vanishes on $X_\xi$.  Now apply
Theorem~\ref{thm:KS}(b).  (b) Reduce pointwise.  (c) The sum of two components
with parallel periods is again periodic (take a common multiple of the two
periods); if only one term were left, $\vartheta$ would be periodic.
\end{proof}

\begin{thm}[Colle {\cite[Theorem~1.9]{Colle23a}}]\label{thm:coneed}
Let $\cA\subset\Z$ be finite and let $\xi\in\cA^{\Z^2}$ have a non-trivial
annihilator.  If for every
line $\ell$ either $-\boldsymbol\ell\notin\mathrm{ONED}(\xi)$ or
$\boldsymbol\ell\notin\mathrm{ONED}(\xi)$, then $\xi$ is fully periodic.  In
particular, if $\xi$ is non-periodic there is a line $\ell$ with
$\pm\boldsymbol\ell\in\mathrm{ONED}(\xi)$.
\end{thm}

The proof of \cite[Theorem~1.9]{Colle23a} uses the periodic-limit theorem of
Kari and Moutot \cite{KM20,KM23} and the Boyle--Lind theorem \cite{BoyleLind};
it does not depend on the equal-order assumption of \cite[\S4.1]{Colle23a}.

\begin{thm}[Colle {\cite[proof of Lemma~4.6]{Colle23a}}]\label{thm:colregion}
Let $\xi$ be a counterexample of minimal order in the sense of
Remark~\ref{rem:minimal}, so that the equal-order assumption of \S4.1 of
\cite{Colle23a} is in force, and let $\ell$ be as in
Theorem~\ref{thm:coneed}.  Then there are a non-periodic $\vartheta\in X_\xi$
and a lattice-convex region $R$ such that $\vartheta|_R$ is fully periodic.
Here $R$ is a $(\boldsymbol\ell,\boldsymbol\ell')$-region: the set of lattice
points of a convex unbounded polygonal region whose boundary consists of two
unbounded edges, along $\ell$ and along $\ell'$, both nonexpansive lines for
$X_\xi$, and finitely many bounded edges.
\end{thm}

Theorem~\ref{thm:colregion} is extracted from Cases 1 and 2 in the proof of
\cite[Lemma~4.6]{Colle23a}.  There, regional periods require equality only
where a point and its translate both belong to the region.  Both cases also
provide a period parallel to each of the two non-parallel unbounded edges.
Orient these periods $h,h'$ along the corresponding outward rays of $C_R$.
Closed convexity puts $h,h'$ in $K_R$: if $x_0+\R_{\ge0}h\subseteq C_R$,
then for $x\in C_R$ and $t\ge0$, letting $s\to\infty$ in
$(1-t/s)x+(t/s)(x_0+sh)\in C_R$ gives $x+th\in C_R$.
Thus $R+h,R+h'\subseteq R$, and the overlap identities hold at every point
of $R$, giving full periodicity in the sense of \S\ref{ssec:halfplanes}.
The source proof uses Lemma~3.5, Lemma~4.2, Propositions~2.13--2.14 and
Claims~4.7--4.16 of \cite{Colle23a}, of which Claims~4.11 and 4.16 depend on the
equal-order assumption.  Only three items of it are used below: $\vartheta$ is
non-periodic, $R$ is a lattice-convex region, and $\vartheta$ is fully periodic
on $R$.  The position of the two asymptotic directions of the region is not
needed: Proposition~\ref{prop:firsthalf}(b) gives the half-plane containments
for its recession cone that are used in Lemma~\ref{lem:available}.

In the Lean formalisation Theorems~\ref{thm:KS}, \ref{thm:coneed} and
\ref{thm:colregion} are not assumed: the forms in which they are used below are
proved, together with the parts of \cite{Colle23a} their proofs rest on.  The one step of the source that had to
be replaced, in the proof of Lemma~3.5(i), is described in
Appendix~\ref{app:lean}.

\subsection{The first half-plane}\label{ssec:firsthalf}

This subsection proves Theorem~\ref{thm:structure}(i).  The only tools are
Lemma~\ref{lem:extend}, the difference operators $Q_i$ of \eqref{eq:Qexp}, and
the finite-state argument of step~(6) of Lemma~\ref{lem:propagation} in
Appendix~\ref{app:szabados}.

\begin{lem}[One-sided finite-state lemma]\label{lem:finitestate}
Let $G\colon\Z^2\to\Fp$ have period $kv$ with $v$ primitive and $k\ge1$, let
$(v,u)$ be a basis of $\Z^2$, so $\det(v,u)=\pm1$, and define the row coordinate
of $z=sv+tu$ by $t=t(z)=\det(v,u)\det(v,z)$.  Let
$H_1,\dots,H_r\in\Z^2$ with $r\ge1$ satisfy $\beta_j:=t(H_j)\ne0$, and put
$\beta_-:=\sum_{\beta_j<0}\beta_j$ and $\beta_+:=\sum_{\beta_j>0}\beta_j$.  If
\[\Delta:=\prod_{j=1}^r\bigl(T^{H_j}-1\bigr)\,G\]
vanishes on the half-plane $\{t\ge c\}$, then there is $q\ge1$ with
$G(z+qu)=G(z)$ for all $z$ with $t(z)\ge c+\beta_-$; in particular $G$ is fully
periodic on $\{t\ge c+\beta_-\}$, with the periods $kv$ and $qu$, neither of
which leaves that half-plane.  If $\Delta\equiv0$, then $G$ is doubly periodic.
\end{lem}

\begin{proof}
Write $z=sv+tu$.  Let $\cW:=\{f\colon\Z\to\Fp:\ f(s+k)=f(s)\}$, a set of
cardinality $p^k$, and for each row $t$ let $G_t\in\cW$ be
$G_t(s):=G(sv+tu)$.  Write $H_j=\alpha_jv+\beta_ju$.  Then
\[(T^{H_j}G)(sv+tu)=G\bigl((s+\alpha_j)v+(t+\beta_j)u\bigr)
  =(\rho^{\alpha_j}G_{t+\beta_j})(s),\]
where $(\rho f)(s):=f(s+1)$ is an invertible linear map of $\cW$.  Since $\rho$
commutes with the translations in the $t$ direction, for every $t$
\[\Delta_t=\sum_{C\subseteq[r]}(-1)^{r-|C|}\rho^{\alpha_C}G_{t+\beta_C},\qquad
  \alpha_C:=\sum_{j\in C}\alpha_j,\quad\beta_C:=\sum_{j\in C}\beta_j .\]
The maximum $\beta_+=\max_C\beta_C$ is attained only at
$C_+:=\{j:\beta_j>0\}$ and the minimum $\beta_-=\min_C\beta_C$ only at
$C_-:=\{j:\beta_j<0\}$; put $L:=\beta_+-\beta_-\ge1$.  From $\Delta_t=0$ for
$t\ge c$ we get:

\smallskip
\noindent\textup{(F)}\quad the forward recursion
\[G_{t+\beta_+}=-(-1)^{r-|C_+|}\rho^{-\alpha_{C_+}}\sum_{C\ne C_+}(-1)^{r-|C|}
 \rho^{\alpha_C}G_{t+\beta_C},\]
whose right-hand side involves only $G_{t+\beta_-},\dots,G_{t+\beta_+-1}$;

\noindent\textup{(B)}\quad likewise $G_{t+\beta_-}$ is determined by
$G_{t+\beta_-+1},\dots,G_{t+\beta_+}$.

\smallskip
Let the state be $\mathbf x_t:=(G_{t+\beta_-},\dots,G_{t+\beta_+-1})\in\cW^L$
and let $\Sigma:=\{\mathbf x_t:t\ge c\}$, a finite set.  By (F), $\mathbf
x_{t+1}$ is determined by $\mathbf x_t$, so $\Psi(\mathbf x_t):=\mathbf x_{t+1}$
is a well-defined map $\Sigma\to\Sigma$; by (B) it is injective on $\Sigma$.  An
injective self-map of a finite set is a bijection, so
$(\mathbf x_t)_{t\ge c}$ is purely periodic: there is $q\ge1$ with
$\mathbf x_{t+q}=\mathbf x_t$ for all $t\ge c$.  Reading off the first
component, $G_{t+q}=G_t$ for $t\ge c+\beta_-$, that is $G(z+qu)=G(z)$ whenever
$t(z)\ge c+\beta_-$.  If $\Delta\equiv0$, take instead the finite set
$\Sigma_{\mathrm{all}}:=\{\mathbf x_t:t\in\Z\}$.  Both recursions hold at every
integer $t$, so the same argument makes the forward map a permutation of
$\Sigma_{\mathrm{all}}$.  A positive power $q$ of this permutation is the
identity, giving $G_{t+q}=G_t$ for every $t\in\Z$.  Thus $G$ has the two
independent periods $qu$ and $kv$.
\end{proof}

\begin{prop}[The first half-plane]\label{prop:firsthalf}
Let $\vartheta\colon\Z^2\to\Fp$ with $\vartheta=\sum_{i=1}^m\vartheta_i$,
$m\ge2$, where $\vartheta_i$ has period $h_i=k_iv_i$ with $v_i$ primitive and
the $v_i$ pairwise non-parallel.  Let $R$ be a lattice-convex region on which
$\vartheta$ is fully periodic, with vectors $a,b$, and put $K:=K_R$.  Then
\begin{enumerate}
\item[\textup{(a)}] $a,b\in K$ and $\dim K=2$;
\item[\textup{(b)}] for each $i$ exactly one of the following holds:
 \begin{itemize}
 \item[$(\alpha)$] the line $\R v_i$ meets $\operatorname{int}K$, and then
 $\vartheta_i$ is doubly periodic;
 \item[$(\beta)$] otherwise there are $\sigma_i\in\{+,-\}$ and $t_i\in\Z$
 such that $K\subseteq\{\sigma_i\pi_i\ge0\}$ and $\vartheta_i$ is fully
 periodic on the half-plane $U_i:=\{\sigma_i\pi_i\ge t_i\}$;
 \end{itemize}
\item[\textup{(c)}] if $\vartheta$ is non-periodic and no component is doubly
periodic, then $K$ is a closed sector of opening strictly between $0$ and
$\pi$, and its interior contains no $\pm v_j$.  In particular $K$ cannot be
a half-plane.
\end{enumerate}
\end{prop}

\begin{proof}
\emph{Step 1 (global extension).}  By Lemma~\ref{lem:extend}, $a,b\in K$,
$\dim K=2$, $g:=a+b\in\operatorname{int}K$, and there is a doubly periodic
$\widetilde\vartheta$ with $\widetilde\vartheta=\vartheta$ on $R$ and
$\Per(\widetilde\vartheta)$ of finite index.  Choose $N_0\ge1$ with
$N_0\Z^2\subseteq\Per(\widetilde\vartheta)$ and put $H_j:=N_0h_j$.  Pick
$x_0\in C_R$, so that $x_0+K\subseteq C_R$.

\emph{Step 2 (differencing).}  Fix $i$ and put $Q_i:=\prod_{j\ne i}(T^{H_j}-1)$,
with formal exponent set $\cE_i=\{H_C:C\subseteq[m]\setminus\{i\}\}$ and the
expansion \eqref{eq:Qexp}.  If $z+\cE_i\subseteq R$, then
\[(Q_i\vartheta)(z)=\sum_{C\subseteq[m]\setminus\{i\}}(-1)^{m-1-|C|}\,
  \widetilde\vartheta\bigl(z+H_C\bigr)
  =\widetilde\vartheta(z)\cdot(1-1)^{m-1}=0,\]
because every $H_C$ lies in $\Per(\widetilde\vartheta)$.  On the other hand
$H_j$ is a period of $\vartheta_j$, so $Q_i\vartheta_j=0$ for $j\ne i$, and
therefore
\[D_i:=Q_i\vartheta_i=Q_i\vartheta\quad\text{vanishes on }
  R_i':=\{z:\ z+\cE_i\subseteq R\}.\]
Put $C':=\bigcap_{e\in\cE_i}(C_R-e)$.  Then $R_i'=C'\cap\Z^2$, and
$C'\supseteq y_0+K$ where $y_0:=x_0+N'g$ with $N'$ large: indeed for
$x\in y_0+K$ and $e\in\cE_i$ one has
$x-x_0+e=N'g+e+(x-y_0)\in K$, because $N'g+e\in K$.

\emph{Step 3 (saturation along $v_i$).}  $D_i$ has period $h_i$, so $D_i$
vanishes on $\mathrm{Sat}_i:=\{z:(z+\Z h_i)\cap R_i'\ne\varnothing\}$.  There
are two cases.

$(\alpha)$ $\R v_i\cap\operatorname{int}K\ne\varnothing$.  Then $v_i$ or $-v_i$
lies in $\operatorname{int}K$, so for any $z$ the point $z+nh_i$ enters
$y_0+K$ as $n\to+\infty$ or as $n\to-\infty$; hence
$\mathrm{Sat}_i=\Z^2$ and $D_i\equiv0$.  Apply Lemma~\ref{lem:finitestate} in
its whole-plane case to the basis $(v_i,u_i)$, the hypothesis
$t(H_j)=N_0\det(v_i,h_j)\ne0$ holding because the directions are pairwise
non-parallel: $\vartheta_i$ is doubly periodic.

$(\beta)$ $\R v_i\cap\operatorname{int}K=\varnothing$.  A convex set
and a line missing its interior can be separated, so $K$ lies in one of the two
closed half-planes bounded by $\R v_i$; choose $\sigma_i$ with
$K\subseteq\{\sigma_i\pi_i\ge0\}$.  For the vector
$g\in\operatorname{int}K$ of Step 1 put $a_i:=\sigma_i\pi_i(g)>0$.
Choose $\varepsilon_i>0$ so small that $g\pm\varepsilon_i v_i\in K$.
For every $\tau>0$, the two points
$\frac{\tau}{a_i}(g\pm\varepsilon_i v_i)$ lie in
$K\cap\{\sigma_i\pi_i=\tau\}$, so this slice contains a segment parallel to
$v_i$ of length $2\varepsilon_i\tau|v_i|/a_i$.
Choose $c_i>0$ with $2\varepsilon_i c_i|v_i|/a_i\ge|h_i|$.
If $\sigma_i\pi_i(z-y_0)\ge c_i$, then
\[(z+\R v_i)\cap(y_0+K)=y_0+\bigl(K\cap\{\sigma_i\pi_i
  =\sigma_i\pi_i(z-y_0)\}\bigr)\]
contains a segment of length $\ge|h_i|$ and hence a point of $z+\Z h_i$, these
points being spaced $|h_i|$ apart along that line; so $z\in\mathrm{Sat}_i$.
Therefore $D_i$ vanishes on the half-plane $\{\sigma_i\pi_i\ge t\}$ with
$t:=\lceil c_i+\sigma_i\pi_i(y_0)\rceil$.  Apply Lemma~\ref{lem:finitestate} to the basis
$(v_i,\sigma_iu_i)$, for which $t(z)=\sigma_i\pi_i(z)$ and
$t(H_j)=\sigma_iN_0\det(v_i,h_j)\ne0$: the field $\vartheta_i$ is fully periodic
on $U_i:=\{\sigma_i\pi_i\ge t_i\}$ with the integer $t_i:=t+\beta_-$.
The choice of $\sigma_i$ also gives $K\subseteq\{\sigma_i\pi_i\ge0\}$,
as required in (b).

\emph{Step 4 (part (c)).}  Suppose $\vartheta$ is non-periodic and no component
is doubly periodic.  Then case~$(\alpha)$ occurs for no $i$, that is, no line
$\R v_i$ meets $\operatorname{int}K$.  Now $K$ is a closed convex cone with
$\dim K=2$.  If $K$ were a half-plane, at most one line $\R v_i$ could lie on
its boundary and all the others would fall into case~$(\alpha)$, a
contradiction; and $K=\R^2$ is excluded in the same way.  A closed convex cone
of dimension two in the plane is either the whole plane, a half-plane, or a
closed sector of opening strictly between $0$ and $\pi$.  Thus $K$ is such a
sector, and its interior contains no $\pm v_j$.
\end{proof}

\begin{cor}[Theorem~\ref{thm:structure}(i)]\label{cor:firsthalf}
Let $\vartheta$ and $R$ be as in Theorem~\ref{thm:colregion}, let $p>\max\cA$ be
a prime, and let $\vartheta=\sum_i\bar\vartheta_i$ be the decomposition of
Lemma~\ref{lem:modp}, with the $h_i$ pairwise non-parallel.  Let $I$ be the
indices of the components that are not doubly periodic and let $B$ be the sum
of all the other components, so $B$ is doubly periodic.  Since $\vartheta$ is
non-periodic, $|I|\ge2$: if $I$ were empty then $\vartheta=B$ would be doubly
periodic, and if $I$ had one element then a common period of that component
and $B$ along its direction would be a period of $\vartheta$.
For each $i\in I$, case~$(\alpha)$ of Proposition~\ref{prop:firsthalf} is
impossible.  Part (b) therefore gives a half-plane
$U_i=\{\sigma_i\pi_i\ge t_i\}$ of full periodicity and
$K_R\subseteq\{\sigma_i\pi_i\ge0\}$.
Choose $i_0\in I$, set $G_{i_0}:=\bar\vartheta_{i_0}+B$, and set
$G_i:=\bar\vartheta_i$ for $i\in I\setminus\{i_0\}$.
The field $G_{i_0}$ has a non-zero period along $h_{i_0}$, obtained by taking
a suitable multiple in $\Per(B)$, and is fully periodic on $U_{i_0}$:
there it agrees with the sum of $B$ and the doubly periodic extension of
$\bar\vartheta_{i_0}|_{U_{i_0}}$ from Lemma~\ref{lem:extend}.
It is not doubly periodic, since subtracting $B$ would then make
$\bar\vartheta_{i_0}$ doubly periodic.
Thus $\vartheta=\sum_{i\in I}G_i$, no component is doubly periodic, and their
periods are pairwise non-parallel.  Retaining the primitive directions and
renaming the period multiples if necessary, every $G_i$ is fully periodic on
$U_i$, with $K_R\subseteq\{\sigma_i\pi_i\ge0\}$.  The cone $K_R$ has
dimension two and is a closed sector of opening $<\pi$ whose interior
contains no remaining component direction, by
Proposition~\ref{prop:firsthalf}(a),(c) applied to this decomposition.
\end{cor}

\begin{rem}\label{rem:relation}
\textup{(1)} Proposition~\ref{prop:firsthalf} is the one-sided version of
Lemma~\ref{lem:fi} of \S\ref{sec:twopoint} — the support of $f_i=Q_i\eta$
lies in a strip — and Lemma~\ref{lem:finitestate} is the general form of
step~(6) of Lemma~\ref{lem:propagation}.
\textup{(2)} In Proposition~\ref{prop:firsthalf}(b) neither the non-periodicity
of $\vartheta$ nor the minimality of the decomposition is required.  The
corollary absorbs all doubly periodic components, including any that fall
into case~$(\beta)$.
\end{rem}

\subsection{The second half-plane}\label{ssec:secondhalf}

This subsection proves Theorem~\ref{thm:structure}(ii), in three steps.
Proposition~\ref{prop:twolimits} uses Theorem~\ref{thm:szabados} to show that,
when $\vartheta$ is pushed to the limit along a suitable direction, every
subsequential limit on the far side of one component is doubly periodic.
Lemma~\ref{lem:propagate2} shows that this forces the component to agree with a
doubly periodic field on a half-plane on that side, the point of the argument
being that there are only finitely many such limits.
Lemma~\ref{lem:available} guarantees that the directions required can always be
found, in a fixed order, so that an induction produces all of the $V_i$.

\begin{thm}[The second half-plane]\label{thm:secondhalf}
In the notation of Corollary~\ref{cor:firsthalf}, for every $i\in I$ there is a
half-plane $V_i$ with $V_i\cap U_i=\varnothing$ such that $G_i|_{V_i}$ is fully
periodic.
\end{thm}

The proof is at the end of the subsection.  By step~$(\gamma)$ of
Lemma~\ref{lem:normalise}, every half-plane of full periodicity of an
$h_i$-periodic field that is not doubly periodic is bounded by a line parallel
to $h_i$; so such a $V_i$ necessarily lies on the side of $U_i$ opposite to
$U_i$, and shrinking either $U_i$ or $V_i$ destroys nothing.

\begin{lem}[Bi-periodicity over $\Fp$]\label{lem:biper}
Let $G\colon\Z^2\to\Fp$ have a non-zero period $h$, let $d\in\Z^2$ be
non-parallel to $h$, and suppose $g:=(T^d-1)G$ is doubly periodic.  Then $G$ is
doubly periodic.
\end{lem}

\begin{proof}
Choose $M\ge1$ with $Md\in\Per(g)$.  From
$T^{Md}-1=\bigl(\sum_{i=0}^{M-1}T^{id}\bigr)(T^d-1)$ we get
$g':=(T^{Md}-1)G=\sum_{i=0}^{M-1}T^{id}g$, and
$T^{Md}g'=\sum_iT^{id}\,T^{Md}g=g'$.  Hence
\[(T^{pMd}-1)G=\Bigl(\sum_{i=0}^{p-1}T^{iMd}\Bigr)g'=p\,g'=0,\]
so $pMd$ is a period of $G$, and it is independent of $h$.
\end{proof}

\smallskip\noindent\textbf{Safe directions.}
In the notation of Corollary~\ref{cor:firsthalf}, whenever $\sigma_j=-$ replace
$(v_j,u_j,h_j)$ by $(-v_j,-u_j,-h_j)$.  This preserves $h_j=k_jv_j$ with
$k_j\ge1$ and $\det(v_j,u_j)=1$, and gives $\sigma_j=+$ throughout, that is
$U_j=\{\pi_j\ge t_j\}$.  For
$j\in I$ let $\cK_j$ be the set of half-planes on which $G_j$ is so far known to
be fully periodic: initially $\cK_j=\{U_j\}$, and a half-plane $V_j$ disjoint
from $U_j$ is added to $\cK_j$ as soon as it has been produced.  For
$\cH\in\cK_j$ write $\widetilde G_j^{\cH}$ for the global extension of
$G_j|_{\cH}$ furnished by Lemma~\ref{lem:extend}.  Fix $N\ge1$ divisible by all
the $k_j$ and such that $N\Z^2$ is contained in the period lattice of every
$\widetilde G_j^{\cH}$.  For a non-zero $d\in N\Z^2$ say that $d$ \emph{fixes}
$j$ if $d\parallel h_j$ — then $d\in\Z k_jv_j$ and $T^dG_j=G_j$ — and that $d$
is \emph{safe for $j$} if there is
$\cH=\{\tau\pi_j\ge t\}\in\cK_j$, $\tau\in\{+,-\}$, with $\tau\pi_j(d)>0$.  If
$d$ is safe for $j$ then $T^{nd}G_j\to\widetilde G_j^{\cH}$ pointwise as
$n\to\infty$, the convergence being eventual equality at each point: for a fixed
$z$ and large $n$ one has $z+nd\in\cH$, whence
$G_j(z+nd)=\widetilde G_j^{\cH}(z+nd)=\widetilde G_j^{\cH}(z)$.

\begin{prop}[Two-component limits]\label{prop:twolimits}
Let $b,c\in I$ with $b\ne c$, let $d\in N\Z^2$ fix $c$, let $\pi_b(d)<0$, and
let $d$ be safe for every $j\in I\setminus\{b,c\}$.  Let
$n_1<n_2<\cdots$ be such that $T^{n_kd}G_b$ converges pointwise — such a
sequence exists by compactness of $\Fp^{\Z^2}$ — and write $G_b^\infty$ for the
limit.  Then $G_b^\infty$ is doubly periodic.
\end{prop}

\begin{proof}
By the preceding paragraph $T^{n_kd}\vartheta\to\Xi:=G_c+B+G_b^\infty$
pointwise, where $B:=\sum_{j\ne b,c}\widetilde G_j^{\cH_j}$ is doubly periodic.
Since $X_\xi$ is closed, $\Xi\in X_\xi$, so $P_\Xi(S)\le P_\vartheta(S)\le|S|$.
Choose $h_c'\in(\Per(B)\cap\Z h_c)\setminus\{0\}$; then
$\theta_1:=G_c+B$ has period $h_c'$ and $\theta_2:=G_b^\infty$ has period $h_b$
(a pointwise limit of translates keeps the period), and $h_c'\nparallel h_b$.
Theorem~\ref{thm:szabados} supplies a non-zero period $d_0$ of
$\Xi=\theta_1+\theta_2$.  There are three cases.

\emph{(i) $d_0$ is parallel to neither $h_b$ nor $h_c$.}  Then
$g:=(T^{d_0}-1)\theta_1=-(T^{d_0}-1)\theta_2$ has both $h_c'$ and $h_b$ as
periods, hence is doubly periodic; Lemma~\ref{lem:biper}, applied to $\theta_1$
with $d_0\nparallel h_c'$, makes $\theta_1$ doubly periodic, so
$G_c=\theta_1-B$ is doubly periodic, contradicting
Corollary~\ref{cor:firsthalf}.

\emph{(ii) $d_0\parallel h_b$.}  Pick $d_1\in(\Z d_0\cap\Z h_b)\setminus\{0\}$.
From $T^{d_1}\Xi=\Xi$ and $T^{d_1}\theta_2=\theta_2$ we get
$T^{d_1}\theta_1=\theta_1$; as $d_1\nparallel h_c'$, $\theta_1$ is doubly
periodic and we conclude as in (i).

\emph{(iii) $d_0\parallel h_c$.}  Pick
$d_1\in(\Z d_0\cap\Z h_c')\setminus\{0\}$.  From $T^{d_1}\theta_1=\theta_1$ and
$T^{d_1}\Xi=\Xi$ we get $T^{d_1}\theta_2=\theta_2$; as $d_1\nparallel h_b$, the
field $\theta_2=G_b^\infty$ is doubly periodic.
\end{proof}

\begin{lem}[Available directions]\label{lem:available}
Let $O\subseteq I$ be the set of components for which no $V_j$ has yet been
produced, $O\ne\varnothing$, and let $N$ be chosen as above.  Then there are
$b\in O$, $c\in I\setminus\{b\}$ and $d\in N\Z^2$ satisfying all the hypotheses
of Proposition~\ref{prop:twolimits}.
\end{lem}

\begin{proof}
If $O=\{b\}$, take any $c\ne b$ and $d:=\tau Nh_c$ with $\tau\in\{+,-\}$ chosen
so that $\pi_b(\tau h_c)<0$; this is possible because $h_c\nparallel h_b$, so
$\pi_b(h_c)\ne0$.  For $j\ne b,c$ we have $j\notin O$, so
$\cK_j=\{U_j,V_j\}$ contains a half-plane on each side of $\pi_j$, and
$\pi_j(d)\ne0$; hence $d$ is safe for $j$.

If $|O|\ge2$, put $K_O:=\bigcap_{j\in O}\{\pi_j\ge0\}$.  By
Corollary~\ref{cor:firsthalf} we have $K_R\subseteq\{\pi_j\ge0\}$ for every
$j\in O$, so $K_O\supseteq K_R$ and $\dim K_O=2$.  Being the
intersection of $|O|\ge2$ closed half-planes through the origin with pairwise
distinct boundary lines $\R v_j$, the cone $K_O$ is a closed sector of opening
$<\pi$.  At each non-zero boundary point at least one defining inequality
is an equality; otherwise all are strict and the point is interior.
Consequently its two boundary rays lie on two of the lines $\R v_a$, $\R v_b$
with $a,b\in O$, $a\ne b$.  Also $\operatorname{int}K_O\subseteq\{\pi_j>0\}$ for
$j\in O$ contains no $\pm v_j$ with $j\in O$.  Write $r\subset\R v_a$ and
$r'\subset\R v_b$ for the two boundary rays, oriented so that turning
anticlockwise from $r$ to $r'$ sweeps exactly $K_O$.  Among the $2m$ rays
$\{\pm v_j:j\in I\}$, let $r_c=\varepsilon_cv_c$ ($c\in I$,
$\varepsilon_c\in\{+,-\}$) be the first one met when turning anticlockwise from
$r'$.  It exists, and $c\ne b$ with $\angle(r',r_c)<\pi$: the angle from $r'$
anticlockwise to $-r$ equals $\pi-\angle(r,r')<\pi=\angle(r',-r')$, so $-r$
occurs before $-r'$.  (The index $c$ may lie in $O$, and may equal $a$.)  Put
$d:=\varepsilon_cNh_c$.

We verify the hypotheses.  \textup{(1)} $d$ fixes $c$.  \textup{(2)} The
interior of the sector $Q:=\operatorname{cone}(r',r_c)$ contains none of the
rays $\pm v_j$, since $r_c$ is the first after $r'$, and its boundary rays are
$r'\subset\R v_b$ and $r_c\subset\R v_c$.  For $j\in O\setminus\{b,c\}$ the line
$\R v_j$ misses $Q\setminus\{0\}$, so $\pi_j$ has constant sign on
$Q\setminus\{0\}$; as $r'\subset K_O\subseteq\{\pi_j\ge0\}$ and
$r'\not\subset\R v_j$, that sign is positive, and in particular $\pi_j(d)>0$, so
$d$ is safe for $j$ (using $U_j\in\cK_j$).  \textup{(3)} $Q$ is adjacent to
$K_O$ along $r'$ and crosses the line $\R v_b$, because
$\angle(r',r_c)<\pi$; hence $Q\subseteq\{\pi_b\le0\}$, and since
$r_c\not\subset\R v_b$ we get $\pi_b(d)<0$.  \textup{(4)} For $j\notin O$ with
$j\ne c$ we have $\cK_j=\{U_j,V_j\}$ and $\pi_j(d)\ne0$, so $d$ is safe for $j$.
\end{proof}

\begin{lem}[Propagation]\label{lem:propagate2}
Let $G\colon\Z^2\to\Fp$ have period $kv$ with $v$ primitive and $k\ge1$, let
$(v,u)$ be a basis of $\Z^2$ with $\det(v,u)=1$ and put $\pi:=\det(v,\cdot)$.
Let $d\in\Z^2$ satisfy $\delta:=-\pi(d)>0$.  If every subsequential limit of the
sequence $(T^{nd}G)_{n\ge1}$ is doubly periodic, then there are $\beta\in\Z$ and
a doubly periodic field $\theta_*$ with $G=\theta_*$ on $\{\pi\le\beta\}$.
\end{lem}

\begin{proof}
Let $\Omega$ be the set of limit points of the sequence: $y\in\Omega$ if and only
if $T^{n_kd}G\to y$ pointwise for some $n_1<n_2<\cdots$.  By compactness of
$\Fp^{\Z^2}$, $\Omega\ne\varnothing$; it is closed; and
$T^{d}\Omega=T^{-d}\Omega=\Omega$, since replacing the indices $n_k$ by
$n_k\pm1$ does not affect the existence of a subsequential limit.  Every element
of $\Omega$ has period $kv$, each $T^{nd}G$ having it and pointwise limits
keeping it, and is doubly periodic by hypothesis.  Call $\{\pi=t\}$ the
\emph{$t$-th row} and put $B_W:=\{sv+tu:|s|,|t|\le W\}$.  Two fields of period
$kv$ that agree at $k$ consecutive lattice points of a row agree on the whole
row; so agreement on $B_W$ with $W\ge k$ implies agreement on the rows $-W$ to
$W$.  For a doubly periodic $y$ of period $kv$ set
\[\Tr(y):=\min\{\pi(\lambda):\lambda\in\Per(y),\ \pi(\lambda)>0\},\]
which is defined because $\Per(y)$ has finite index.

\emph{Step 1: $\Omega$ is finite.}  First, $\{y\in\Omega:\Tr(y)\le P\}$ is finite
for every $P$.  Indeed take $\lambda\in\Per(y)$ with $\pi(\lambda)=\Tr(y)$;
subtracting a multiple of $kv$ we may assume $\lambda=s'v+\Tr(y)\,u$ with
$0\le s'<k$.  For any $z$ write $\pi(z)=m\Tr(y)+r$ with $0\le r<\Tr(y)$; then
$y(z)=y(z-m\lambda)$ and $\pi(z-m\lambda)=r$.  So $y$ is determined by $\lambda$
together with the values of $y$ on
$\{sv+tu:0\le s<k,\ 0\le t<\Tr(y)\}$, and there are only finitely many such
data.

Suppose $\Omega$ were infinite.  Then there are $y_n\in\Omega$ with
$\Tr(y_n)\to\infty$; passing to a subsequence, $y_n\to\theta\in\Omega$.  Now
$\theta$ is doubly periodic; put $P:=\Tr(\theta)$ and choose
$\varpi\in\Per(\theta)$ with $\pi(\varpi)=P$.  From $y_n\to\theta$ we may choose
$W_n\to\infty$ with $y_n=\theta$ on $B_{W_n}$, hence, once $W_n\ge k$, on the
rows $-W_n$ to $W_n$.

Call the $t$-th row \emph{$\varpi$-good for $y_n$} if $y_n(z+\varpi)=y_n(z)$ for
every $z$ with $\pi(z)=t$.  The rows $-W_n$ to $W_n-P$ are $\varpi$-good: both
$z$ and $z+\varpi$ lie in rows where $y_n$ and $\theta$ agree, and
$\theta(z+\varpi)=\theta(z)$.  Let $J_n$ be the largest integer interval
containing $[-W_n,W_n-P]$ and consisting of $\varpi$-good rows.  If
$J_n=\Z$ then $\varpi\in\Per(y_n)$ and $\Tr(y_n)\le P$, contradicting
$\Tr(y_n)\to\infty$ for large $n$.  So $J_n=[a_n,b_n]$ has at least one finite
endpoint; passing to a subsequence we may assume $a_n$ finite for every $n$, the
case of finite $b_n$ being symmetric (replace $T^{q_nd}$ below by $T^{-q_nd}$
and $\{\pi\ge\delta\}$ by $\{\pi\le-\delta\}$).  Then $a_n\le-W_n$ and the row
$a_n-1$ is not $\varpi$-good: there is $z_n$ with $\pi(z_n)=a_n-1$ and
$y_n(z_n+\varpi)\ne y_n(z_n)$.

Put $q_n:=\lceil-a_n/\delta\rceil\ge1$, so that $r_n:=a_n+q_n\delta\in[0,\delta)$,
and put $y_n':=T^{q_nd}y_n\in\Omega$.  Since $\pi(z+q_nd)=\pi(z)-q_n\delta$, the
$t$-th row is $\varpi$-good for $y_n'$ exactly when the row $t-q_n\delta$ is
$\varpi$-good for $y_n$.  Hence the rows $r_n$ to $b_n+q_n\delta$ are
$\varpi$-good for $y_n'$, where
$b_n+q_n\delta\ge(W_n-P)+(-a_n)\ge2W_n-P\to\infty$, while the row $r_n-1$ is
not: the point $z_n':=z_n-q_nd$ satisfies $\pi(z_n')=r_n-1$ and
$y_n'(z_n'+\varpi)\ne y_n'(z_n')$.  Using the $kv$-periodicity, translate $z_n'$
along $v$ into $\{sv+tu:0\le s<k\}$; this yields a point $z_n''$ of the finite
set $F:=\{sv+tu:0\le s<k,\ -1\le t\le\delta-2\}$ with
$y_n'(z_n''+\varpi)\ne y_n'(z_n'')$.  Passing to a further subsequence we may
assume $z_n''=z_*$ and $r_n=r_*$ constant and $y_n'\to y\in\Omega$, $\Omega$
being closed.

Then: \textup{(i)} $y(z+\varpi)=y(z)$ whenever $\pi(z)\ge\delta$ — for large $n$
one has $r_n\le\delta-1<\pi(z)\le b_n+q_n\delta$, so the row $\pi(z)$ is
$\varpi$-good for $y_n'$, and one lets $n\to\infty$; and \textup{(ii)}
$y(z_*+\varpi)\ne y(z_*)$, the two values lying in the finite set $\Fp$ and
converging, so that the inequality survives in the limit.  But $y\in\Omega$ is
doubly periodic: taking $\lambda\in\Per(y)$ with $\pi(\lambda)>0$ and $M$ with
$\pi(z_*+M\lambda)\ge\delta$, part~(i) gives
\[y(z_*+\varpi)=y(z_*+\varpi+M\lambda)=y(z_*+M\lambda)=y(z_*),\]
contradicting (ii).  So $\Omega$ is finite.

\emph{Step 2: eventual agreement.}  Put $\Lambda:=\bigcap_{y\in\Omega}\Per(y)$,
of finite index, and choose $N'\ge1$ with $N'\Z^2\subseteq\Lambda$; every
element of $\Omega$ then has period $N'u$.  Put $W:=\max(k,N'+\delta)$.

\textup{(a)} For all large $n$ the field $T^{nd}G$ agrees on $B_W$ with some
element of $\Omega$.  Otherwise there are $n_1<n_2<\cdots$ with
$T^{n_kd}G|_{B_W}$ outside the finite set $\{y|_{B_W}:y\in\Omega\}$; passing to
a convergent subsequence with limit $y\in\Omega$ gives
$T^{n_kd}G|_{B_W}=y|_{B_W}$ for large $k$, a contradiction.  For such $n$ choose
$y^{(n)}\in\Omega$ with $T^{nd}G=y^{(n)}$ on $B_W$, hence, $W\ge k$, on the rows
$-W$ to $W$.

\textup{(b)} Let $n\ge n_0$ be such $n$.  For $\pi(z)\in[-W+\delta,W]$ one has
$\pi(z+d)\in[-W,W]$, so
\[y^{(n+1)}(z)=(T^{(n+1)d}G)(z)=(T^{nd}G)(z+d)=y^{(n)}(z+d)=(T^{d}y^{(n)})(z).\]
Thus $T^dy^{(n)}$ and $y^{(n+1)}$ agree on $2W-\delta+1\ge N'$ consecutive
entire rows.  Both lie in $\Omega$, since $T^d\Omega=\Omega$, and both have
period $N'u$; any row can be translated by a multiple of $N'u$ into those $N'$
consecutive rows, so the two fields agree everywhere: $T^dy^{(n)}=y^{(n+1)}$.

\textup{(c)} Hence $y^{(n)}=T^{(n-n_0)d}y^{(n_0)}$ for $n\ge n_0$.  For
$\pi(w)\in[-n\delta-W,\,-n\delta+W]$ one has $\pi(w-nd)\in[-W,W]$, so
\[G(w)=(T^{nd}G)(w-nd)=y^{(n)}(w-nd)
  =y^{(n_0)}\bigl(w-nd+(n-n_0)d\bigr)=\bigl(T^{-n_0d}y^{(n_0)}\bigr)(w).\]
As $n$ runs over the integers $\ge n_0$, the row intervals
$[-n\delta-W,-n\delta+W]$ cover $(-\infty,-n_0\delta+W]$, consecutive intervals
overlapping because $2W\ge\delta$.  Take $\theta_*:=T^{-n_0d}y^{(n_0)}$, which is
doubly periodic, and $\beta:=-n_0\delta+W$.
\end{proof}

\begin{proof}[Proof of Theorem~\textup{\ref{thm:secondhalf}}]
Induct, starting from $O:=I$, the set of components for which no $V_j$ has yet
been produced.  While $O\ne\varnothing$, choose $N$ as in the paragraph on safe
directions — it depends on the current $\cK_j$ — and take $b\in O$, $c$ and $d$
as in Lemma~\ref{lem:available}.  By Proposition~\ref{prop:twolimits} every
subsequential limit of $(T^{nd}G_b)_{n\ge1}$ is doubly periodic, so
Lemma~\ref{lem:propagate2}, applied to $G:=G_b$, $v:=v_b$, $u:=u_b$ and this $d$
(for which $\pi_b(d)<0$), yields $\beta$ and a doubly periodic $\theta_*$ with
$G_b=\theta_*$ on $\{\pi_b\le\beta\}$.  Put
$V_b:=\{\pi_b\le\min(\beta,t_b-1)\}$.  Then $V_b\cap U_b=\varnothing$ and
$G_b|_{V_b}$ is fully periodic, with the periods $k_bv_b$ and any
$\lambda\in\Per(\theta_*)$ with $\pi_b(\lambda)<0$, neither of which leaves
$V_b$.  Add $V_b$ to $\cK_b$ and remove $b$ from $O$.  Each step decreases
$|O|$ by one, so after $|I|$ steps $O=\varnothing$ and every component has the
required $V_i$.
\end{proof}

\subsection{Normalisation}\label{ssec:normalisation}

\begin{lem}[Normalisation]\label{lem:normalise}
Let $\theta\colon\Z^2\to\Fp$ be non-periodic and suppose
$\theta=\sum_{i=1}^mG_i$, where each $G_i\colon\Z^2\to\Fp$ has a non-zero period
$h_i$ and there are disjoint half-planes $U_i,V_i$ on which $G_i$ is fully
periodic.  Then $\theta$ is a star configuration.
\end{lem}

Call a family $(G_i,h_i,U_i,V_i)_{i=1}^m$ with these properties an
\emph{admissible decomposition} of $\theta$, of length $m$.  Since $\theta$ is
non-periodic, every admissible decomposition has $m\ge2$: for $m=1$ we would
have $\theta=G_1$ periodic.

\begin{proof}
\emph{Step $(\gamma)$: alignment of the components that are not doubly
periodic.}  Let $G$ have a non-zero period $h$, not be doubly periodic, and be
fully periodic on disjoint half-planes $U$ with normal $n$ and $V$ with normal
$n'$.  \emph{Claim: $\langle h,n\rangle=0$.}  Let $\widetilde G$ be the global
extension of $G|_U$ from Lemma~\ref{lem:extend} and put
$\Gamma_U:=\Per(\widetilde G)$, of finite index.  Suppose
$\langle h,n\rangle\ne0$; replacing $h$ by $-h$ we may assume
$\langle h,n\rangle>0$.  For $z\in U$ we have $z+h\in U$, so
$\widetilde G(z+h)=G(z+h)=G(z)=\widetilde G(z)$: the $\Gamma_U$-periodic
function $\widetilde G(\cdot+h)-\widetilde G$ vanishes on $U$.  But a
$\Gamma_U$-periodic function vanishing on a half-plane vanishes everywhere:
$\Gamma_U$ contains some $N'\Z^2$ and hence a vector $g$ with
$\langle g,n\rangle>0$, and any $z$ is translated into $U$ by $z+Ng$ for $N$
large.  So $\widetilde G$ has period $h$.  Now for any $z$ pick $N$ with
$z+Nh\in U$: then $G(z)=G(z+Nh)=\widetilde G(z+Nh)=\widetilde G(z)$, that is
$G=\widetilde G$ is doubly periodic, a contradiction.

Hence $n\perp h$.  Write $h=kv$ with $v$ primitive and $k\ge1$, after changing
the sign of $h$ if necessary, and put $\pi=\det(v,\cdot)$.  Since
$v^\perp=(-v_2,v_1)$ satisfies $\langle z,v^\perp\rangle=\pi(z)$ and
$n\parallel v^\perp$, we have $\langle z,n\rangle=\kappa\,\pi(z)$ for some
$\kappa\in\R^\times$, so $U$ is either $\{\pi<\ell\}$ or $\{\pi>r\}$ with
$\ell,r\in\Z$.  The same applies to $V$.  If $U$ and $V$ were both of the type
$\{\pi<\cdot\}$, or both of the type $\{\pi>\cdot\}$, they would intersect; so
one is $\{\pi<\ell\}$ and the other is $\{\pi>r\}$, and disjointness says
exactly $\ell\le r+1$.  Taking $\widetilde L,\widetilde R$ to be the global
extensions of $G$ on these two half-planes, $G$ satisfies \textup{(S1)} with
period $kv$, \textup{(S2)}, and \textup{(S3)}.

\emph{Operation A: absorbing a doubly periodic component.}  Suppose $G_i$ is
doubly periodic and $m\ge2$; pick $k\ne i$.  Since $\Per(G_i)$ has finite index,
$\Per(G_i)\cap\Z h_k$ contains some $h_k'\ne0$, and $G_k':=G_k+G_i$ has period
$h_k'$.  On $U_k$ we have $G_k'=\widetilde G_k+G_i$ by
Lemma~\ref{lem:extend}, the restriction of a doubly periodic field, hence fully
periodic; likewise on $V_k$.  Deleting $G_i$ and replacing $G_k$ by $G_k'$ gives
an admissible decomposition of length $m-1$.

\emph{Operation C: merging parallel components.}  Suppose $G_i,G_k$ with $i\ne
k$ are both not doubly periodic and $h_i\parallel h_k$.  By Step $(\gamma)$
there is a primitive $v$ with $h_i=\pm k_iv$, $h_k=\pm k_kv$, and
$\{U_i,V_i\}=\bigl\{\{\pi<\ell_i\},\{\pi>r_i\}\bigr\}$,
$\{U_k,V_k\}=\bigl\{\{\pi<\ell_k\},\{\pi>r_k\}\bigr\}$, where
$\pi=\det(v,\cdot)$.  Then $G_i+G_k$ has period $\lcm(k_i,k_k)\,v$, equals
$\widetilde L_i+\widetilde L_k$ on $\{\pi<\min(\ell_i,\ell_k)\}$ and
$\widetilde R_i+\widetilde R_k$ on $\{\pi>\max(r_i,r_k)\}$, both restrictions of
doubly periodic fields and hence fully periodic; and these two half-planes are
disjoint, since
$\min(\ell_i,\ell_k)\le\ell_i\le r_i+1\le\max(r_i,r_k)+1$.  Merging gives an
admissible decomposition of length $m-1$.

\emph{Termination and conclusion.}  Starting from the given admissible
decomposition, repeat: if some component is doubly periodic, apply Operation~A;
otherwise, if some pair is parallel, apply Operation~C.  Each step strictly
decreases the length, so the process terminates.  At termination no component is
doubly periodic, the periods are pairwise non-parallel, and $m\ge2$.  By
Step~$(\gamma)$ every component satisfies \textup{(S1)}--\textup{(S3)} and the
primitive directions $v_i$ are pairwise non-parallel.  So $\theta$ is a star
configuration.
\end{proof}

\subsection{The convex Nivat conjecture}\label{ssec:nivat}

\begin{thm}[Convex Nivat conjecture]\label{thm:convex-nivat}
Let $\xi\colon\Z^2\to\cA$ with $\cA$ finite, and suppose there is a non-empty
finite lattice-convex $S$ with $P_\xi(S)\le|S|$.  Then $\xi$ is periodic.
\end{thm}

\begin{proof}
Suppose the conjecture has a counterexample.  Embed the alphabet into $\Z_{>0}$
by an injective encoding, which changes neither the number of patterns nor
periodicity, and use Remark~\ref{rem:minimal} to choose a counterexample of
minimal $\ord$, still written $\xi$; it is non-periodic and has low convex
complexity.

\emph{(0) The annihilator.}  Theorem~\ref{thm:KS} supplies pairwise
non-parallel $h_1,\dots,h_m$ and $\varphi=\prod_i(X^{h_i}-1)$ annihilating
$\xi$; we may take $m=\ord(\xi)$.

\emph{(1) The region.}  Theorem~\ref{thm:coneed} supplies $\ell$, and
Theorem~\ref{thm:colregion} a non-periodic $\vartheta\in X_\xi$ together with a
lattice-convex region $R$ on which $\vartheta$ is fully periodic.

\emph{(2) The decomposition.}  Take any prime $p>\max\cA$; then
Lemma~\ref{lem:modp} gives $\vartheta=\sum_i\bar\vartheta_i$ over $\Fp$ with
each $\bar\vartheta_i$ of period $h_i$, the directions being pairwise distinct.

\emph{(3) The first half-plane.}  Corollary~\ref{cor:firsthalf} gives
$\vartheta=\sum_{i\in I}G_i$ with $|I|\ge2$, no component doubly periodic,
pairwise non-parallel periods, and each $G_i$ fully periodic on $U_i$.

\emph{(4) The second half-plane.}  Theorem~\ref{thm:secondhalf} gives
half-planes $V_i$ disjoint from $U_i$ with $G_i|_{V_i}$ fully periodic.  At this
point the whole of Theorem~\ref{thm:structure} holds.

\emph{(5) The complexity does not increase.}  Since $\cA\subset\dbr{p}$, reading
the values of $\vartheta$ in $\Fp$ merges no colours: as an $\Fp$-valued
configuration, $\vartheta$ has the same pattern set and the same periodicity as
it has as an $\cA$-valued one.  Every finite-window pattern of a configuration
in $X_\xi$ is a pattern of $\xi$, a limit in the product topology agreeing with
some translate on any finite window.  Hence
\begin{equation}
P_\vartheta(S)\ \le\ P_\xi(S)\ \le\ |S| .
\label{eq:complexity}
\end{equation}

\emph{(6) Normalisation.}  $\vartheta$ is non-periodic,
$\vartheta=\sum_{i\in I}G_i$ over $\Fp$ with each $G_i$ of period $h_i$ and
fully periodic on the disjoint half-planes $U_i,V_i$.  By
Lemma~\ref{lem:normalise}, $\vartheta$ is a star configuration.

\emph{(7) Contradiction.}  By Theorem~\ref{thm:main},
$P_\vartheta(S)\ge|S|+1$, contradicting \eqref{eq:complexity}.
\end{proof}

\begin{cor}[Nivat's conjecture]\label{cor:nivat}
If $P_\xi\bigl([1,n]\times[1,k]\bigr)\le nk$ for some $n,k\ge1$, then $\xi$ is
periodic.  (A rectangle is lattice-convex.)
\end{cor}

Theorem~\ref{thm:convex-nivat} and Corollary~\ref{cor:nivat} are
\texttt{Nivat.convex\_nivat} and \texttt{Nivat.nivat\_conjecture} in the Lean
development, whose top-level proof follows Steps (0)--(7) above
(Appendix~\ref{app:lean}).

\appendix

\section{Skeleton of the proof}\label{app:skeleton}

\[
\begin{array}{rl}
\text{\S\ref{sec:firstcase}} &
  \exists a:DI_a\ne0\ \Longrightarrow\ \text{Theorem~\ref{thm:main} at once, no
  convexity needed;}\\
  &\text{otherwise }DI_a\equiv0\text{ and }D\eta=0 .\\[3pt]
\text{\S\ref{sec:spectrum}} &
  \Lambda_i=\text{exceptional spectrum},\ A=\prod_iA_i(T^{v_i}),\
  Z=\Newt(A);\\
  & f(T)\eta=\text{const}\Rightarrow A\mid f\ \Longrightarrow\
  \dim U_S\ge|S|-|R_Z(S)|+1 .\\[3pt]
\text{\S\ref{sec:background}} &
  Ae\text{ agrees on adjacent realised sectors}\Rightarrow Ae_\theta=b,\
  \text{doubly periodic, all }H_i\text{-invariant}.\\[3pt]
\text{\S\ref{sec:twopoint}} &
  f_i=Q_i\eta\ne0\text{ and transversally bounded}\ \Longrightarrow\
  \exists d:\ J(d,\cdot)=D_z[\eta(z)\eta(z+d)]\ne0 .\\[3pt]
\text{\S\ref{sec:quotient}} &
  A(T_d)J=A(T_zT_d^{-1})J=0;\ \ R/(a,c)\text{ spanned by }Y^{Z-Z}
  \ \Longrightarrow\ d\in(Z-Z)\cap\Z^2 .\\[3pt]
\text{\S\ref{sec:zonotope}} &
  (Z-Z)\cap\Z^2=(Z\cap\Z^2)-(Z\cap\Z^2)\ \Longrightarrow\ d=q_1-q_0,\
  q_0,q_1\in Z .\\[3pt]
\text{\S\ref{sec:quadratic}} &
  |R_Z(S)|\text{ quadratic observables independent modulo }U_S
  \ \Longrightarrow\ P_\theta(S)\ge|S|+1 .\\[3pt]
\text{\S\ref{sec:reduction}} &
  \text{\cite{KS},\cite{Colle23a}}
  \xrightarrow{\text{Prop.\ \ref{prop:firsthalf}}}U_i
  \xrightarrow{\text{Thm.\ \ref{thm:secondhalf}}}V_i
  \xrightarrow{\text{Lem.\ \ref{lem:normalise}}}\text{star configuration}\\
  &\xrightarrow{\text{Thm.\ \ref{thm:main}}}\text{convex Nivat}
  \Longrightarrow\text{Nivat}.
\end{array}
\]

Dependencies: \S\S\ref{sec:notation}--\ref{sec:quadratic} are self-contained,
using only the standard facts that $\C[T^\pm]$ is a unique factorisation domain,
Newton-polygon additivity (Ostrowski), the Chinese remainder theorem, the
Nullstellensatz, and the existence of unimodular triangulations of lattice
polygons.  The external inputs to \S\ref{sec:reduction} are Kari and Szabados
\cite{KS} and Colle \cite{Colle23a} only: the first supplies the annihilator and
the periodic decomposition, the second a line both of whose directions are
one-sided nonexpansive (\cite[Theorem~1.9]{Colle23a}) and a non-periodic
configuration in the orbit closure that is fully periodic on a region
(the proof of \cite[Lemma~4.6]{Colle23a}).  The $\Fp$ and convex form of Szabados' theorem,
used by Proposition~\ref{prop:twolimits}, is proved in
Appendix~\ref{app:szabados}.

\section{Independent edge-increment results, not part of the proof
chain}\label{app:edge}

Before the argument of \S\S\ref{sec:notation}--\ref{sec:reduction} took its
present form, the same question was approached along a different route, by
comparing $P_\theta(S)$ with $P_\theta(S\setminus E)$ for an edge $E$ of the
polygon $\Conv(S)$ and controlling the increment through the combinatorics of
the edge directions of $\Conv(S)$ against the directions $v_i$.  That route
produced a number of statements of independent interest: an exact identity
for the lattice width of $\Conv(S)$ in a given direction in terms of its edge
data, a divisibility relation between the determinants of three consecutive
edge directions, a propagation principle for sums of functions each of which
depends on a single one of the projections $\pi_i$ and which vanish on enough
consecutive layers, a description of the
periods surviving after repeated differencing, a count of the far-field
patterns along an edge, and lower bounds of the form
$P_\theta(S)-P_\theta(S\setminus E)\ge2$ for short edges and of the form
$P_\theta(S)-P_\theta(S\setminus E)\ge$ (a row-length quantity) in general.
None of them is used above, their proofs are not reproduced here, and the
notation they require is not set up in this paper; they will be treated
separately.  Nothing in \S\S\ref{sec:notation}--\ref{sec:reduction} depends
on any of them.

\section{Numerical verification}\label{app:numerics}

The scripts are \texttt{star.py} (construction, enumeration and all checks),
\texttt{verify.py} and \texttt{verify2.py} (the two randomised experiments),
and \texttt{lemma31.py} (the counterexample of \S\ref{app:counterexample}).
They depend only on \texttt{numpy} and are supplied as ancillary files together
with their output for the stated seeds.  The reported experimental values are
read from that output; the supplementary exact checks described below are
included with the source material.  The latter include
\texttt{check\_case\_support.py}, \texttt{check\_periods\_and\_spectra.py},
and \texttt{check\_encodings.py}; their locations and saved results are listed
in the accompanying \texttt{BUILD.md}.

\subsection{Method}\label{app:method}

\begin{itemize}
\item \emph{Configurations.}  Each component is specified by a primitive
direction $v_i$, a tangential period $k_i$, an arbitrary value table on the
transition strip $[\ell_i,r_i]$, and left and right tails $L_i,R_i$ given, in
the basis $(v_i,u_i)$ of \S\ref{ssec:conv}, as doubly periodic tables in
$(t\bmod a,\ s\bmod k_i)$.  Axioms \textup{(S1)} and \textup{(S3)} hold by
construction.  Axiom \textup{(S2)} is decided exactly: $F_i$ is doubly
periodic if and only if it coincides with the global extension of its left
tail (a doubly periodic function vanishing on a half-plane vanishes
identically), and that coincidence is tested on the strip and on one full
period beyond it.  With $M$ the least common multiple of all the $k_j$ and of
all the tail period indices, the lattice $\Gamma$ of \S\ref{ssec:conv}
contains $M\Z^2$, and the program uses $H_i=Mk_iv_i$; by
Remark~\ref{rem:dichotomy-stable} this choice, which need not be the minimal
one, does not affect the Case~A/Case~B decision, and it does not affect
$\Lambda_i$ either, since refining the period only adds zero Fourier
coefficients.

\item \emph{Pattern counts.}  The $S$-patterns are enumerated at all anchors
of the square box of radius $R=90$ and again of radius $R=180$.
The entries labelled $P_\theta(S)$ below are these observed counts, hence lower
bounds for the global complexity.  The two counts agreed on every reported
window, but this agreement alone does not certify exhaustiveness: that would
require covering the strip crossings and all periodic phase classes outside
them, with the window offsets included.

\item \emph{Deciding the case.}  $DI_a$ is computed on the box of radius $40$
from the formal expansion \eqref{eq:Dexp}.  A non-zero value certifies Case~A;
an all-zero result only suggests Case~B unless the box covers the support.
For the reported configurations, the case assignments were additionally
confirmed on boxes containing all pairwise intersections of the strips
\[\ell_i-\max_{C\subseteq[m]}\pi_i(H_C)\le\pi_i(z)
  \le r_i-\min_{C\subseteq[m]}\pi_i(H_C).\]
The strip argument of Lemma~\ref{lem:finite-support}, with $W=\{0\}$, places
the entire support of every $DI_a$ inside this union.

\item \emph{Spectrum and zonotope.}  The fields $\Xi_i^\pm$ and
$G_i^{\sigma,\epsilon}$ are built as in \S\ref{ssec:isolating}; for each
difference field $d_{i,\sigma,\epsilon,a}$ a length-$\kappa_i$ DFT is taken
row by row over all rows within $2M+2$ of the strip (outside the strip the
rows repeat with period $M$).  Frequencies with coefficient of modulus above
$10^{-9}$ give numerical estimates of $\Lambda_i$, $d_i$ and $Z$.  For all
$570$ components in the two experiments, the frequency sets were also checked
exactly by reducing the integer row polynomials modulo the appropriate
cyclotomic polynomials; every set agreed with the DFT output.  The set
$R_Z(S)$ is decided by testing whether all subset sums of the $d_iv_i$,
translated by $r$, lie in $\Conv(S)$.

\item \emph{Ranks.}  The encoding $w$ is a random injection into
$\{1,\dots,100\}$.  Injectivity alone does not imply the spectrum-preservation
condition of Lemma~\ref{lem:encoding}; for the reported encodings that condition
was checked separately by exact cyclotomic remainders and held for all $570$
components.  The rank of the $|\cL|\times(|S|+1)$ matrix
$\bigl(1,\,w(\gamma(s))_{s\in S}\bigr)_{\gamma\in\cL}$ is computed over
$\mathbb F_q$ with $q=2^{31}-1$.  The entries labelled $\dim U_S$ below are
these modular ranks of the sampled pattern matrices; they are lower bounds
for the rational ranks on the full pattern set.  Meeting the lower bound in
Lemma~\ref{lem:budget} is therefore sufficient to verify that inequality for
the sampled configuration.  In Case~B with $R_Z(S)\ne\varnothing$ the matrix is
extended by the columns $\Phi_r(\gamma)=w(\gamma(r+q))\,w(\gamma(r+q+d))$,
$r\in R_Z(S)$, of \S\ref{sec:quadratic}, and the rank $H_2(S)$ of the
extended matrix is compared with $\dim U_S+|R_Z(S)|$
(Lemma~\ref{lem:independence}) and with $|S|+1$ (Theorem~\ref{thm:main}).

\item \emph{Witness and background.}  In Case~B, for every non-zero
$d\in(Z-Z)\cap\Z^2$ the function $J(d,\cdot)$ is computed on the box of
radius $40$ and the non-zero instances are collected
(Proposition~\ref{prop:witness}); the first of them, with a $q$ as in
Corollary~\ref{cor:pair}, is the pair $(d,q)$ used for the $\Phi_r$.  The
polynomial $A$ is built with complex coefficients from the $\Lambda_i$, and
$Ae_\theta$ is checked to be $M\Z^2$-periodic on the box
(Proposition~\ref{prop:Ae-b}).
\end{itemize}

\subsection{Results}\label{app:results}

\noindent\textbf{Experiment 1} (script \texttt{verify.py}, seed $1$).  We
generated $84$ random star configurations with $p\in\{2,3,5\}$, $m\in\{2,3\}$,
directions drawn without
repetition from the $8$ primitive vectors
$(1,0),(0,1),(1,\pm1),(1,\pm2),(2,\pm1)$, $k_i\in\{1,2,3\}$, tail periods of
$1$ or $2$ rows, strip widths $1$ to $3$ with $\ell_i\in\{-1,0,1\}$, all
tables uniformly random over $\mathbb F_p$ subject to \textup{(S2)}; component
$1$ has independently drawn left and right tails, the other components have
$L_i=R_i$.  Windows: the $2\times3$, $3\times3$, $2\times4$ and $4\times2$
rectangles, the $6$-point triangle with vertices $(0,0),(2,0),(0,2)$, the
$8$-point hexagon with vertices $(1,0),(2,0),(3,1),(2,2),(1,2),(0,1)$, the
$5$-point rhombus with vertices $(0,0),(1,1),(2,0),(1,-1)$, and the $3$-point
segment $\{0,1,2\}\times\{0\}$.  Result: $81$ configurations in Case~A and
$3$ in Case~B; $P_\theta(S)\ge|S|+1$ held on all $8$ windows of all $84$
configurations.  The three Case-B instances have $d_i\in\{3,6\}$, so $Z$ is
larger than every window and $R_Z(S)=\varnothing$ throughout; on them
$\dim U_S$ was exactly $|S|+1$ on every window, with $192$, $310$ and $28$
non-zero witnesses $d$ found respectively, and the periodicity residual of
$Ae_\theta$ on the tested box was at most $2\cdot10^{-15}$.  These are instances of
Remark~\ref{rem:empty-RZ}; Experiment~2 targets $R_Z(S)\ne\varnothing$.

\noindent\textbf{Experiment 2} (script \texttt{verify2.py}, seed $7$).  We
generated $157$ configurations of the form $\sum_iw_i\circ\pi_i$, so $k_i=1$
and constant
tails: each $w_i\colon\Z\to\mathbb F_p$ takes a constant value $c_i$ outside a
window of length $1$ to $4$ and arbitrary values inside it, with $L_i=R_i=c_i$;
$p\in\{2,3,5\}$, $m=2$ with probability $0.7$ and $m=3$ with probability
$0.3$, directions drawn from the $11$ primitive vectors
$(1,0),(0,1),(1,\pm1),(1,\pm2),(2,\pm1),(1,\pm3),(3,1)$.  Windows: the
$3\times4$, $4\times4$, $5\times3$ and $4\times3$ rectangles, the $15$-point
triangle with vertices $(0,0),(4,0),(0,4)$, the $17$-point hexagon with
vertices $(1,0),(3,0),(4,2),(3,4),(1,4),(0,2)$, the $13$-point rhombus with
vertices $(0,0),(2,2),(4,0),(2,-2)$, and the $12$-point parallelogram with
vertices $(0,0),(3,0),(5,2),(2,2)$.  Result: $128$ in Case~A and $29$ in
Case~B, all with $d_i=1$; $P_\theta(S)\ge|S|+1$ held on all $8$ windows of all
$157$ configurations, and every check listed in \S\ref{app:method} passed on
every Case-B instance.  Of the $29$ Case-B configurations, $14$ have
$R_Z(S)\ne\varnothing$ for at least one window, giving $35$
configuration--window pairs with $R_Z(S)\ne\varnothing$ and
$|R_Z(S)|\in\{1,2,3,4,5\}$; all $35$ have $m=2$ (no Case-B instance with
$m=3$ and $R_Z(S)\ne\varnothing$ was produced by this generator).
Representative rows of the output:

\begin{center}\small
\begin{tabular}{@{}cllcccccccc@{}}
\toprule
$p$ & directions & $S$ & $|S|$ & $P_\theta(S)$ & $d_i$ & $|R_Z(S)|$ &
$\dim U_S$ & $|S|+1-|R_Z(S)|$ & $H_2(S)$ & $d$\\
\midrule
$3$ & $(1,2),(1,0)$ & $4\times4$ & $16$ & $35$ & $(1,1)$ & $4$ & $13$ & $13$ & $17$ & $(-1,-1)$\\
$3$ & $(1,2),(1,0)$ & hexagon & $17$ & $37$ & $(1,1)$ & $5$ & $13$ & $13$ & $18$ & $(-1,-1)$\\
$3$ & $(1,2),(1,-2)$ & hexagon & $17$ & $45$ & $(1,1)$ & $3$ & $15$ & $15$ & $18$ & $(-1,-1)$\\
$5$ & $(1,-3),(1,0)$ & $4\times4$ & $16$ & $37$ & $(1,1)$ & $2$ & $15$ & $15$ & $17$ & $(-1,1)$\\
$2$ & $(2,-1),(1,2)$ & hexagon & $17$ & $41$ & $(1,1)$ & $2$ & $16$ & $16$ & $18$ & $(-2,0)$\\
$5$ & $(2,1),(1,2)$ & $4\times4$ & $16$ & $54$ & $(1,1)$ & $1$ & $16$ & $16$ & $17$ & $(-2,-2)$\\
\bottomrule
\end{tabular}
\end{center}

\noindent
On all $35$ pairs with $R_Z(S)\ne\varnothing$, the sampled base matrix had
modular rank $|S|+1-|R_Z(S)|$, and the extended matrix had modular rank
$H_2(S)=|S|+1$.  Thus their rank difference was exactly $|R_Z(S)|$, matching
the dimension comparison in \S\ref{sec:quadratic}.  These numerical equalities
alone do not identify the full rational space of affine relations.  On the
$116$ two-component configurations of Experiment~2 the criterion of
Observation~\ref{obs:m2} below agreed with the computed case in all $116$
instances, and the Case-B direction pairs had $|\det(v_1,v_2)|\in
\{2,3,4,5,6,7,10\}$.

\begin{obs}[A criterion for Case~B when $m=2$]\label{obs:m2}
Suppose $m=2$ and both components satisfy $L_i=R_i$; write
$\delta_i=F_i-L_i$ for the deviation field, supported in strip $i$.  Replace
$H_i$ by a sufficiently large multiple, which by
Remark~\ref{rem:dichotomy-stable} does not change the Case~A/Case~B decision,
and expand the four-point template $\{z,z+H_1,z+H_2,z+H_1+H_2\}$ at an anchor
$z$: $\delta_1$ is constant along the $H_1$-pairs and $\delta_2$ along the
$H_2$-pairs, so $De_\theta(z)=0$ if and only if at least one of the two pairs of
constants agrees.  Hence \emph{Case~B holds if and only if
$\supp\delta_1\cap\supp\delta_2=\varnothing$}, that is, no point is deviated by
both components at once.  In particular, when $k_i=1$, the tails are constant
and equal, and $|\det(v_1,v_2)|=1$, one is always in Case~A: the map
$(\pi_1,\pi_2)\colon\Z^2\to\Z^2$ is then surjective and the two deviation sets
must meet.  Under these additional hypotheses, Case~B occurs only for
$|\det|\ge2$.  The sample of Experiment~2 was
generated at random and then sorted by the numerical test; the criterion
explains why Case~B appeared only on direction pairs with $|\det|\ge2$.
\end{obs}

\subsection{A counterexample to the generalisation of
Lemma~\ref{lem:Ae-sectors}}\label{app:counterexample}

(\texttt{lemma31.py}.)  Take $p=2$, $m=3$, $v_1=(1,0)$, $v_2=(0,-1)$,
$v_3=(1,-1)$.  Tails: $L_1=0$ and $R_1=\one[x\text{ even}]$; $L_2=L_3=0$ and
$R_2=R_3=\one[3\mid x]$; the transition strip of each component is constant
$1$, of width $1$.  Here $\pi_2(v_1)=\pi_3(v_1)=1>0$, so
$B_1^+=R_2+R_3\equiv0$ and $B_1^-=L_2+L_3=0$, and the program computes
$\Lambda_1=\{\pm1\}$ (with $\kappa_1=12$ and frequencies $\{0,6\}$), so that
$A_1(T^{v_1})=T^{2v_1}-1$.

\begin{itemize}
\item For the pairs of realised sectors adjacent along the rays $\pm v_1$,
namely $(R_1,R_2,R_3)$ with $(L_1,R_2,R_3)$ and $(R_1,L_2,L_3)$ with
$(L_1,L_2,L_3)$: the difference after applying $A_1(T^{v_1})$ is $0$, up to a
floating-point error of $1.3\cdot10^{-16}$, for both colours.

\item For the mixed choice $(R_1,R_2,L_3)$ against $(L_1,R_2,L_3)$ — here
$B=R_2+L_3=\one[3\mid x]$, and the sign vector $(+,+,-)$ corresponds to the cone
$\{y>0,\ x>0,\ x+y<0\}=\varnothing$, which is not realised: the difference after
applying $A_1(T^{v_1})$ has maximum absolute value $2$.
\end{itemize}

By hand: along the $x$-direction,
$\Delta_1=\one[R_1+B=1]-\one[L_1+B=1]$ takes the values $(-1,0,1,0,1,0)$ for
$x\bmod6=0,\dots,5$, whereas its translate by $x\mapsto x+2$ is
$(1,0,1,0,-1,0)$.  The two differ, so $\Delta_1$ is not killed by
$T^{2v_1}-1$: its spectrum contains cube roots of unity, and those are not in
$\Lambda_1$.  This is why \S\ref{sec:background} can only assert that $Ae$ is
constant across the \emph{realised} sectors, where $B=B_i^\sigma$ — and that is
exactly, and only, what Lemma~\ref{lem:Ae-finite} needs.

\section{The \texorpdfstring{$\Fp$}{Fp} and convex form of Szabados'
theorem}\label{app:szabados}

This appendix proves in full the $\Fp$ and convex version of a theorem of
Szabados.  Szabados \cite[Theorem~1]{Sz} proved that a non-periodic sum of two $\Z$-valued
periodic configurations with different period directions does not have low
rectangular complexity; the convex generalisation is
\cite[Theorem~1.7]{Colle23a}, and the $\Fp$ and convex version is stated, with a
sketch, as \cite[Theorem~2.8]{CG}.  Theorem~\ref{thm:szabados} below is the
latter, proved here in full.  The proof follows Szabados \cite{Sz} —
pigeonhole, balanced sets, and the one-dimensional Morse--Hedlund theorem — with
two changes.  The coefficient ring is $\Fp$: the algebraic steps use addition,
subtraction, equality and periodicity, without division or a property specific
to the characteristic.  The window is an arbitrary lattice-convex set instead
of a rectangle: in
Lemma~\ref{lem:balanced} one deletes, at each step, whichever of the two extreme
rows is the shorter.  In addition, Lemma~3 of \cite{Sz} rests on Lemma~39 of
Kari and Szabados \cite{KS}; we replace it by two self-contained arguments,
Lemma~\ref{lem:strip-period}, which uses the explicit shape of the annihilator,
and step~(6) of Lemma~\ref{lem:propagation}, which uses the two-sided recursion
supplied by the generating property.  Apart from the one-dimensional
Morse--Hedlund theorem, nothing external is quoted here.

\setcounter{subsection}{-1}
\newpage
\subsection{The statements}\label{app:szab-statements}

\begin{thm}\label{thm:szabados}
Let $p$ be a prime, let $\theta_1,\theta_2\colon\Z^2\to\Fp$ have periods
$h_1,h_2\in\Z^2\setminus\{0\}$ respectively, with $h_1$ and $h_2$ linearly
independent over $\R$, and put $\theta=\theta_1+\theta_2$.  If $P_\theta(S)\le|S|$
for some non-empty finite lattice-convex $S$, then $\theta$ is periodic.
\end{thm}

\begin{cor}\label{cor:szabados}
Let $\eta\colon\Z^2\to\Fp$ and let $\eta=\eta_1+\eta_2$ be an $\Fp$-minimal
periodic decomposition, of order $2$ — that is, a decomposition into periodic
$\Fp$-configurations with as few summands as possible, here two.  Then $\eta$
does not have low convex complexity.
\end{cor}

\begin{proof}
Order $2$ forces $\eta$ to be non-periodic, a periodic configuration having a
minimal decomposition of length $1$, and it forces the two periods to point in
different directions, components with parallel periods being mergeable into one,
against minimality.  If $\eta$ had low convex complexity,
Theorem~\ref{thm:szabados} would make it periodic, a contradiction.
\end{proof}

\subsection{Notation}\label{app:szab-notation}

Let $\su$ be the primitive vector in the direction of $h_1$, so that
$h_1=c_1\su$ with $c_1\ge1$, and let $\pi:=\det(\su,\cdot)\colon\Z^2\to\Z$,
which is surjective.  The \emph{rows} are the fibres $L_t:=\pi^{-1}(t)$;
neighbouring lattice points of a row differ by $\su$.  Put $\Delta:=\pi(h_2)\ne0$
and, replacing $h_2$ by $-h_2$ if necessary, assume $\Delta>0$.  For integers
$a\le b$ the \emph{strip} $\mathrm{St}[a,b]:=\{v:a\le\pi(v)\le b\}$, of
\emph{width} $b-a$, is the union of the rows $L_a,\dots,L_b$ and is invariant
under $\pm\su$; we also set $\mathrm{St}[a,a-1]:=\varnothing$.
For a non-empty finite $F\subset\Z^2$ put
$\operatorname{ht}(F):=\max\pi(F)-\min\pi(F)$; $F$ \emph{fits into}
$\mathrm{St}[a,b]$ if some translate of it is contained there, that is if
$\operatorname{ht}(F)\le b-a$.  For a non-empty finite lattice-convex $F$ write
$E_+(F):=F\cap L_{\max\pi(F)}$ and $E_-(F):=F\cap L_{\min\pi(F)}$, the top and
bottom rows; lattice convexity makes every non-empty row of $F$ a run of
consecutive lattice points.

For a finite $F\subset\Z^2$ write
$\mathrm{Pat}(F,\theta):=\{\theta|_{v+F}:v\in\Z^2\}$, the patterns being
regarded as functions $w\mapsto\theta(v+w)$ on $F$, so that
$P_\theta(F)=|\mathrm{Pat}(F,\theta)|$ in the notation of \S\ref{ssec:main},
with $P_\theta(\varnothing)=1$; $X_\theta$
is the orbit closure, so that every finite-window pattern of an $x\in X_\theta$
is a pattern of $\theta$.  ``$x=\theta$ on $U$'' means $x(v)=\theta(v)$ for all
$v\in U$.  Inside this appendix the letter $D$ denotes auxiliary finite sets and
never the difference operator of \S\ref{ssec:conv}, which does not occur here.

\subsection{Two algebraic facts}\label{app:szab-algebra}

For $s\in\Z$, the $h_2$-periodicity of $\theta_2$ gives
\begin{equation}
\theta(v+sh_2)-\theta(v)=\theta_1(v+sh_2)-\theta_1(v),
\quad\text{which is $h_1$-periodic in $v$.}
\label{eq:D-diff}
\end{equation}
Consequently: if $T^{sh_2}\theta$, that is $v\mapsto\theta(v+sh_2)$, agrees with
$\theta$ on a set $F$, then it agrees with $\theta$ on $F+\Z h_1$.

\begin{lem}[Periodicity on a strip implies periodicity]\label{lem:strip-period}
Suppose $b-a\ge\Delta$ and $q\ge1$, and that $\theta(v+q\su)=\theta(v)$ for all
$v\in\mathrm{St}[a,b]$.  Then $\theta$ has period $qh_1$.
\end{lem}

\begin{proof}
The strip is invariant under $+\su$, so $\theta(v+qh_1)=\theta(v)$ for
$v\in\mathrm{St}[a,b]$, applying the hypothesis $c_1$ times.  Put
$\psi(v):=\theta(v+qh_1)-\theta(v)$; since $\theta_1$ has period $h_1$ we have
$\psi(v)=\theta_2(v+qh_1)-\theta_2(v)$, so $\psi$ has period $h_2$, and
$\psi=0$ on $\mathrm{St}[a,b]$.  Write $h_2=gh_2'$ with $h_2'$ primitive and
take any lattice line $L=\{v_0+sh_2':s\in\Z\}$ in the direction $h_2'$.  Then
$\psi|_L$ is periodic in $s$ with period $g$, while
$\pi(v_0+sh_2')=\pi(v_0)+s\,\pi(h_2')$ runs through an arithmetic progression of
common difference $\pi(h_2')=\Delta/g\ge1$; a closed interval $[a,b]$ of length
$b-a\ge\Delta=g\cdot(\Delta/g)$ contains at least $g$ of its terms, so $L$ meets
$\mathrm{St}[a,b]$ in at least $g$ consecutive of its points.  A function of
period $g$ vanishing at $g$ consecutive points vanishes identically, so
$\psi|_L\equiv0$.  As $L$ was arbitrary, $\psi\equiv0$.
\end{proof}

\subsection{Balanced sets}\label{app:szab-balanced}

\begin{dfn}\label{dfn:balanced}
Let $B$ be a non-empty finite lattice-convex set, let $\sigma\in\{+,-\}$ and put
$E:=E_\sigma(B)$.  Call $B$ \emph{$\sigma$-balanced} if
\begin{itemize}
\item[\textup{(i)}] $P_\theta(B)\le|B|$;
\item[\textup{(ii)}] $P_\theta(B)<P_\theta(B\setminus E)+|E|$;
\item[\textup{(iii)}] for every integer
      $t\in[\min\pi(B),\max\pi(B)]$, one has $|B\cap L_t|\ge|E|-1$.
\end{itemize}
\end{dfn}

\begin{lem}[Existence]\label{lem:balanced}
Let $S$ be a non-empty finite lattice-convex set with $P_\theta(S)\le|S|$.  Then
there are a sign $\sigma$ and a $\sigma$-balanced set $B\subseteq S$.
\end{lem}

\begin{proof}
Construct $D_0:=S\supsetneq D_1\supsetneq\cdots\supsetneq D_K=\varnothing$ as
follows: if $D_i\ne\varnothing$, choose $\sigma_i$ with
$|E_{\sigma_i}(D_i)|\le|E_{-\sigma_i}(D_i)|$ — either sign will do when $D_i$
has a single row, where $E_+=E_-=D_i$ — and put
$D_{i+1}:=D_i\setminus E_{\sigma_i}(D_i)$.  Each $D_i$ is again lattice-convex:
after an extreme row has been deleted the remaining points lie in an open
half-plane, so the lattice points of their convex hull still belong to $D_i$ and
do not lie on the deleted row.  Put $f(i):=P_\theta(D_i)-|D_i|$; then $f(0)\le0$
and $f(K)=1$, so there is a smallest $i$ with $f(i+1)>0$, and for it
$f(i)\le0<f(i+1)$.  Let $B:=D_i$, $\sigma:=\sigma_i$ and $E:=E_\sigma(B)$.
Condition~(i) is $f(i)\le0$.  For (ii),
$P_\theta(B)\le|B|=|B\setminus E|+|E|<P_\theta(B\setminus E)+|E|$.

For (iii), write $n=|E|$ and $E'=E_{-\sigma}(B)$, so $|E'|\ge n$.
If $n=1$, the required lower bound is zero and holds for every integer $t$.
If $B$ has only one row, the only relevant row contains $n$ points, so (iii)
also holds.  Hence assume $n\ge2$ and $\min\pi(B)<\max\pi(B)$.
For a real
$t\in[\min\pi(B),\max\pi(B)]$ let $\lambda(t)$ be the length of the intersection
of the convex polygon $\Conv(B)$ with the line $\{\pi=t\}$; $\lambda$ is concave,
by the concavity of the length of parallel chords of a convex set.  On the two
extreme rows that intersection is exactly $\Conv(E)$, respectively $\Conv(E')$,
a supporting line meeting the polygon in the convex hull of the vertices it
carries, of lengths $(n-1)|\su|$ and $(|E'|-1)|\su|\ge(n-1)|\su|$.  Hence
$\lambda(t)\ge(n-1)|\su|$ for every $t$ in that interval.  For every integer
$t$ there, surjectivity of $\pi$ gives a lattice point on $L_t$, and the lattice
points of $L_t$ are spaced $|\su|$ apart.  Its intersection with $\Conv(B)$ is
a closed segment of length at least $(n-1)|\su|$, so it contains at least
$n-1$ lattice points.  These belong to $B$ by lattice convexity.  This proves
(iii) at every intermediate integer height, including those not assumed in
advance to meet $B$.
\end{proof}

\subsection{Ambiguity strips and the propagation lemma}\label{app:szab-prop}

\begin{dfn}\label{dfn:ambiguity}
Say that $\theta$ \emph{has the $+$-ambiguity strip} $\mathrm{St}[a,b]$, where
$a\le b$, if there is an $x\in X_\theta$ with $x=\theta$ on
$\mathrm{St}[a,b-1]$ — an empty condition when $a=b$ — and $x\ne\theta$ at some
point of $L_b$.  The $-$-ambiguity strips are defined symmetrically, replacing
$\su$ by $-\su$.
\end{dfn}

\begin{lem}[Propagation]\label{lem:propagation}
Let $B$ be $+$-balanced and put $n:=|E_+(B)|$.  If $\theta$ has some
$+$-ambiguity strip $\mathrm{St}[a,b]$ into which $B$ fits, that is with
$b-a\ge\operatorname{ht}(B)$, then $\theta$ is periodic.
\end{lem}

\begin{proof}
\emph{(1) Placing $B$.}  Translate $B$ so that $\max\pi(B)=b$; then
$B\subseteq\mathrm{St}[a,b]$, and $E:=E_+(B)=B\cap L_b=\{e_1,\dots,e_n\}$ with
$e_{l+1}=e_l+\su$, while $B_0:=B\setminus E\subseteq\mathrm{St}[a,b-1]$, possibly
empty.  Put $t_-:=\min\pi(B)$.  Take a witness $x\in X_\theta$ of the ambiguity:
$x=\theta$ on $\mathrm{St}[a,b-1]$ and $x\ne\theta$ somewhere on $L_b$.  Note
that $B_0+i\su\subseteq\mathrm{St}[a,b-1]$ for every $i\in\Z$, so $x=\theta$ on
each $B_0+i\su$.

\emph{(2) Two determination rules.}  For $0\le i\le n$ put
$D_i:=B_0\cup\{e_1,\dots,e_i\}$ and $D'_i:=B_0\cup\{e_{n-i+1},\dots,e_n\}$.  By
(ii) of Definition~\ref{dfn:balanced}, $P_\theta(D_n)-|D_n|<P_\theta(D_0)-|D_0|$,
so there is a $k\in\{0,\dots,n-1\}$ with
$P_\theta(D_{k+1})-|D_{k+1}|<P_\theta(D_k)-|D_k|$, that is
$P_\theta(D_{k+1})<P_\theta(D_k)+1$; the restriction map
$\mathrm{Pat}(D_{k+1},\theta)\to\mathrm{Pat}(D_k,\theta)$ being onto, this forces
$P_\theta(D_{k+1})=P_\theta(D_k)$ and makes that map a bijection.  Hence:
\begin{itemize}
\item[\textup{(R$_+$)}] for every $y\in X_\theta$ and $v\in\Z^2$, the value
$y(e_{k+1}+v)$ is determined by $y|_{D_k+v}$.
\end{itemize}
In the same way there is a $k'\in\{0,\dots,n-1\}$ with
\begin{itemize}
\item[\textup{(R$_-$)}] for every $y\in X_\theta$ and $v\in\Z^2$, the value
$y(e_{n-k'}+v)$ is determined by $y|_{D'_{k'}+v}$.
\end{itemize}

\emph{(3) Claim: $x|_{D_k+i\su}\ne\theta|_{D_k+i\su}$ for every $i\in\Z$.}
Suppose $x=\theta$ on $D_k+i_0\su$.  If $k=0$, then $D_0=B_0$, and (1) together
with \textup{(R$_+$)} gives $x(e_1+i\su)=\theta(e_1+i\su)$ for every $i$, that is
$x=\theta$ on $L_b$, against the choice of the witness.  So let $k\ge1$.  By
\textup{(R$_+$)}, $x(e_{k+1}+i_0\su)=\theta(e_{k+1}+i_0\su)$; hence $x=\theta$ on
$D_k+(i_0+1)\su$, whose points are those of $B_0+(i_0+1)\su$, by (1), together
with the points $e_l+(i_0+1)\su=e_{l+1}+i_0\su$ for $1\le l\le k$, which lie in
$D_k+i_0\su$ when $l\le k-1$ and are the point just treated when $l=k$.  By
induction $x=\theta$ on $D_k+i\su$ for every $i\ge i_0$, hence on the half-line
$\{e_1+s\su:s\ge i_0\}\subset L_b$.  Now go left.  If $k'=0$, then (1) and
\textup{(R$_-$)} give $x=\theta$ on $\{e_n+i\su:i\in\Z\}=L_b$, a contradiction.
If $k'\ge1$, choose $i_1$ with $n-k'+i_1\ge i_0$; the points
$e_1+(n-k'+i_1)\su,\dots,e_1+(n-1+i_1)\su$ of the top row of $D'_{k'}+i_1\su$
then lie on that half-line, so $x=\theta$ on $D'_{k'}+i\su$ for every $i\ge i_1$.
By \textup{(R$_-$)}, $x(e_{n-k'}+i_1\su)=\theta(e_{n-k'}+i_1\su)$, whence
$x=\theta$ on $D'_{k'}+(i_1-1)\su$, whose top-row points
$e_l+(i_1-1)\su=e_{l-1}+i_1\su$, $n-k'+1\le l\le n$, either lie in
$D'_{k'}+i_1\su$ or are the point just treated.  Descending induction gives
$x=\theta$ on $D'_{k'}+i\su$ for every $i\in\Z$, and the top rows of these sets
cover all of $L_b$, against the choice of the witness.  The claim follows.

\emph{(4) Counting.}  By (3) and (1), for each $i$ the patterns $x|_{B+i\su}$ and
$\theta|_{B+i\su}$ are two $B$-patterns of $\theta$ which agree on $B_0+i\su$
and differ somewhere.  So every pattern $\gamma$ in
$\Gamma:=\{\theta|_{B_0+i\su}:i\in\Z\}\subseteq\mathrm{Pat}(B_0,\theta)$ has at
least two extensions to a $B$-pattern, $N(\gamma)\ge2$.  But
\[\sum_{\gamma\in\mathrm{Pat}(B_0,\theta)}\bigl(N(\gamma)-1\bigr)
  =P_\theta(B)-P_\theta(B_0)<|E|=n\]
by (ii), so $|\Gamma|\le n-1$; in particular $n\ge2$, since otherwise $\Gamma$
would be empty, which it is not.

\emph{(5) Periodicity of the rows of $B_0$.}  Assume $B_0\ne\varnothing$; when
$B_0=\varnothing$, one has $t_-=b$, the initial strip
$\mathrm{St}[t_-,b-1]$ is empty, and one takes $Q_0:=1$.
For every integer $t$ with $t_-\le t\le b-1$, condition (iii) and $n\ge2$
from (4) give $|B\cap L_t|\ge n-1\ge1$.  Thus no row of this strip is missing.
By lattice convexity, $B\cap L_t=B_0\cap L_t$ is a run of $m_t\ge n-1$
consecutive points $w_t,w_t+\su,\dots,w_t+(m_t-1)\su$.  Put
$\xi_t(s):=\theta(w_t+s\su)$ for $s\in\Z$.  As $i$ varies, $\theta|_{B_0+i\su}$
takes at most $n-1$ values, so its restriction to the first $n-1$ points of row
$t$ — that is, the word of length $n-1$ read in $\xi_t$ from position $i$ — takes
at most $n-1$ values, and $\xi_t$ has at most $n-1$ subwords of length $n-1$.  By
the one-dimensional Morse--Hedlund theorem \cite{MH38} (a bi-infinite sequence
with at most $m$ subwords of some length $m$ is periodic; see also
\cite[Theorem~1.1]{Colle23a}), $\xi_t$ is periodic.  Let $Q_0$ be a common
multiple of the periods of these rows; then $\theta(v+Q_0\su)=\theta(v)$ for all
$v\in\mathrm{St}[t_-,b-1]$.

\emph{(6) Extending upwards one row at a time.}  If $\Delta<b-t_-$, then
$b-1-t_-\ge\Delta$, so the strip obtained in (5) already contains
$\mathrm{St}[t_-,t_-+\Delta]$; take $Q^*:=Q_0$ and proceed to (7).
Otherwise $\Delta\ge b-t_-$.  Suppose $R\ge b$ and $\theta$
has period $Q\su$ on $\mathrm{St}[t_-,R-1]$, which for $R=b$ is step~(5) with
$Q=Q_0$.  Choose $w$ with $\pi(w)=R-b$, so that the top row of $B+w$ lies on
$L_R$ and $B_0+w+i\su\subseteq\mathrm{St}[t_-,R-1]$; then the pattern
$\theta|_{B_0+w+i\su}$ depends only on $i\bmod Q$.  Put
$y_s:=\theta(e_1+w+s\su)$ and apply \textup{(R$_+$)} and \textup{(R$_-$)} to
$y:=\theta$ and $v=w+i\su$:
\begin{gather*}
y_{i+k}\ \text{is determined by}\ \bigl(i\bmod Q;\ y_i,\dots,y_{i+k-1}\bigr),\\
y_{i-1}\ \text{is determined by}\ \bigl(i\bmod Q;\ y_i,\dots,y_{i+k'-1}\bigr),
\end{gather*}
the second being \textup{(R$_-$)} at $v=w+(i-n+k')\su$.  Put
$K:=\max(k,k',1)$ and let
$\sigma_i:=(i\bmod Q;\,y_i,\dots,y_{i+K-1})\in\Z/Q\times\Fp^{K}$.  Then
$\sigma_{i+1}$ is determined by $\sigma_i$, and so is $\sigma_{i-1}$.  Let
$\Phi$ be the forward map, defined on the set $\Sigma$ of states that actually
occur; backward determinacy makes $\Phi$ injective on $\Sigma$, since
$\Phi(\sigma)=\Phi(\sigma')$ forces $\sigma=\sigma'$, and $\Sigma$ is finite, so
$\Phi$ is a bijection of $\Sigma$ and the bi-infinite orbit $(\sigma_i)$ is
periodic.  Its period $Q'$ is a multiple of $Q$, by the phase component, and
$y_{s+Q'}=y_s$.  Hence $\theta$ has period $Q'\su$ on $\mathrm{St}[t_-,R]$.
Iterating from $R=b$ up to $R=t_-+\Delta$, finitely many steps, gives a $Q^*$
for which $\theta$ has period $Q^*\su$ on the strip
$\mathrm{St}[t_-,t_-+\Delta]$, of width $\Delta$.

\emph{(7)}  By Lemma~\ref{lem:strip-period}, $\theta$ has period $Q^*h_1$.
\end{proof}

\begin{rem}\label{rem:szab-steps}
Step~(5) of Lemma~\ref{lem:propagation} is the argument of \cite[Lemma~4]{Sz}
and of \cite[Lemma~2.24]{CyrKra15}; steps~(6) and (7) replace \cite[Lemma~3]{Sz},
which appeals to Lemma~39 of \cite{KS}, and use nothing beyond the generating
property and the explicit annihilator.
\end{rem}

\subsection{Proof of Theorem \texorpdfstring{\ref{thm:szabados}}{D.1}}
\label{app:szab-proof}

Suppose $\theta$ is not periodic.  Replacing $h_2$ by $ch_2$ for an integer
$c\ge1$, which changes neither $\pi$ nor $\su$, we may assume
$\Delta\ge\operatorname{ht}(S)+1$.  Put
\[D:=\{\alpha h_1+\beta h_2:\ \alpha,\beta\in[0,1)\}\cap\Z^2,\qquad
D_j:=D+jh_2,\qquad S_j:=\mathrm{St}\bigl[j\Delta,(j+1)\Delta-1\bigr].\]
Every $v\in\Z^2$ is uniquely $v=\alpha h_1+\beta h_2$ with $\alpha,\beta\in\R$,
and $\pi(v)=\beta\Delta$, so $v\in S_j$ if and only if $\beta\in[j,j+1)$; in that
case $v-\lfloor\alpha\rfloor h_1=\{\alpha\}h_1+\beta h_2$ is a lattice point of
$D_j$.  Hence $S_j=D_j+\Z h_1$, and distinct values of $\lfloor\alpha\rfloor$
give distinct points, so the union is disjoint.  Put $N:=P_\theta(D)<\infty$.

\emph{(a) A balanced set, and unambiguity.}  Take a $\sigma$-balanced
$B\subseteq S$ by Lemma~\ref{lem:balanced}.  We may assume $\sigma=+$: otherwise
replace $\su,h_1,\pi,h_2$ by $-\su,-h_1,-\pi,-h_2$, respectively.  This
keeps $h_1=c_1\su$ with $c_1\ge1$ and $\Delta>0$, and condition (iii) is
unchanged when the order of the integer heights is reversed.  Use these
oriented periods in the definitions of $D,D_j,S_j$ and $N$ above.
Since $\theta$ is not
periodic, the contrapositive of Lemma~\ref{lem:propagation} says: \emph{for all
$a\le b$ with $b-a\ge\operatorname{ht}(B)$ and every $x\in X_\theta$, if
$x=\theta$ on $\mathrm{St}[a,b-1]$ then $x=\theta$ on $\mathrm{St}[a,b]$.}
Iterating row by row,
\begin{equation}
x=\theta\ \text{on}\ \mathrm{St}[a,b],\ b-a\ge\operatorname{ht}(B)-1
\ \Longrightarrow\ x=\theta\ \text{on}\ \{\pi\ge a\}.
\label{eq:D-iterate}
\end{equation}

\emph{(b) A witness of non-periodicity.}  $N!\,h_2$ is not a period of $\theta$,
so there is a $v_0$ with $\theta(v_0+N!h_2)\ne\theta(v_0)$.  Put
$k:=\lfloor\pi(v_0)/\Delta\rfloor-N$, so that $\pi(v_0)\ge(k+N)\Delta$.

\emph{(c) Pigeonhole.}  Among the $N+1$ patterns $\theta|_{D_{k+j}}$,
$0\le j\le N$, at most $N$ are distinct, so there are $0\le j<j'\le N$ with
$\theta(v+(j'-j)h_2)=\theta(v)$ for all $v\in D_{k+j}$.  By \eqref{eq:D-diff}
this holds for all $v\in S_{k+j}$, that is
$x:=T^{(j'-j)h_2}\theta\in X_\theta$ satisfies $x=\theta$ on
$S_{k+j}=\mathrm{St}[(k+j)\Delta,(k+j+1)\Delta-1]$.  That strip has width
$\Delta-1\ge\operatorname{ht}(S)\ge\operatorname{ht}(B)$, so
\eqref{eq:D-iterate} gives $x=\theta$ on $\{\pi\ge(k+j)\Delta\}$, that is
\[\theta(v+(j'-j)h_2)=\theta(v)\qquad\text{whenever }\pi(v)\ge(k+j)\Delta .\]
Since $\pi(h_2)=\Delta>0$, that half-plane is invariant under $+(j'-j)h_2$;
applying the identity repeatedly, and using $(j'-j)\mid N!$, we get
$\theta(v+N!h_2)=\theta(v)$ for every $v$ with $\pi(v)\ge(k+j)\Delta$.  But
$\pi(v_0)\ge(k+N)\Delta\ge(k+j)\Delta$, so $\theta(v_0+N!h_2)=\theta(v_0)$,
contradicting (b).\hfill$\square$

\begin{rem}\label{rem:szab-general}
The proof of Theorem~\ref{thm:szabados} works for any finite-valued configuration
with a decomposition into two periodic configurations over an abelian group.
The required algebraic identity is \eqref{eq:D-diff}: the difference
$\theta(\cdot+sh_2)-\theta$ is $h_1$-periodic.  Over $\Z$ it is
the convex generalisation of \cite[Theorem~1]{Sz} recorded as
\cite[Theorem~1.7]{Colle23a}.
\end{rem}

\section{The Lean formalisation}\label{app:lean}

The formalisation is a Lean~4 project built on Mathlib; the toolchain is
\texttt{v4.33.1} and Mathlib is pinned at commit \texttt{0df444a}.  Its source
is the ancillary directory \texttt{anc/lean}.  The main theorem is
\texttt{Nivat.nivat\_conjecture}, in the file
\texttt{Nivat/Section8/Main.lean}, where it is derived from
\texttt{Nivat.convex\_nivat}.

\subsection*{The statement}
A reader who wants to trust the formal result needs to check only the statement
and the handful of definitions it mentions; everything else is checked by the
kernel.  They are, in the notation of this paper:
\begin{itemize}
\item \texttt{Config $\alpha$} is the type of maps $\Z\times\Z\to\alpha$, and
  \texttt{T u f} is the translate $z\mapsto f(z+u)$;
\item \texttt{Per f} is the subgroup $\{u: \texttt{T u f}=f\}$ and
  \texttt{IsPeriodic f} means that \texttt{Per f} contains a non-zero vector;
\item \texttt{pattern $\theta$ Q u} is the map $Q\to\alpha$, $q\mapsto\theta(u+q)$;
  \texttt{patterns $\theta$ Q} is the set of all of them, $u\in\Z^2$; and
  \texttt{P $\theta$ Q} is its cardinality (finite because $\alpha$ is);
\item \texttt{rectangle n k} is $\{1,\dots,n\}\times\{1,\dots,k\}$, of
  cardinality $nk$;
\item \texttt{LatticeConvex S} says that every $z\in\Z^2$ whose image in
  $\R^2$ lies in the convex hull of $S$ belongs to $S$.
\end{itemize}
With these, \texttt{nivat\_conjecture} assumes \texttt{[Finite $\alpha$]},
$0<n$, $0<k$ and $\texttt{P}\,\xi\,(\texttt{rectangle n k})\le nk$, and
concludes \texttt{IsPeriodic $\xi$}; \texttt{convex\_nivat} assumes
\texttt{[Finite $\alpha$]}, $S$ non-empty and lattice-convex and
$\texttt{P}\,\xi\,S\le|S|$, with the same conclusion.  These are
Corollary~\ref{cor:nivat} and Theorem~\ref{thm:convex-nivat}, without extra
hypotheses.

\subsection*{The architecture}
The top-level proof in \texttt{Nivat/Section8/} follows the proof of
Theorem~\ref{thm:convex-nivat}: an integer encoding of the alphabet
(\texttt{exists\_intEncoding}); a counterexample of minimal order
(\texttt{exists\_minimalCounterexample}, Remark~\ref{rem:minimal}); the
structure of \S\S\ref{ssec:firsthalf}--\ref{ssec:secondhalf}
(\texttt{structure\_input}); normalisation to a star configuration
(\texttt{exists\_starConfig\_of\_admissible}, Lemma~\ref{lem:normalise}); and
the contradiction with Theorem~A, \texttt{star\_complexity\_lower\_bound} in
\texttt{Nivat/MainTheorem.lean}.  Sections
\ref{sec:notation}--\ref{sec:quadratic} are formalised section by section in
the root of \texttt{Nivat/} (Laurent polynomials and Newton polygons in
\texttt{Laurent/}, zonotopes and unimodular triangulations in
\texttt{Lattice/}), and Appendix~\ref{app:szabados} in
\texttt{AppendixD.lean}.

The results quoted in \S\ref{ssec:external} are proved in
\texttt{Nivat/External/}: Kari--Szabados in \texttt{KSAnnihilator},
\texttt{KSDecomposition} and \texttt{KSCorollary}; Colle's theorem on
one-sided nonexpansive directions, with Boyle--Lind and the Cyr--Kra induction
behind it, in \texttt{Colle/Theorem114Final}, \texttt{Colle/BoyleLind} and
\texttt{Colle/Prop210}; and Colle's region, through his Lemma~3.5 and the
Claims of his \S4, in \texttt{Colle/ColleRegion} and about forty further
modules.  The formalisation follows the arXiv version of \cite{Colle23a}, in
which \cite[Theorem~1.9]{Colle23a} is numbered Theorem~1.14 and the region is
built in the proof of Lemma~4.5.

\subsection*{What formalising the source changed}
In the proof of Lemma~3.5(i) of the arXiv version of \cite{Colle23a}, sets
$\hat A_i^{(\varepsilon)}$ are obtained from the layers $\hat A_i$ by sweeping
along $-\vec v_{\ell_{J+1}}$ and intersecting with $\hat A_\infty^{(\varepsilon)}$,
and are asserted, without argument, to be enveloped and to lie between the
strips $B_i-k_i\vec v_\ell$ and $H_{B_i}(\ell)-k_i\vec v_\ell$.  In a numerical
model of $(\ell,\ell_J)$-regions this fails: over $2448$ instances
(ten families of generators and several choices of $\iota$, $J$, $\varepsilon$
and $i$), envelopedness fails in $1083$ and the strip containment in $768$.
The formalisation does not sweep.  It pushes only the $J$-face of each layer
outward by $\det(\text{prev},J)\cdot\det(J,\text{next})$ lattice layers and
verifies the seven properties a shell must have directly
(\texttt{Colle/Hole1PushShell}, theorem \texttt{exists\_chainData\_closed}).  The
statement of Lemma~3.5 is unchanged.  Similarly the monotonicity
$\hat A_i\subseteq\hat A_j$, which the source uses in a compactness argument but
records only in the caption of a figure, is proved in \texttt{Colle/AhatMono}.

\subsection*{Checking it}
With \texttt{elan} installed, from the root of the project:
\begin{center}
\texttt{lake exe cache get};\quad \texttt{lake build};\quad
\texttt{bash scripts/sorries.sh};\quad
\texttt{bash scripts/check\_axioms.sh Nivat.nivat\_conjecture}.
\end{center}
The \texttt{lakefile} builds every module under \texttt{Nivat/}, not only those
the main theorem imports.  On a fresh copy of the ancillary source the build
completes ($3368$ jobs, all $449$ modules), \texttt{sorries.sh} reports
\texttt{TOTAL: 0}, and \texttt{\#print axioms} gives
\texttt{[propext, Classical.choice, Quot.sound]} for both
\texttt{Nivat.nivat\_conjecture} and \texttt{Nivat.convex\_nivat}.  No constant
of the development is an \texttt{axiom} or \texttt{opaque} declaration.
Replaying the compiled development through the kernel once more with
\texttt{leanchecker}, both module by module and from an empty environment up to
\texttt{Nivat.Section8.Main}, reports no error.

\section{How this work was done}\label{app:ai}

This paper and its formalisation were produced by the author together with a
personal AI research system, an orchestration of language-model agents running
in Claude Code.  The agents recorded in the working sessions are
Claude Opus~5, with Claude Fable~5.1 on two formalisation lanes.  We record what the system did, how it was organised, and where it
was wrong, because a reader is entitled to know which parts of a proof were
searched for by machine and how they were checked.

\subsection*{Timeline}
The campaign on Nivat's conjecture began on 1 September 2026; the author chose
the target.  The proof of \S\S\ref{sec:notation}--\ref{sec:reduction} and the
paper were finished on 12 September.  The formalisation began on 13 September;
the three results quoted from the literature, at first stated as Lean axioms,
were all proved and their axioms deleted by 18 September; the development was
complete, with no \texttt{sorry} and no axiom beyond Lean's own, on
30 September.  It was rebuilt from scratch and re-audited independently on
1 October.

\subsection*{Organisation}
The formalisation ran as parallel lanes, each an agent owning a group of
lemmas of the blueprint (the chain construction, the placement of the region,
the tower of heights, the leaf lemmas of Colle's Lemma~3.5, and so on), up to
seven at a time.  A single integrating agent was the only one allowed to build
the main repository and to wire lanes into the main chain.  Separate agents
were tasked only with refutation: given a proposed lemma, find a
kernel-checked counterexample.  Auditors, given only the artefact and the
criteria and not the author's reasoning, checked statements against the source
papers.  The recorded formalisation sessions contain $277$ subagent runs.
The author set the target, ruled on every change to a theorem statement, and
directed the strategy: on 18 September the author cut the number of parallel
lanes from seven to two because open goals were not closing while the records
grew, and on 20 September required the formalisation to follow the source
papers as written.

\subsection*{Where the system was wrong, and how it was caught}
\begin{itemize}
\item \emph{A false clause in a draft of \S\ref{sec:reduction}.}  An
  adversarial audit of the draft on 12 September found a counterexample over
  $\mathbb F_2$ to a geometric clause of one proposition, and a gap in
  Appendix~\ref{app:szabados}.  Both were repaired the same day, the
  proposition by using a weaker conclusion that suffices; Theorems~A and~B
  were not affected.
\item \emph{False leaf lemmas.}  On 14 September three of the six leaf lemmas
  into which the region construction had been split were refuted by
  independent agents with kernel-checked counterexamples.  The cause was
  hypotheses dropped when the lemmas were extracted from the source.
\item \emph{A mis-stated external result.}  The first Lean statement of Colle's
  full-periodicity theorem was not derivable from the result of Kari and Moutot
  alone, as had been assumed; it needs a further step of Colle's argument.
\item \emph{Errors of the integrating agent.}  It drafted wrong signatures,
  cited a wrong line of the source, and twice chose data for a chain
  construction that a kernel-checked refutation later showed impossible.
\item \emph{Gaps in the source.}  The two points of Appendix~\ref{app:lean}
  — the shell sweep in Colle's Lemma~3.5(i) and the monotonicity recorded only
  in a figure caption — were found because the formalisation could not go
  through them.
\end{itemize}
None of these errors survives in the final development.  The kernel accepts
the main theorem, \texttt{Nivat.nivat\_conjecture}, with no assumption beyond
the axioms of Lean's logic.

% Inline bibliography; citation numbering is preserved.

\smallskip\noindent\emph{Theorem and lemma numbers of the cited works are those
of the originals.}

\clearpage
\begin{thebibliography}{99}

\bibitem{BoyleLind}
M.~Boyle and D.~Lind,
\emph{Expansive subdynamics},
Trans. Amer. Math. Soc. \textbf{349} (1997), no.~1, 55--102.

\bibitem{Colle23a}
C.~F. Colle,
\emph{On periodic decompositions, one-sided nonexpansive directions and Nivat's
conjecture},
Discrete Contin. Dyn. Syst. \textbf{43} (2023), no.~12, 4299--4327.

\bibitem{Colle23b}
C.~F. Colle,
\emph{On periodic decompositions and nonexpansive lines},
Math. Z. \textbf{303} (2023), art.~80.

\bibitem{CGalpha}
C.~F. Colle and E.~Garibaldi,
\emph{An alphabetical approach to Nivat's conjecture},
Nonlinearity \textbf{33} (2020), 3620--3652.

\bibitem{CG}
C.~F. Colle and E.~Garibaldi,
\emph{A modular structure theorem for minimal periodic decompositions and
periodicity of configurations with $P_\eta(4,n)\le 4n$},
preprint, \texttt{arXiv:2606.10193v1} [math.DS], 8 June 2026.

\bibitem{CyrKra15}
V.~Cyr and B.~Kra,
\emph{Nonexpansive $\Z^2$-subdynamics and Nivat's conjecture},
Trans. Amer. Math. Soc. \textbf{367} (2015), no.~9, 6487--6537.

\bibitem{CyrKra16}
V.~Cyr and B.~Kra,
\emph{Complexity of short rectangles and periodicity},
European J. Combin. \textbf{52} (2016), 146--173.

\bibitem{EKM}
C.~Epifanio, M.~Koskas and F.~Mignosi,
\emph{On a conjecture on bidimensional words},
Theoret. Comput. Sci. \textbf{299} (2003), 123--150.

\bibitem{KM20}
J.~Kari and E.~Moutot,
\emph{Decidability and periodicity of low complexity tilings},
in 37th International Symposium on Theoretical Aspects of Computer Science
(STACS 2020), LIPIcs vol.~154, Art.~No.~14, 2020.

\bibitem{KM23}
J.~Kari and E.~Moutot,
\emph{Decidability and periodicity of low complexity tilings},
Theory Comput. Syst. \textbf{67} (2023), 125--148.

\bibitem{KS}
J.~Kari and M.~Szabados,
\emph{An algebraic geometric approach to Nivat's conjecture},
Inform. and Comput. \textbf{271} (2020), 104481.

\bibitem{Khe}
A.~Khetan,
\emph{A counterexample to Nivat's conjecture for a non-convex window of full
affine span},
preprint, \texttt{arXiv:2607.09830v2}, July 2026.

\bibitem{Lean4}
L.~de~Moura and S.~Ullrich,
\emph{The Lean~4 theorem prover and programming language},
in: Automated Deduction -- CADE~28, Lecture Notes in Comput. Sci.
\textbf{12699}, Springer, 2021, 625--635.

\bibitem{Mathlib}
The mathlib Community,
\emph{The Lean mathematical library},
in: Proceedings of the 9th ACM SIGPLAN International Conference on Certified
Programs and Proofs (CPP 2020), 2020, 367--381.

\bibitem{MH38}
M.~Morse and G.~A. Hedlund,
\emph{Symbolic dynamics},
Amer. J. Math. \textbf{60} (1938), 815--866.

\bibitem{Nivat}
M.~Nivat,
\emph{Invited talk},
24th International Colloquium on Automata, Languages and Programming
(ICALP 1997), Bologna, 1997.

\bibitem{QZ}
A.~Quas and L.~Zamboni,
\emph{Periodicity and local complexity},
Theoret. Comput. Sci. \textbf{319} (2004), 229--240.

\bibitem{ST}
J.~W. Sander and R.~Tijdeman,
\emph{The rectangle complexity of functions on two-dimensional lattices},
Theoret. Comput. Sci. \textbf{270} (2002), 857--863.

\bibitem{Sz}
M.~Szabados,
\emph{Nivat's conjecture holds for sums of two periodic configurations},
in SOFSEM 2018: Theory and Practice of Computer Science,
Lecture Notes in Comput. Sci. vol.~10706, Springer, 2018, pp.~539--551.

\end{thebibliography}

\end{document}
